%% =====================================================================
%%  Molecular Simulation on Quantum Computing Systems
%%  CSE 402 -- Numerical Analysis, Simulation and Modeling
%%  Group 09, Section A -- ACM sigconf template
%%
%%  Build:   latexmk -pdf report.tex      (or pdflatex/bibtex/pdflatex x2)
%%  Tables:  python gen_tables.py         (regenerates tables/*.tex from the
%%           repository's saved scan files; nothing is typed in by hand)
%% =====================================================================
\documentclass[sigconf,screen,nonacm]{acmart}

\usepackage{mathtools}
\usepackage[version=4]{mhchem}
\usepackage{multirow}
\usepackage{array}
\usepackage{adjustbox}
\usepackage{listings}
\usepackage{algorithm}
\usepackage{algpseudocode}
\usepackage{placeins}
\usepackage{enumitem}
\usepackage{fontawesome5}

\graphicspath{{figures/}}

% chemical formulae inside headings should not upset the PDF bookmarks
\pdfstringdefDisableCommands{\def\ce#1{#1}}

% ---------------------------------------------------------------------
%  Float placement: let figures and tables sit close to the text that
%  discusses them (a placeins barrier before every section and molecule
%  subsection keeps them from drifting into the next topic).
% ---------------------------------------------------------------------
\renewcommand{\topfraction}{0.9}
\renewcommand{\bottomfraction}{0.7}
\renewcommand{\textfraction}{0.1}
\renewcommand{\floatpagefraction}{0.7}
\renewcommand{\dbltopfraction}{0.9}
\renewcommand{\dblfloatpagefraction}{0.7}
\setcounter{topnumber}{3}
\setcounter{bottomnumber}{2}
\setcounter{totalnumber}{6}
\setcounter{dbltopnumber}{2}
\setlength{\textfloatsep}{10pt plus 2pt minus 2pt}
\setlength{\floatsep}{8pt plus 2pt minus 2pt}
\setlength{\intextsep}{8pt plus 2pt minus 2pt}
\setlength{\dbltextfloatsep}{10pt plus 2pt minus 2pt}
\setlength{\emergencystretch}{2em}

% ---------------------------------------------------------------------
%  Spacing fixes.
%  (1) \raggedbottom: never stretch the space between paragraphs and
%      headings to fill a column (acmart defaults to \flushbottom, which
%      produced the large blank gaps inside columns).
%  (2) After \FloatBarrier flushes pending floats onto float-only pages,
%      LaTeX can start the next column with a stale, too-small column
%      height (\@colroom), so the column ended after a few lines.  When a
%      column starts empty and holds no top/bottom floats, reset its
%      height to the full column height.
% ---------------------------------------------------------------------
\raggedbottom
\makeatletter
\newcommand{\@resetcolroom}{%
  \ifdim\pagetotal=\z@
    \ifx\@toplist\@empty\ifx\@botlist\@empty\ifx\@midlist\@empty
      \global\@colroom\@colht
      \global\vsize\@colht
    \fi\fi\fi
  \fi}
\let\@origFloatBarrier\FloatBarrier
\renewcommand{\FloatBarrier}{\@origFloatBarrier\@resetcolroom}
%  (3) \MoleculeBarrier, used between the molecule subsections: like
%      \FloatBarrier, but it only waits for column-width floats.  Page-wide
%      floats (figure*, table*) can only go to the top of a later page
%      anyway, so waiting for them used to leave the rest of the page
%      blank.  Sections and subsections keep the strict \FloatBarrier.
\newif\if@narrowpending
\newcommand{\@checknarrow}[1]{\ifdim\wd#1>\columnwidth\else\global\@narrowpendingtrue\fi}
\newcommand{\MoleculeBarrier}{\par\begingroup
  \global\@narrowpendingfalse
  \let\@elt\@checknarrow \@botlist\@deferlist
  \if@narrowpending
    \ifx\@fltovf\relax
      \if@firstcolumn\clearpage\else\null\newpage\MoleculeBarrier\fi
    \else
      \newpage\let\@fltovf\relax\MoleculeBarrier
    \fi
  \fi\endgroup\@resetcolroom}
\makeatother

% ---------------------------------------------------------------------
%  Project banner (course / group / section / supervisor + GitHub link).
%  Used under the title, at the end of the abstract and in the
%  "Code and Installation" section.
% ---------------------------------------------------------------------
\newcommand{\BannerCourse}{CSE 402: Numerical Analysis, Simulation and Modeling \textbullet{} Group 09 \textbullet{} Section A \textbullet{} Supervisor: Ishrat Jahan}
\newcommand{\BannerGitHub}{\href{https://github.com/QuitttCat/CSE-402-Numerical-Analysis-Simulation-and-Modeling-Project}{\faGithub\hspace{0.3em}https://github.com/QuitttCat/CSE-402-Numerical-Analysis-Simulation-and-Modeling-Project}}
\newcommand{\ProjectBanner}{{\centering\small\BannerCourse\par\smallskip\BannerGitHub\par}}

% ---------------------------------------------------------------------
%  Figure helpers.  Plots are one column wide; the two-panel surface
%  renders are far too detailed for a column, so they span the page.
%  \figpair places two column-width figures one after the other.
% ---------------------------------------------------------------------
\newcommand{\figone}[3]{%
  \begin{figure}[!htbp]\centering
  \includegraphics[width=\columnwidth]{#1}
  \Description{Plot or render described in the caption.}
  \caption{#2}\label{fig:#3}
  \end{figure}}

\newcommand{\figwide}[3]{%
  \begin{figure*}[!t]\centering
  \includegraphics[width=0.94\textwidth]{#1}
  \Description{Plot or render described in the caption.}
  \caption{#2}\label{fig:#3}
  \end{figure*}}

\newcommand{\figpair}[6]{\figone{#1}{#2}{#3}\figone{#4}{#5}{#6}}

% ---------------------------------------------------------------------
%  Code listings (floating, one column wide)
% ---------------------------------------------------------------------
\definecolor{lstbg}{gray}{0.96}
\definecolor{lstkw}{rgb}{0.10,0.25,0.60}
\definecolor{lstcm}{rgb}{0.30,0.45,0.30}
\definecolor{lststr}{rgb}{0.60,0.20,0.15}
\lstset{
  language=Python,
  basicstyle=\ttfamily\fontsize{6.6}{7.9}\selectfont,
  keywordstyle=\color{lstkw}\bfseries,
  commentstyle=\color{lstcm}\itshape,
  stringstyle=\color{lststr},
  backgroundcolor=\color{lstbg},
  numbers=left, numberstyle=\tiny\color{gray}, numbersep=4pt,
  frame=leftline, framerule=0.5pt, rulecolor=\color{gray},
  breaklines=true, breakatwhitespace=false,
  postbreak=\mbox{\textcolor{gray}{$\hookrightarrow$}\space},
  showstringspaces=false, columns=fullflexible,
  captionpos=b, xleftmargin=9pt, aboveskip=2pt, belowskip=2pt,
  float=tbp,
  literate={--}{{-{}-}}2
}

\newcommand{\mHa}{\,\text{mHa}}
\newcommand{\Ha}{\,\text{Ha}}
\newcommand{\Ang}{\text{\AA}}

% ---------------------------------------------------------------------
%  Title block
% ---------------------------------------------------------------------
\acmConference[CSE 402]{CSE 402 Project Report}{Group 09, Section A}{BUET, Dhaka, Bangladesh}
\settopmatter{authorsperrow=3}

\begin{document}

\title{Molecular Simulation on Quantum Computing Systems}
\subtitle{Variational Ground-State Energy Surfaces and Their Reconstruction by Polynomial and Spline Interpolation\\[6pt]{\normalfont\normalsize\BannerCourse}\\[3pt]{\normalfont\small\BannerGitHub}}

\author{Shahraz Hussain Siyam}
\affiliation{\position{Student ID: 2105023}\department{Department of CSE}\institution{Bangladesh University of Engineering and Technology}\city{Dhaka}\country{Bangladesh}}
\author{Nabiha Tahseen}
\affiliation{\position{Student ID: 2105031}\department{Department of CSE}\institution{Bangladesh University of Engineering and Technology}\city{Dhaka}\country{Bangladesh}}
\author{Debashri Roy}
\affiliation{\position{Student ID: 2105035}\department{Department of CSE}\institution{Bangladesh University of Engineering and Technology}\city{Dhaka}\country{Bangladesh}}
\author{Niloy Kumar Mondal}
\affiliation{\position{Student ID: 2105044}\department{Department of CSE}\institution{Bangladesh University of Engineering and Technology}\city{Dhaka}\country{Bangladesh}}
\author{Amit Saha}
\affiliation{\position{Student ID: 2105049}\department{Department of CSE}\institution{Bangladesh University of Engineering and Technology}\city{Dhaka}\country{Bangladesh}}

% ---------------------------------------------------------------------
%  Abstract (body 150--250 words) + project banner at the end
% ---------------------------------------------------------------------
\begin{abstract}
This report simulates the ground-state potential energy of eleven molecules with the Variational Quantum Eigensolver (VQE) and then asks a numerical-analysis question: given a few expensive samples of a smooth function, how well can it be filled in between them? We ran 23 scans, 692 independent variational optimisations over active spaces of six to ten qubits, and checked every energy against exact diagonalisation of the same active-space Hamiltonian. Fourteen of the 23 scans fall inside the 1.6~mHa chemical-accuracy threshold; the rest miss it because the optimiser ran out of budget on the deepest circuits or the ansatz could not represent the state, not because the variational principle failed. Each scan was reconstructed by piecewise-linear interpolation, a single global polynomial, and cubic or bicubic splines, all scored out of sample: leave-one-out cross-validation in one dimension, a half-grid holdout in two. Splines reduce the error of the linear baseline by a factor of about 2 to 15 on most scans and place equilibrium bond lengths within 0.001--0.09~\AA{} of the reference from grids spaced 0.075--0.25~\AA{} apart, roughly an order of magnitude of resolution for no extra quantum cost. One result ran against expectation: on \ce{LiH} the global polynomial, which textbook theory says should collapse into Runge's phenomenon, beats the cubic spline by a factor of 38. We show why, find the node spacing at which the failure does appear, and use the full campaign to decide which reconstruction method to trust and when.
\par\medskip
\ProjectBanner
\end{abstract}

\ccsdesc[500]{Computing methodologies~Quantum mechanic simulation}
\ccsdesc[300]{Mathematics of computing~Interpolation}
\ccsdesc[300]{Applied computing~Chemistry}
\ccsdesc[100]{Hardware~Quantum computation}

\keywords{Variational quantum eigensolver, potential energy surface, polynomial interpolation, cubic spline, Runge's phenomenon, cross-validation, numerical simulation, quantum chemistry}

\maketitle

% =====================================================================
\FloatBarrier
\section{Introduction and Base Paper Review}
\label{sec:intro}

\subsection{Background and motivation}
The central quantity in computational chemistry is the ground-state
electronic energy of a molecule as a function of where its nuclei sit. That
function, the \emph{potential energy surface} (PES), decides equilibrium
geometries, reaction barriers and vibrational frequencies. Almost everything
a chemist wants to predict comes from the shape of this surface, not from any
single point on it.

Two facts make the PES a clean subject for a numerical-analysis project, and
this report leans on both.

The first is that computing one point on the surface is a well-posed
eigenvalue problem. Under the Born--Oppenheimer approximation the nuclei are
held fixed, and the electronic Hamiltonian becomes a finite Hermitian matrix
whose smallest eigenvalue is the energy we want. For the small active spaces
used here, that matrix can be diagonalised exactly, so every approximate
result in this report can be checked against ground truth, a comparison
simulation work rarely gets to make.

The second is that each point is expensive, and gets more expensive as the
system grows. The matrix has dimension $2^{n_q}\times2^{n_q}$ in the qubit
representation, so exact diagonalisation scales exponentially with the number
of spin-orbitals kept. The Variational Quantum Eigensolver
(VQE)~\cite{peruzzo2014,mcclean2016} sidesteps that exponential storage cost
by replacing diagonalisation with minimisation of a Rayleigh quotient over a
family of trial states. It is approximate, it is why ``quantum computing for
chemistry'' is an active research area, and, seen from the outside, it is
simply an expensive black-box function $\theta \mapsto E(\theta)$ handed to a
classical optimiser.

That expense is what turns building a PES into an interpolation problem. A
one-dimensional scan of a bond length at $0.1\Ang$ resolution over a
$1.5\Ang$ window costs sixteen variational optimisations. A two-dimensional
surface at the same resolution on both axes costs the square of that. The
most expensive single scan in this project, the \ce{BeH2} surface, needed
$110$ optimisations; halving the grid spacing on both axes would have needed
$440$. Meanwhile, what a chemist actually wants from the surface, the
location of the minimum, the curvature at the bottom of the well, the height
of a barrier, is a property of the continuous function, not of the sampled
points. Interpolation bridges that gap, and how well it bridges it is a
numerical-analysis question with a measurable answer.

\subsection{Base paper review}
\label{sec:basepaper}

\paragraph{Citation.}
The base paper for this project is A.~Peruzzo, J.~McClean, P.~Shadbolt,
M.-H.~Yung, X.-Q.~Zhou, P.~J.~Love, A.~Aspuru-Guzik and J.~L.~O'Brien,
``A variational eigenvalue solver on a photonic quantum processor,''
\emph{Nature Communications} 5, 4213 (2014)~\cite{peruzzo2014}, the paper
that introduced the Variational Quantum Eigensolver.

\paragraph{Problem addressed.}
The standard quantum route to molecular ground-state energies, quantum phase
estimation, needs long, deep, fully coherent circuits, far beyond what
near-term devices can run. The base paper asks whether the same quantity can
be obtained with only short quantum circuits, by moving most of the work to a
classical computer.

\paragraph{Method.}
The paper splits the calculation into a hybrid loop. The quantum processor
prepares a parametrised trial state and estimates the expectation value of
each Pauli term of the qubit Hamiltonian; the classical computer sums the
terms into an energy and hands it to a classical optimiser (Nelder--Mead in the
paper), which proposes new parameters. The variational principle guarantees
that every energy the loop reports is an upper bound on the true ground-state
energy, so minimising it drives the trial state towards the ground state. Each
energy evaluation needs only a short state preparation, at the price of many
repeated measurements.

\paragraph{Demonstration and findings.}
The algorithm was run on a two-qubit photonic chip to compute the
ground-state energy of the helium hydride ion \ce{HeH+} in a minimal basis
over a range of bond lengths. The paper's central claim is that variational
hybrid algorithms trade circuit depth for repetitions, which makes molecular
simulation a realistic target for small, noisy quantum hardware. McClean
\emph{et~al.}~\cite{mcclean2016} later formalised the theory, and Kandala
\emph{et~al.}~\cite{kandala2017} introduced the hardware-efficient ansatz
used in this report.

\paragraph{How this work builds on it.}
We implement the same loop (Algorithm~\ref{alg:vqe}) and rely on the same
variational bound as a correctness check. We extend the base paper in four
directions: from one two-qubit molecule to eleven molecules on six to ten
qubits; from a single bond coordinate to bends, a dihedral twist and five
two-dimensional surfaces; from a device-specific parametrisation to UCCSD and
hardware-efficient ansatze driven by classical optimisers, including an exact
parameter-shift gradient; and, most importantly for this course, from
sampling a curve point by point to asking how accurately the curve can be
reconstructed \emph{between} the sampled points. The base paper samples the
dissociation curve but never asks how to fill it in; that interpolation
question is the numerical-analysis core of this report.

\paragraph{Related work.}
Our project plan also surveyed the broader VQE literature. The review of
Tilly \emph{et~al.}~\cite{tilly2022} covers methods and best practices for the
algorithm. Two open difficulties it highlights, the choice of a good ansatz
and optimisation landscapes that flatten into barren
plateaus~\cite{larocca2024bp} where gradient variance vanishes, motivated
generative approaches such as the Generative Quantum
Eigensolver~\cite{nakaji2024gqe}, which trains a GPT-2 style model to propose
circuits and was demonstrated on \ce{H2}, \ce{LiH}, \ce{BeH2} and \ce{N2}.
This report does not pursue those directions. It takes the standard VQE as
given and studies a question all of them leave untouched: what to do with the
expensive samples once they exist.

\subsection{Core objective}

The core objective of this work is twofold. First, to implement the
Variational Quantum Eigensolver as a numerical algorithm and verify it, point
by point, against exact diagonalisation of the same active-space Hamiltonian
across a broad campaign of molecules and geometric coordinates. Second, to
treat the resulting potential energy surfaces as a reconstruction problem, a
smooth function known only at a few expensive samples, and to measure
strictly out of sample how well piecewise-linear, global polynomial and spline
interpolation recover it, so as to decide which reconstruction method to
trust, and when.

\subsection{What this report covers}
\label{sec:contributions}

This report documents a full simulation campaign and the numerical analysis
built on top of it.

\textbf{Simulation.} We implement VQE in two forms. A core library
(\texttt{src/}) is built from scratch on CUDA-Q~\cite{cudaq} and
PySCF~\cite{pyscf}: a hardware-efficient ansatz, a UCCSD ansatz, seven SciPy
optimisers and a hand-written gradient-descent routine that computes its
gradients exactly through the parameter-shift rule~\cite{schuld2019} rather
than through finite differences. The $23$-scan campaign itself is produced by
per-molecule drivers that use Qiskit Nature's UCCSD
ansatz~\cite{qiskit}, an exact statevector estimator and SciPy's SLSQP or
COBYLA. Section~\ref{sec:theory} develops the theory and
Section~\ref{sec:impl} the implementation.

\textbf{Campaign.} We run $23$ scans over $11$ molecules, \ce{H2}, \ce{LiH},
\ce{BeH2}, \ce{H2O}, \ce{NH3}, \ce{N2}, \ce{O2}, \ce{F2}, \ce{H2O2}, \ce{NO2}
and \ce{N2O}, covering bond stretches, bend angles, a dihedral twist, and
five two-dimensional surfaces. Table~\ref{tab:campaign} lists all of them.
Every energy is checked against exact diagonalisation of the same
active-space Hamiltonian (Table~\ref{tab:vqeacc}).

\textbf{Reconstruction.} We reconstruct every scan with piecewise-linear
interpolation, a global polynomial, and a cubic or bicubic spline, and score
each one out of sample, leave-one-out in 1-D, half-grid holdout in 2-D.
Results are collected in Table~\ref{tab:interp} and discussed per molecule in
Section~\ref{sec:results}.

\textbf{Geometry recovery.} We use the reconstructions to find equilibrium
geometries at sub-grid resolution and compare them against experimental and
model reference values (Table~\ref{tab:geom}).

\textbf{An honest negative result.} On \ce{LiH} the global polynomial does
not show Runge's phenomenon over the scanned interval, and beats the cubic
spline by a factor of $38$. We report this, explain why, and find the node
density at which the expected failure does show up
(Section~\ref{sec:lih}).

\subsection{What this project is not}

Two things are worth clearing up before the rest of the report. First, none
of the main results need physical quantum hardware: the quantum circuits are simulated exactly on a classical CPU by a statevector simulator. This is
a numerical simulation project, and VQE appears here as a numerical
algorithm, an ansatz, a cost function and an optimiser, not as a piece of
hardware engineering. A supplementary demonstration that runs the same
calculation on real IBM Quantum hardware is discussed in
Section~\ref{sec:ibm}, mainly to show the gap between noiseless simulation
and present-day devices.

Second, the absolute energies reported here are not meant to be chemically
predictive. Every calculation uses a minimal STO-3G basis and a small
complete active space. The energies are internally consistent and exactly
checkable, but they are not converged to the basis-set limit, and the
equilibrium geometries we recover deviate from experiment by amounts that
belong to the \emph{model}, not to the solver. We say explicitly, for every
number, which reference it is being measured against.

\subsection{How the rest of the report is organised}

Section~\ref{sec:method} is the methodology: the quantum-chemical and
variational theory behind VQE (Section~\ref{sec:theory}), then the
interpolation methods and how we validate them, with pseudocode for the main
algorithms (Section~\ref{sec:numerics}). Section~\ref{sec:impl} covers
implementation and experiments: the software, with code excerpts
(Section~\ref{sec:software}), the simulation setup and computational protocol
(Section~\ref{sec:protocol}), and the parameter choices, datasets and
benchmarks (Section~\ref{sec:params}). Section~\ref{sec:resdisc} presents the
results molecule by molecule, with every figure placed next to the discussion
it belongs to (Section~\ref{sec:results}), then the campaign-wide accuracy and
reconstruction results (Sections~\ref{sec:accuracy} and~\ref{sec:recon}) and
the discussion that only makes sense across molecules, including a comparison
with the base paper (Section~\ref{sec:discussion}).
Section~\ref{sec:conclusion} gives the key takeaways, the limitations and the
future work, and Section~\ref{sec:code} explains how to obtain and run the
code.

% =====================================================================
\FloatBarrier
\section{Methodology}
\label{sec:method}

\subsection{Background: quantum chemistry and the variational method}
\label{sec:theory}
\subsubsection{The electronic structure problem}

For a molecule with $M$ nuclei of charge $Z_A$ at positions $\mathbf{R}_A$ and
$N$ electrons at positions $\mathbf{r}_i$, the non-relativistic Hamiltonian in
atomic units~\cite{szabo,helgaker} is
\begin{equation}
\hat{H} = -\sum_{i=1}^{N}\frac{\nabla_i^2}{2}
          -\sum_{i=1}^{N}\sum_{A=1}^{M}\frac{Z_A}{|\mathbf{r}_i-\mathbf{R}_A|}
          +\sum_{i<j}\frac{1}{|\mathbf{r}_i-\mathbf{r}_j|},
\label{eq:elec-ham}
\end{equation}
where nuclear kinetic energy has been dropped and nuclear repulsion enters as
an additive constant. This is the Born--Oppenheimer approximation: a proton
is roughly $1836$ times heavier than an electron, so the electrons settle
into equilibrium essentially instantly at any nuclear configuration, and the
electronic energy can be computed at a fixed $\{\mathbf{R}_A\}$. The map
\begin{equation}
\{\mathbf{R}_A\} \;\longmapsto\; E_0(\{\mathbf{R}_A\})
\end{equation}
is the potential energy surface this project samples and reconstructs.

The difficulty sits entirely in the last term of~\eqref{eq:elec-ham}. Drop
electron--electron repulsion and the problem separates into $N$ independent
one-electron problems. Keep it, and the exact wavefunction lives in a space
whose dimension grows combinatorially with the number of electrons and basis
functions.

\subsubsection{Second quantisation and the active space}

Expanding the wavefunction in a finite one-particle basis
$\{\phi_p\}_{p=1}^{K}$ (here the minimal STO-3G set) and moving to second
quantisation turns~\eqref{eq:elec-ham} into
\begin{equation}
\hat{H} = \sum_{pq} h_{pq}\, a_p^\dagger a_q
        + \tfrac{1}{2}\sum_{pqrs} h_{pqrs}\, a_p^\dagger a_q^\dagger a_r a_s
        + E_{\text{nuc}},
\label{eq:secondquant}
\end{equation}
where $a_p^\dagger,a_p$ create and annihilate an electron in spin-orbital
$p$, and the one- and two-body integrals $h_{pq}$, $h_{pqrs}$ are computed
classically by PySCF. The Hilbert space now has dimension $2^{2K}$ before any
symmetry restriction, still exponential.

The standard remedy, used throughout this work, is the \emph{complete active
space} (CAS) approximation. Orbitals split into a set held doubly occupied, a
set held empty, and an active set of $n_{\text{act}}$ orbitals holding
$n_{\text{el}}$ electrons, where full configuration interaction is carried
out. We write this as $(n_{\text{el}}\text{e}, n_{\text{act}}\text{o})$. The
choice trades accuracy for tractability; Table~\ref{tab:campaign} records the
active space used for every scan, ranging from $(2\text{e},5\text{o})$ for
\ce{LiH} to $(10\text{e},6\text{o})$ for \ce{F2}.

This is the single most consequential approximation in the report. It is why
the computed equilibrium geometries differ from experiment by tenths of an
\AA ngstr\"om in some cases, and it is also why exact diagonalisation stays
feasible, which is why every variational result here can be checked.

\subsubsection{Fermion-to-qubit mapping}

To write~\eqref{eq:secondquant} in terms of qubits, the fermionic operators
must be mapped onto Pauli operators in a way that preserves the
anticommutation relations $\{a_p, a_q^\dagger\} = \delta_{pq}$. The
Jordan--Wigner transformation~\cite{jordanwigner} does this with
\begin{equation}
a_p^\dagger = \Big(\bigotimes_{k<p} Z_k\Big)\otimes \frac{X_p - iY_p}{2},
\end{equation}
where the leading string of $Z$ operators enforces antisymmetry under
exchange. The Bravyi--Kitaev encoding~\cite{bravyikitaev} does the same job
with $\mathcal{O}(\log K)$ rather than $\mathcal{O}(K)$ operator weight.
Either way, the result is a qubit Hamiltonian
\begin{equation}
\hat{H} = \sum_{j} c_j\, \hat{P}_j,
\qquad
\hat{P}_j \in \{I,X,Y,Z\}^{\otimes n_q},
\label{eq:pauli}
\end{equation}
a weighted sum of Pauli strings with real coefficients $c_j$. Every energy
evaluation in this project is an expectation value of the
form~\eqref{eq:pauli}, and the qubit counts $n_q$ in Table~\ref{tab:campaign}
(six, eight or ten) are twice the number of active spatial orbitals minus
two: the parity mapping used by the experiment drivers removes two qubits by
exploiting the conserved particle-number parities.

\subsubsection{The variational principle}

The whole VQE algorithm rests on one fact from linear algebra, which a
numerical-analysis reader already knows as the Rayleigh quotient bound. For
any normalised state $|\psi\rangle$,
\begin{equation}
\boxed{\;\langle\psi|\hat{H}|\psi\rangle \;\geq\; E_0
= \min\operatorname{spec}(\hat{H}),\;}
\label{eq:variational}
\end{equation}
with equality only if $|\psi\rangle$ is a ground eigenvector. The proof is
short: expand $|\psi\rangle=\sum_k c_k|\phi_k\rangle$ in the eigenbasis, so
$\langle\psi|\hat{H}|\psi\rangle=\sum_k |c_k|^2 E_k \geq E_0\sum_k|c_k|^2 =
E_0$.

Equation~\eqref{eq:variational} turns an eigenvalue problem into a
minimisation problem, and minimisation is something numerical analysis knows
how to handle. Given a family of trial states
$|\psi(\boldsymbol\theta)\rangle$ that we can cheaply prepare and measure,
\begin{equation}
E(\boldsymbol\theta)
= \langle\psi(\boldsymbol\theta)|\hat{H}|\psi(\boldsymbol\theta)\rangle,
\end{equation}
minimising $E$ over $\boldsymbol\theta$ gives the best estimate of $E_0$
available inside that family. Equation~\eqref{eq:variational} also
guarantees the estimate is an \emph{upper bound}: the optimiser can never
report an energy below the true ground state, so any energy that comes out
lower than the exact reference signals a bug, not a discovery. We use this
one-sided error as a correctness check throughout the report.

\subsubsection{Ansatz design}

The trial-state family is the \emph{ansatz}. Choosing one is the same
trade-off as choosing a polynomial degree in Section~\ref{sec:numerics}:
expressive power against optimisation difficulty. We implement two ansatz
families.

\paragraph{Hardware-efficient ansatz.}
The hardware-efficient ansatz~\cite{kandala2017} prepares the Hartree--Fock
reference, then applies $d$ layers of an independent $R_y$ rotation on every
qubit followed by a CNOT ladder:
\begin{equation}
|\psi(\boldsymbol\theta)\rangle = \prod_{\ell=1}^{d}
\Big[\, U_{\text{CNOT}} \prod_{q=1}^{n_q} R_y(\theta_{\ell q}) \Big]
\,|\psi_{\text{HF}}\rangle .
\end{equation}
The parameter count is exactly $p = n_q d$. The rotations give it flexibility;
the entangling ladder gives it the correlation a Hartree--Fock product state
cannot represent on its own. This family knows nothing about the chemistry
it is fitting: it is cheap and shallow, and it comes with no guarantee of
containing a good approximation to the true ground state.

\paragraph{UCCSD ansatz.}
The unitary coupled-cluster singles-and-doubles ansatz is chemically
motivated:
\begin{equation}
|\psi(\boldsymbol\theta)\rangle
= e^{\hat{T}(\boldsymbol\theta)-\hat{T}^\dagger(\boldsymbol\theta)}
  |\psi_{\text{HF}}\rangle,
\qquad
\hat{T} = \hat{T}_1 + \hat{T}_2,
\end{equation}
where $\hat{T}_1$ and $\hat{T}_2$ generate single and double excitations out
of the Hartree--Fock determinant. Its parameter count is not a free choice:
it is fixed by the number of symmetry-allowed excitations, which is why the
$p$ column of Table~\ref{tab:campaign} ranges so widely, from $20$ parameters
for \ce{NO2} in a $(5\text{e},4\text{o})$ space to $117$ for \ce{N2} in
$(6\text{e},6\text{o})$. UCCSD is far more expressive than the
hardware-efficient ansatz, but its circuits are deeper and its optimisation
landscape harder, which shows up directly in the accuracy results of
Section~\ref{sec:protocol}.

\subsubsection{Classical optimisation and the parameter-shift rule}

From the optimiser's point of view, $E(\boldsymbol\theta)$ is just an
expensive black-box scalar function. The codebase wires up eight optimisers
in total. Three are gradient-free methods from \texttt{scipy.optimize}~\cite{scipy},
COBYLA~\cite{powell1994}, Nelder--Mead and Powell; four more are gradient-based or
quasi-Newton methods from the same library, BFGS, L-BFGS-B, CG and SLSQP,
which by default estimate their gradients by finite differences. Which one
was used for each scan is recorded in Table~\ref{tab:campaign}.

The eighth optimiser is a gradient-descent routine we wrote from scratch, for
a reason specific to this course: the gradient of $E$ can be computed
\emph{exactly}, with no finite-difference truncation error. Because each
$\theta_i$ enters the circuit only through a rotation
$\exp(-i\theta_i \hat{P}/2)$ with $\hat{P}^2 = I$, we have the \emph{parameter-shift rule}~\cite{schuld2019}
\begin{equation}
\frac{\partial E}{\partial \theta_i}
= \frac{1}{2}\Big[
  E\big(\boldsymbol\theta + \tfrac{\pi}{2}\mathbf{e}_i\big)
- E\big(\boldsymbol\theta - \tfrac{\pi}{2}\mathbf{e}_i\big)\Big].
\label{eq:paramshift}
\end{equation}
This looks like a central difference, but the step is a fixed, macroscopic
$\pi/2$, and the result is \emph{exact}, not $\mathcal{O}(h^2)$. There is no
truncation error to trade off against floating-point cancellation, so the
usual optimal-step-size analysis simply does not apply. It is a rare case in
which a numerical derivative carries no discretisation error at all. Its
cost is $2p$ energy evaluations per gradient, which is why COBYLA remains the
practical default for production scans.

Putting the last three subsections together gives the full VQE loop that
every scan in this report runs at every geometry, summarised in
Algorithm~\ref{alg:vqe}. Steps 3 and 4 are the only quantum part; everything
else is classical bookkeeping around a scalar minimisation.

\begin{algorithm}[!ht]
\caption{Variational Quantum Eigensolver at one fixed geometry}
\label{alg:vqe}
\begin{algorithmic}[1]\small
\Require qubit Hamiltonian $\hat H=\sum_j c_j\hat P_j$, ansatz
  $|\psi(\boldsymbol\theta)\rangle$ with $p$ parameters, optimiser, tolerance
  $\varepsilon$
\Ensure estimate $E^\star \geq E_0$ of the ground-state energy
\State $\boldsymbol\theta \gets$ initial guess (Hartree--Fock reference, all
  rotation angles zero)
\Repeat
  \State prepare $|\psi(\boldsymbol\theta)\rangle$ on the (simulated) circuit
  \State measure $E(\boldsymbol\theta) = \sum_j c_j
    \langle\psi(\boldsymbol\theta)|\hat P_j|\psi(\boldsymbol\theta)\rangle$
  \If{optimiser needs a gradient}
    \For{$i = 1, \dots, p$}
      \State $g_i \gets \tfrac12\big[E(\boldsymbol\theta+\tfrac{\pi}{2}
        \mathbf e_i) - E(\boldsymbol\theta-\tfrac{\pi}{2}\mathbf e_i)\big]$
        \Comment{parameter-shift rule, exact}
    \EndFor
  \EndIf
  \State $\boldsymbol\theta \gets$ optimiser step using $E(\boldsymbol\theta)$
    (and $\mathbf g$, if used)
\Until{$|E(\boldsymbol\theta) - E_{\text{prev}}| < \varepsilon$ or the
  evaluation budget is spent}
\State \Return $E^\star \gets E(\boldsymbol\theta)$
\end{algorithmic}
\end{algorithm}

A full scan repeats Algorithm~\ref{alg:vqe} once per sampled geometry,
building up exactly the small, expensive, unevenly reliable data set that
Section~\ref{sec:numerics} then has to reconstruct.

Optimiser choice interacts with circuit depth in a way that shows up
throughout the campaign. \texttt{NEW\_MOLECULE\_PROMPT.md}, the project's own
notes for extending the pipeline to a new molecule, records a concrete
comparison on \ce{O2}'s hardest geometry: tripling COBYLA's evaluation budget
improved the energy by $0.6\mHa$ ($1.7\%$), while warm-starting the optimiser
from the previous geometry's converged parameters improved it by $4.0\mHa$,
roughly \emph{seven times} more, for no extra cost at all. More budget helps
a stuck optimiser a little; a better starting point helps it a lot. Both
tools appear in the codebase, and Section~\ref{sec:discussion} returns to
this trade-off when it explains which scans do and do not reach chemical
accuracy.

\subsubsection{The reference calculations}

Three energies are recorded at every geometry, and keeping them apart matters
for reading every table in this report.

\begin{itemize}[leftmargin=1.5em]
\item $E_{\text{HF}}$, the Hartree--Fock mean-field energy. Each electron
      moves in the averaged field of the others, and correlation is absent
      entirely. This is the starting point the variational calculation is
      meant to improve on.
\item $E_{\text{VQE}}$, the variational result: the quantity under test.
\item $E_{\text{exact}}$, exact diagonalisation of the identical active-space
      Hamiltonian, computed with \texttt{numpy.linalg.eigvalsh}. This is
      ground truth for the model, not for the molecule.
\end{itemize}

Two error measures follow:
\begin{align}
\delta &= E_{\text{VQE}} - E_{\text{exact}}
   && \text{(variational error, $\geq 0$)},\\
E_{\text{corr}} &= E_{\text{HF}} - E_{\text{exact}}
   && \text{(correlation energy)} .
\end{align}
$E_{\text{corr}}$ is what the whole exercise is trying to recover; $\delta$
is how much of it VQE missed. A run is only doing useful work when $\delta
\ll E_{\text{corr}}$, and Table~\ref{tab:vqeacc} reports both side by side so
that ratio can be read off directly. The conventional target is
\emph{chemical accuracy}, $\delta < 1.6\mHa$
($1\,\text{kcal}\,\text{mol}^{-1}$), the threshold below which errors in
predicted reaction rates become tolerable.

\subsection{Numerical methods for surface reconstruction}
\label{sec:numerics}
Once we have $n$ expensive samples $\{(x_i, y_i)\}_{i=1}^{n}$ of a smooth
function, $y_i$ an energy and $x_i$ a geometric coordinate, the rest of the
project is a classical reconstruction problem: fill in the function between
the samples. This section covers the three methods we use and, just as
important, how we grade them.

\subsubsection{The piecewise-linear baseline}

Every raw scan plot in this project joins neighbouring points with straight
lines. That default is itself an interpolation scheme, and the crudest one
available:
\begin{equation}
L(x) = y_i + \frac{y_{i+1}-y_i}{x_{i+1}-x_i}(x - x_i),
\quad x \in [x_i, x_{i+1}].
\end{equation}
Its weaknesses matter because they motivate everything that follows. $L$ is
continuous but not differentiable at the nodes; its second derivative is
zero on every interval and undefined at every node, so it carries no
curvature information; and the minimum it implies is always exactly one of
the sampled points. A piecewise-linear reconstruction can never locate an
equilibrium geometry more finely than the grid spacing, which for these scans
is $0.1$--$0.25\Ang$, one to two orders of magnitude coarser than the
precision chemists usually quote. For equally spaced nodes with spacing $h$,
the classical bound is
\begin{equation}
\|f - L\|_\infty \leq \frac{h^2}{8}\,\|f''\|_\infty ,
\label{eq:linbound}
\end{equation}
so halving the grid quarters the error, but doubles a cost that is already
quadratic on a surface. That is exactly the wrong direction, and it is why a
better interpolant is worth having. The \emph{Linear} columns in
Table~\ref{tab:interp} report this baseline throughout.

\subsubsection{Global polynomial interpolation}

Given $n$ distinct nodes, exactly one polynomial of degree at most $n-1$
passes through all of them~\cite{burden,trefethen}. In Lagrange form,
\begin{equation}
P_{n-1}(x) = \sum_{i=1}^{n} y_i \ell_i(x),
\qquad
\ell_i(x) = \prod_{j\neq i}\frac{x-x_j}{x_i-x_j}.
\label{eq:lagrange}
\end{equation}
On paper this looks like the ideal answer: one smooth closed-form curve,
infinitely differentiable, exact at every node, with no piecewise
bookkeeping. The classical error theorem gives, for $f \in C^n$,
\begin{equation}
f(x) - P_{n-1}(x) = \frac{f^{(n)}(\xi)}{n!}\prod_{i=1}^{n}(x - x_i)
\label{eq:polyerror}
\end{equation}
for some $\xi$ in the interval. The $1/n!$ factor is reassuring; the factors
$f^{(n)}(\xi)$ and the node polynomial are not. For equally spaced nodes the
node polynomial grows sharply near the ends of the interval, and $f^{(n)}$
can grow faster than $n!$ shrinks.

\paragraph{Runge's phenomenon.}
The resulting failure is Runge's phenomenon~\cite{runge1901}: on equally
spaced nodes, a high-degree interpolant can oscillate violently between
nodes, and the oscillation gets worse as more nodes are added. Runge's own
example, $f(x)=1/(1+25x^2)$ on $[-1,1]$, diverges in the supremum norm as
$n\to\infty$.

Two points matter for reading our results correctly, and Section~\ref{sec:lih}
comes back to both of them.

First, a zero fit residual is not evidence of accuracy. The interpolant
passes through every node by construction, so the in-sample residual is zero
to machine precision, our runs record maxima as low as $1.4\times
10^{-11}\mHa$, while the interpolant may be doing anything at all between
nodes. Any honest check has to be out of sample.

Second, Runge's phenomenon is a property of the function being interpolated,
not of the degree by itself. It shows up when the function has a
singularity in the complex plane close to the real interval; Runge's example
has poles at $\pm i/5$. A function that stays analytic over a wide
neighbourhood of the interval is approximated very well by high-degree
polynomials. Whether a given PES falls into the good case or the bad case is
an empirical question, and we answer it separately for each scan.

\paragraph{Degree sweeps.}
Because the best degree is not known in advance, each 1-D scan is fitted at
every degree from $1$ up to $n-2$ and scored out of sample. The resulting
curves, the \texttt{*\_degree\_sweep.png} figures throughout
Section~\ref{sec:results}, tabulated in the \texttt{tab:deg:*} tables, show
error falling as the model gains the flexibility to follow the function,
then, in the Runge regime, rising again as it gains the flexibility to
oscillate. The minimum of that curve is the degree we actually use.

\subsubsection{Cubic spline interpolation}

The spline resolves the same tension differently: instead of one
high-degree polynomial over the whole interval, fit a separate cubic on each
subinterval and force the pieces to join smoothly. On $[x_i,x_{i+1}]$,
\begin{equation}
S_i(x) = a_i + b_i(x-x_i) + c_i(x-x_i)^2 + d_i(x-x_i)^3,
\end{equation}
subject to the interpolation conditions $S_i(x_i)=y_i$,
$S_i(x_{i+1})=y_{i+1}$ and, at every interior node, the continuity conditions
\begin{equation}
S_{i-1}'(x_i) = S_i'(x_i),
\qquad
S_{i-1}''(x_i) = S_i''(x_i),
\end{equation}
\cite{deboor}. These constraints fix all coefficients up to two boundary
conditions (here the \texttt{not-a-knot} condition, SciPy's default). The
resulting linear system is tridiagonal and diagonally dominant, so it solves
in $\mathcal{O}(n)$ operations and is unconditionally stable.

The structural difference from~\eqref{eq:lagrange} that matters is
\emph{locality}: each cubic piece is determined mostly by nearby data, so a
perturbation at one node decays geometrically along the spline rather than
propagating globally. There is no mechanism for a Runge-style blow-up, and
the classical bound
\begin{equation}
\|f - S\|_\infty \leq \frac{5}{384}\,h^4\,\|f^{(4)}\|_\infty
\label{eq:splinebound}
\end{equation}
is $\mathcal{O}(h^4)$ rather than the $\mathcal{O}(h^2)$
of~\eqref{eq:linbound}: halving the grid cuts the error by sixteen. Because
the spline is $C^2$, it also has a meaningful second derivative, which is
what makes estimating vibrational frequencies and barrier curvature possible
at all.

\subsubsection{Two-dimensional reconstruction}

The five surface scans need 2-D versions of all three methods.

\paragraph{Bilinear} The straight-line baseline generalises to bilinear
interpolation on a rectangular grid: interpolate along one axis, then along
the other. It inherits every weakness of its 1-D counterpart.

\paragraph{Bicubic spline} The tensor-product cubic spline
\begin{equation}
S(x,y) = \sum_{i=0}^{3}\sum_{j=0}^{3} c_{ij}\,(x-x_k)^i (y-y_l)^j
\end{equation}
is fitted on the rectangular grid, cubic in each direction, with SciPy's
\texttt{RectBivariateSpline}. The spline order on each axis drops
automatically when that axis has too few points to support a cubic, a real
constraint on the $10\times5$ \ce{H2O2} grid.

\paragraph{Total-degree polynomial} The 2-D analogue of the polynomial fits
\begin{equation}
P(x,y) = \sum_{i+j\leq d} c_{ij}\, x^i y^j,
\end{equation}
which has $\binom{d+2}{2}$ coefficients. Unlike the 1-D case this is
generally a \emph{least-squares} fit rather than an exact interpolant,
because a rectangular grid of $N$ points does not in general admit an exact
total-degree interpolant. The in-sample residual is therefore nonzero, and we
report it alongside the out-of-sample error, a distinction that turns out to
matter for \ce{N2O} and \ce{NO2}.

\subsubsection{Validation methodology}
\label{sec:holdout}

In-sample residuals tell us nothing about an interpolant, so every
reconstruction error in this report is measured out of sample.

\paragraph{Leave-one-out cross-validation (1-D).}
For each interior node $i$, we rebuild the interpolant on the remaining
$n-1$ points and evaluate it at the withheld $x_i$:
\begin{equation}
e_i = \mathcal{I}_{\setminus i}(x_i) - y_i .
\end{equation}
Algorithm~\ref{alg:loo} states this precisely. Endpoints are excluded:
removing an endpoint turns interpolation into extrapolation, a different and
harder question. Nodes where the remaining data cannot support the requested
degree are excluded too and marked \texttt{n/a}. This is why the fully
interpolating degree $n-1$ often cannot be scored at all, and why the tables
report a \emph{best} degree from the sweep instead.

\begin{algorithm}[!ht]
\caption{Leave-one-out cross-validation at the interior nodes}
\label{alg:loo}
\begin{algorithmic}[1]\small
\Require nodes $x_1,\dots,x_n$, values $y_1,\dots,y_n$, polynomial degree $d$
\Ensure error $e_i$ at every interior node $i=2,\dots,n-1$
\For{$i = 2, \dots, n-1$}
  \State $X \gets \{x_1,\dots,x_n\} \setminus \{x_i\}$,\quad
         $Y \gets \{y_1,\dots,y_n\} \setminus \{y_i\}$
  \If{$|X| \geq d+1$}
    \State fit degree-$d$ polynomial $P$ to $(X,Y)$
    \State $e_i^{\text{poly}} \gets P(x_i) - y_i$
  \Else
    \State $e_i^{\text{poly}} \gets$ \texttt{n/a}
  \EndIf
  \State $e_i^{\text{linear}} \gets \mathrm{linear\_interp}(X,Y,x_i) - y_i$
\EndFor
\end{algorithmic}
\end{algorithm}

\paragraph{Half-grid holdout (2-D).}
Leave-one-out does not work on a grid interpolant: remove one cell and the
grid stops being rectangular, so \texttt{RectBivariateSpline} cannot be built
on what remains. Instead, we keep every other value on each axis, which keeps
the subgrid rectangular while withholding about three quarters of the
points. The interpolant is built on the coarse subgrid and scored on
everything held out, as in Algorithm~\ref{alg:holdout}.

\begin{algorithm}[!ht]
\caption{Half-grid holdout for a 2-D surface}
\label{alg:holdout}
\begin{algorithmic}[1]\small
\Require full grid $Z[i,j]$ over axes $x_1,\dots,x_m$ and $y_1,\dots,y_k$
\Ensure error at every held-out grid point
\State $B \gets \{1,3,5,\dots\}$, \quad $A \gets \{1,3,5,\dots\}$
  \Comment{every other index on each axis}
\State build coarse grid from $x[B]$, $y[A]$, $Z[B,A]$
\State fit bicubic spline $S$ and bilinear interpolant $\ell$ on the coarse
  grid
\For{each $(i,j)$ \textbf{not} in $B\times A$}
  \State $e^{\text{spline}}_{ij} \gets S(x_i,y_j) - Z[i,j]$
  \State $e^{\text{linear}}_{ij} \gets \ell(x_i,y_j) - Z[i,j]$
\EndFor
\end{algorithmic}
\end{algorithm}

This holdout answers the question that actually matters for an $N^2$-cost
campaign: \emph{if we had run only half the grid on each axis, at a quarter
of the quantum cost, how much would we have lost?} A good answer is what
licenses coarse grids on molecules too expensive to sample densely, and it
is the single most practically useful measurement in this report.

\paragraph{Error metrics.}
All errors are reported in milli-Hartree over the withheld set
$\mathcal{S}$:
\begin{align}
\text{max} &= \max_{i\in\mathcal{S}} |e_i| \times 10^3, \\
\text{RMS} &= \Big(\tfrac{1}{|\mathcal{S}|}\textstyle\sum_{i\in\mathcal{S}} e_i^2\Big)^{1/2}\times 10^3, \\
\text{MAE} &= \tfrac{1}{|\mathcal{S}|}\textstyle\sum_{i\in\mathcal{S}} |e_i| \times 10^3 .
\end{align}
We rank by RMS because it penalises the isolated large excursions typical of
polynomial failure, which MAE would water down.

\subsubsection{Locating the minimum}

We extract the equilibrium geometry from a reconstruction by dense sampling
to bracket the minimum, then a bounded scalar minimiser. The search is kept
inside the data range on purpose: outside it, a high-degree polynomial
diverges quickly, and any minimum found there would just be an artefact of
extrapolation. We also count the number of interior minima of the fitted
polynomial, since a physically sensible curve over these intervals should
have exactly one; a count above one is direct evidence of spurious
oscillation. Results are in Table~\ref{tab:geom}.

% =====================================================================
\FloatBarrier
\section{Implementation and Experiments}
\label{sec:impl}

\subsection{Software architecture}
\label{sec:software}

The codebase is roughly $24{,}400$ lines of Python in three layers: a
molecular-modelling and VQE core, a per-molecule experiment layer, and a
numerical post-processing layer. This section covers the design decisions
that affect the results we report.

\subsubsection{Layer 1: core (\texttt{src/})}

\texttt{molecule.py} (1887 lines) wraps NVIDIA CUDA-Q Solvers'
\texttt{create\_molecule}~\cite{cudaq}. It runs the classical
pre-computation, Hartree--Fock, MP2, CASCI, CASSCF, CCSD and FCI through
PySCF~\cite{pyscf}, and exposes the qubit Hamiltonian together with
energies, one- and two-body integrals, qubit counts, the Pauli decomposition,
an exact-diagonalisation reference, the Hartree--Fock reference state and
UCCSD operator pools. It also renders molecules three different ways: through
py3Dmol, through Plotly, and through a dependency-free 3Dmol.js HTML page
that needs nothing but a browser. A \texttt{visualize()} helper picks a
backend automatically, py3Dmol inside a notebook, Plotly otherwise, or a
standalone HTML file if neither is available, and these are the renderings
shown alongside the energy curves in the live scan windows. The four
generated files in \texttt{out/} (\texttt{caffeine\_ball\_and\_stick.html},
\texttt{benzene\_spacefill.html}, \texttt{h2o\_ball\_and\_stick.html} and
\texttt{nh3\_plotly.html}) are checked-in examples of exactly this: two
standalone 3Dmol.js pages, one HTML export from py3Dmol, and one Plotly
export, produced by \texttt{examples/molecule\_demo.py}.

\texttt{vqe.py} (740 lines) implements the algorithm of
Section~\ref{sec:theory}, including both ansatze, all eight optimisers, and
the parameter-shift gradient. Listing~\ref{lst:hea} shows the
hardware-efficient ansatz: the Hartree--Fock reference is prepared by
flipping the first $n_{\text{el}}$ qubits, after which each layer applies one
independent $R_y$ rotation per qubit and a nearest-neighbour CNOT ladder.

\begin{lstlisting}[caption={Hardware-efficient ansatz. The parameter count is
exactly \texttt{n\_qubits * depth}; increasing \texttt{depth} enlarges the
trial-state family at the cost of a harder optimisation, the same
flexibility-versus-conditioning trade-off as raising a polynomial degree.},
label={lst:hea}]
def build_hardware_efficient_ansatz(
        n_qubits, n_electrons, depth=2):
    n_params = n_qubits * depth

    @cudaq.kernel
    def kernel(thetas: list[float]):
        qubits = cudaq.qvector(n_qubits)
        # (a) Hartree-Fock reference
        for i in range(n_electrons):
            x(qubits[i])
        # (b) `depth` layers of rotations + entanglers
        idx = 0
        for _layer in range(depth):
            for q in range(n_qubits):
                ry(thetas[idx], qubits[q])
                idx += 1
            for q in range(n_qubits - 1):
                x.ctrl(qubits[q], qubits[q + 1])

    return kernel, n_params
\end{lstlisting}

Listing~\ref{lst:psr} shows the hand-written optimiser implementing the
parameter-shift rule of~\eqref{eq:paramshift} and Algorithm~\ref{alg:vqe}. We
kept it in the codebase on purpose, so the gradient computation is visible
rather than hidden inside a library call.

\begin{lstlisting}[caption={Gradient descent with parameter-shift gradients.
The inner loop costs $2p$ energy evaluations per gradient, which is why
COBYLA is the practical default for production scans.},
label={lst:psr}]
theta = np.array(theta0, dtype=float)
p = len(theta)
prev_energy = cost(theta)
for k in range(max_iterations):
    grad = np.empty(p)
    for i in range(p):
        shift_vec = np.zeros(p)
        shift_vec[i] = shift          # pi/2, exact
        grad[i] = 0.5 * (cost(theta + shift_vec)
                         - cost(theta - shift_vec))
    theta = theta - learning_rate * grad
    energy = cost(theta)
    if abs(prev_energy - energy) < tol:
        prev_energy = energy
        break
    prev_energy = energy
\end{lstlisting}

Two short scripts under \texttt{examples/} exercise this layer end to end
without touching any single molecule's experiment folder:
\texttt{molecule\_demo.py} builds caffeine, benzene, water and ammonia and
renders each one through all three visualisation backends, and
\texttt{vqe\_demo.py} runs \ce{H2} three ways back to back, UCCSD with
COBYLA, UCCSD with the hand-written gradient descent, and the
hardware-efficient ansatz with COBYLA, plotting the live loss curves and
printing a side-by-side comparison. They are the fastest way to see every
ansatz, optimiser and visualisation backend work without running a full
molecular scan.

One platform note affects reproducibility: \texttt{cudaq-solvers} is
declared in \texttt{requirements.txt} as a Linux/WSL-only dependency
(\texttt{sys\_platform != "win32"}), so the \texttt{src/} core and the demo
scripts need WSL or a Linux machine. The campaign drivers use Qiskit and PySCF
instead (Section~\ref{sec:params}) and do not depend on it.

\subsubsection{Layer 2: experiments (\texttt{experiment/})}

Each molecule has a directory with a \texttt{*\_ground\_state\_%
estimation.py} driver. A run sweeps a geometric coordinate, solves the
variational problem at every point, and plots the curve live next to a
rendering of the deforming molecule. When the sweep finishes, the window
switches to \emph{review mode}: arrow keys or a slider walk back and forth
along the completed curve while the molecular rendering redraws to match the
selected geometry.

\paragraph{The chemistry engine of the drivers.}
Every driver follows the same recipe. PySCF builds the Hamiltonian for the
geometry, restricted to a complete active space by Qiskit Nature's
\texttt{ActiveSpaceTransformer} where the full problem is too large. The parity
mapper, with two-qubit reduction, maps it onto qubits. A UCCSD circuit built on
the Hartree--Fock reference state is transpiled once, outside the optimiser
loop, and its energy is evaluated by Qiskit's exact
\texttt{StatevectorEstimator}, with the nuclear-repulsion and inactive-orbital
constants added back. SciPy's \texttt{minimize} then drives the parameters,
warm-started from the previous geometry. The \texttt{src/} core of
Layer~1 implements the same algorithm on CUDA-Q instead and is exercised by the
demo scripts.

\paragraph{Persistence and the \texttt{scan\_io} contract.}
VQE points are expensive: a fraction of a second each for \ce{H2}, but two to
four minutes each for the larger active spaces, so a single \ce{H2O} stretch
scan runs for roughly forty minutes and the \ce{BeH2} surface considerably
longer. No result should ever be paid for twice. Every completed scan
therefore saves itself automatically to two files:

\begin{itemize}[leftmargin=1.5em]
\item \texttt{*\_scan.json}, energies plus full metadata (coordinate, fixed
      values, active space, optimiser, range). This is what \texttt{--load}
      reads, and what every table in this report was generated from.
\item \texttt{*\_scan.csv}, the same energies as a flat table, for
      spreadsheets and external plotting.
\end{itemize}

The format is shared across all molecules through
\texttt{experiment/}\allowbreak\texttt{scan\_io.py} (140 lines), so the drivers cannot drift into
writing files each other cannot read. Each file records a
\texttt{format\_version} and is refused outright on a mismatch instead of
half-loading and failing confusingly later, a small decision that saves real
time on a project with several people touching the same code. Reloading a
saved scan re-enters the interactive review window in seconds, with no
chemistry re-executed.

\paragraph{Animating a surface.}
Three scripts, \texttt{beh2\_surface\_gif.py}, \texttt{h2o2\_surface\_gif.py}
and \texttt{h2o2\_twist\_gif.py}, turn a saved scan into an animated GIF
showing the surface rising out of the grid point by point (the animation
lives in \texttt{assets/}). The three \texttt{*\_anim\_still.png} figures in
this report, marked as animation frames wherever they appear, are single
frames pulled from those GIFs, since a printed page cannot show motion.

\paragraph{Extending the pipeline.}
\texttt{NEW\_MOLECULE\_PROMPT.md} is the project's internal playbook for
adding a new molecule: how to decide whether a coordinate needs a curve or a
surface, which code patterns to reuse (the \texttt{ParityMapper}, transpiling
once instead of per point, warm-starting from the previous geometry, a
serpentine traversal order for surface grids), and a list of concrete traps
found the hard way, among them an open-shell molecule silently defaulting to
a restricted reference, an active space too small to keep the energy curve
monotonic, and the cost of SLSQP's finite-difference gradient scaling with
the number of parameters. It functioned as the project's shared debugging
notebook as much as a template for new code.

\subsubsection{Layer 3: numerical post-processing}

Each molecule additionally has \texttt{*\_polynomial/} and \texttt{*\_spline/}
directories. These run no chemistry at all: they load a saved scan and do
pure numerical analysis on it, so they finish in seconds and can be re-run
freely. Listing~\ref{lst:loo} shows the leave-one-out driver behind
Algorithm~\ref{alg:loo}, the single most important measurement in the
post-processing layer.

\begin{lstlisting}[caption={Leave-one-out cross-validation at the interior
nodes. Endpoints return \texttt{NaN} because withholding one turns the
problem into extrapolation; nodes where the remaining data cannot support the
requested degree are excluded the same way.},
label={lst:loo}]
def leave_one_out(x, y, degree):
    n = len(x)
    poly_err   = np.full(n, np.nan)
    linear_err = np.full(n, np.nan)

    for i in range(1, n - 1):
        keep = np.ones(n, dtype=bool)
        keep[i] = False
        xk, yk = x[keep], y[keep]
        # fitting degree d needs at least d+1 points
        if len(xk) >= degree + 1:
            poly_err[i] = float(
                Polynomial.fit(xk, yk, degree)(x[i])) - y[i]
        linear_err[i] = float(np.interp(x[i], xk, yk)) - y[i]

    return {"polynomial": poly_err, "linear": linear_err}
\end{lstlisting}

Listing~\ref{lst:holdout} shows the 2-D half-grid holdout of
Algorithm~\ref{alg:holdout}, including the automatic spline-order reduction
needed when a coarse axis has fewer than four points.

\begin{lstlisting}[caption={Half-grid holdout for surfaces. Taking every
other value on each axis keeps the subgrid rectangular, which leave-one-out
cannot do, while withholding roughly three quarters of the points.},
label={lst:holdout}]
bi = np.arange(0, len(bond), 2)
aj = np.arange(0, len(angle), 2)
coarse_bond, coarse_angle = bond[bi], angle[aj]
coarse_Z = Z[np.ix_(bi, aj)]

# spline order must stay below the number of points per axis
kx = min(3, len(coarse_bond)  - 1)
ky = min(3, len(coarse_angle) - 1)
spline = RectBivariateSpline(coarse_bond, coarse_angle,
                             coarse_Z, kx=kx, ky=ky)

held_out = np.ones_like(Z, dtype=bool)
held_out[np.ix_(bi, aj)] = False
for i in range(len(bond)):
    for j in range(len(angle)):
        if not held_out[i, j] or not np.isfinite(Z[i, j]):
            continue
        spline_err[i, j] = float(
            spline(bond[i], angle[j])[0, 0]) - Z[i, j]
        linear_err[i, j] = _bilinear(
            coarse_bond, coarse_angle, coarse_Z,
            bond[i], angle[j]) - Z[i, j]
\end{lstlisting}

\subsubsection{Reproducibility}

Because every scan is saved with its metadata, every number and every figure
in this report can be regenerated without re-running any quantum simulation.
The tables in this document were generated programmatically from the stored
JSON files rather than typed in by hand, which rules out an entire class of
transcription error. Section~\ref{sec:repro} lists the commands.

\subsection{Simulation setup and computational protocol}
\label{sec:protocol}
All calculations use the minimal STO-3G basis with an active space chosen
per molecule. Geometries are built analytically from the scanned internal
coordinates, with every other coordinate held at the fixed value recorded in
the scan metadata. For each geometry the driver builds the active-space
Hamiltonian with PySCF, maps it to qubits, starts the ansatz parameters at
the Hartree--Fock reference, minimises the energy with the recorded
optimiser, and stores $E_{\text{HF}}$, $E_{\text{VQE}}$, $E_{\text{exact}}$,
the qubit count, the parameter count, the converged parameter vector, and,
where the optimiser reports it, the number of energy evaluations used.

Table~\ref{tab:campaign} is the full inventory: $23$ scans over $11$
molecules, $692$ variational optimisations in total. Qubit counts range from
six to ten, and ansatz parameter counts from $20$ to $117$.
Table~\ref{tab:surfaces} gives the detail on the five two-dimensional grids.

% GENERATED by gen_tables.py -- do not edit by hand (tab:campaign)
\begin{table*}[!t]
\caption{Complete computational campaign. Every VQE scan stored in the repository, with the active space, qubit count and ansatz parameter count actually used. $N$ is the number of geometries at which the variational problem was solved. Bond lengths are in \AA{} and angles in degrees.}
\label{tab:campaign}
\centering\small\setlength{\tabcolsep}{3.5pt}

\adjustbox{max width=\textwidth}{%
\begin{tabular}{llllrrrlll}
\toprule
Molecule & Coordinate & Type & Range / step & $N$ & $n_q$ & $p$ & Active space & Optimizer & Basis \\
\midrule
\ce{H2} & bond length & curve & 0.3--3 / 0.05 & 55 & -- & -- & full & default & sto3g \\
\addlinespace[1pt]
\ce{LiH} & bond length & curve & 0.9--3.9 / 0.25 & 13 & 8 & 24 & (2e, 5o) & slsqp & sto3g \\
\addlinespace[1pt]
\ce{BeH2} & bend angle & curve & 90--180 / 5 & 19 & 6 & 26 & (4e, 4o) & default & sto3g \\
\ce{BeH2} & symmetric stretch & curve & 0.9--2.4 / 0.1 & 16 & 6 & 26 & (4e, 4o) & default & sto3g \\
\ce{BeH2} & 2-D surface & surface & 11$\times$10 grid & 110 & 6 & 26 & (4e, 4o) & default & sto3g \\
\addlinespace[1pt]
\ce{H2O} & bend angle & curve & 70--160 / 5 & 19 & 6 & 26 & (4e, 4o) & default & sto3g \\
\addlinespace[1pt]
\ce{NH3} & bend angle & curve & 96--120 / 3 & 9 & 8 & 54 & (6e, 5o) & slsqp & sto3g \\
\ce{NH3} & symmetric stretch & curve & 0.85--1.6 / 0.075 & 11 & 8 & 54 & (6e, 5o) & slsqp & sto3g \\
\ce{NH3} & 2-D surface & surface & 7$\times$7 grid & 49 & 8 & 54 & (6e, 5o) & slsqp & sto3g \\
\addlinespace[1pt]
\ce{N2} & bond length & curve & 0.8--2.4 / 0.2 & 9 & 10 & 117 & (6e, 6o) & cobyla & sto3g \\
\addlinespace[1pt]
\ce{O2} & bond length & curve & 0.9--2.5 / 0.2 & 9 & 10 & 68 & (8e, 6o) & cobyla & sto3g \\
\addlinespace[1pt]
\ce{F2} & bond length & curve & 0.9--2.5 / 0.1 & 17 & 10 & 35 & (10e, 6o) & default & sto3g \\
\addlinespace[1pt]
\ce{H2O2} & bend angle & curve & 80--140 / 5 & 13 & 10 & 92 & (8e, 6o) & cobyla & sto3g \\
\ce{H2O2} & symmetric stretch & curve & 1--3 / 0.1 & 21 & 10 & 92 & (8e, 6o) & cobyla & sto3g \\
\ce{H2O2} & 2-D surface & surface & 10$\times$5 grid & 50 & 10 & 92 & (8e, 6o) & cobyla & sto3g \\
\ce{H2O2} & dihedral twist & curve & 0--180 / 10 & 19 & 10 & 92 & (8e, 6o) & cobyla & sto3g \\
\addlinespace[1pt]
\ce{NO2} & bend angle & curve & 90--180 / 5 & 19 & 6 & 20 & (5e, 4o) & cobyla & sto3g \\
\ce{NO2} & symmetric stretch & curve & 0.9--2.4 / 0.1 & 16 & 8 & 37 & (7e, 5o) & cobyla & sto3g \\
\ce{NO2} & 2-D surface & surface & 11$\times$10 grid & 110 & 6 & 20 & (5e, 4o) & cobyla & sto3g \\
\addlinespace[1pt]
\ce{N2O} & bend angle & curve & 90--180 / 5 & 19 & 6 & 26 & (4e, 4o) & default & sto3g \\
\ce{N2O} & N--N stretch & curve & 0.9--2.2 / 0.1 & 14 & 8 & 54 & (6e, 5o) & default & sto3g \\
\ce{N2O} & N--O stretch & curve & 0.9--2.3 / 0.1 & 15 & 8 & 54 & (6e, 5o) & default & sto3g \\
\ce{N2O} & 2-D surface & surface & 10$\times$6 grid & 60 & 6 & 26 & (4e, 4o) & default & sto3g \\
\addlinespace[1pt]
\bottomrule
\end{tabular}}
\end{table*}

% GENERATED by gen_tables.py -- do not edit by hand (tab:surfaces)
\begin{table}[!htbp]
\caption{The five two-dimensional scans. $N$ is the total number of variational solves, which grows as the product of the two axis lengths; this quadratic cost is the reason the holdout experiment of Section~\ref{sec:holdout} matters.}
\label{tab:surfaces}
\centering\footnotesize\setlength{\tabcolsep}{3pt}
\adjustbox{max width=\columnwidth}{%
\begin{tabular}{llrr}
\toprule
Molecule & Grid (axis $\times$ axis) & $N$ & $n_q$ \\
\midrule
\ce{BeH2} & bond: 11 pts 0.9--2.4, angle: 10 pts 90--180 & 110 & 6 \\
\ce{NH3} & bond: 7 pts 0.85--1.6, angle: 7 pts 96--120 & 49 & 8 \\
\ce{H2O2} & dihedral: 10 pts 0--180, angle: 5 pts 80--140 & 50 & 10 \\
\ce{NO2} & bond: 11 pts 0.9--2.4, angle: 10 pts 90--180 & 110 & 6 \\
\ce{N2O} & rnn: 10 pts 0.9--2.25, rno: 6 pts 0.95--1.7 & 60 & 6 \\
\bottomrule
\end{tabular}}
\end{table}


\subsection{Parameter choices, datasets and benchmarks}
\label{sec:params}

\paragraph{Chemistry model.}
Every calculation uses the minimal STO-3G basis. PySCF supplies the orbitals,
from restricted Hartree--Fock for the closed-shell molecules and restricted
open-shell Hartree--Fock for the open-shell \ce{O2} and \ce{NO2}. An active
space $(n_{\text{el}}\text{e},n_{\text{act}}\text{o})$ chosen per molecule,
from $(2\text{e},5\text{o})$ for \ce{LiH} to $(10\text{e},6\text{o})$ for
\ce{F2}, keeps the problem at six to ten qubits (Table~\ref{tab:campaign}).
Where a bond breaks, the larger active space is used: our own tests showed
that $(4\text{e},4\text{o})$ is too small to describe a breaking bond, and
yields a curve that dips back down past its minimum instead of rising towards
dissociation. This is why the \ce{N2O} stretches use $(6\text{e},5\text{o})$
while its bend and surface use $(4\text{e},4\text{o})$.

\paragraph{Ansatz, optimiser and starting point.}
Every campaign scan uses a UCCSD ansatz. Its parameter count is fixed by the
active space, from $20$ for \ce{NO2} to $117$ for \ce{N2} (the $p$ column of
Table~\ref{tab:campaign}). The first geometry of a scan starts from the
Hartree--Fock reference with all parameters zero; every later geometry
warm-starts from the previous geometry's converged parameters, falling back to
zeros if the parameter count changes. SciPy's SLSQP is used for the smaller
ans\"atze (\ce{H2}, \ce{LiH}, \ce{BeH2}, \ce{H2O}, \ce{NH3}, \ce{F2},
\ce{N2O}) and COBYLA for the deepest ones (\ce{N2}, \ce{O2}, \ce{H2O2},
\ce{NO2}), each with a ceiling of $1000$ iterations. The \emph{Optimizer}
column of Table~\ref{tab:campaign} records the choice.

\paragraph{Scan grids.}
The grids are a deliberate budget decision, because coarse sampling is
exactly what the reconstruction experiment is about. Bond lengths are sampled
from $0.05\Ang$ steps for \ce{H2} up to $0.25\Ang$ for \ce{LiH}; bend angles
at $3^\circ$ or $5^\circ$; the \ce{H2O2} dihedral at $10^\circ$; and the five
surfaces on grids of between $7\times7$ and $11\times10$ points
(Table~\ref{tab:surfaces}).

\paragraph{Datasets.}
No external dataset is used. The data are the $23$ scans produced by this
project, stored in the repository as \texttt{*\_scan.json} (energies plus
full metadata) and \texttt{*\_scan.csv}. Together they hold the $692$ VQE
energies, each with its Hartree--Fock and exact reference energy and the
converged parameter vector. Every table and figure in this report is
generated from those files.

\paragraph{Benchmarks and metrics.}
Three energies are recorded at every geometry, $E_{\text{HF}}$,
$E_{\text{VQE}}$ and $E_{\text{exact}}$ (exact diagonalisation of the
identical active-space Hamiltonian), and the variational error is judged
against the chemical-accuracy threshold of $1.6\mHa$. Reconstructions are
scored strictly out of sample, by leave-one-out cross-validation in one
dimension and the half-grid holdout in two (Section~\ref{sec:holdout}), using
max, RMS and MAE errors, and are always compared with the piecewise-linear
baseline. Recovered equilibrium geometries are compared with the experimental
value where the scan metadata records one, and otherwise with the model's own
optimised geometry (Table~\ref{tab:geom}).

% =====================================================================
\FloatBarrier
\section{Results and Discussion}
\label{sec:resdisc}

\subsection{Results by molecule}
\label{sec:results}
Molecules appear in order of increasing electron count, from \ce{H2} to
\ce{N2O}. Each subsection states the model setup, gives the variational
results, and then discusses the reconstruction, with every figure placed
immediately next to the paragraph that discusses it. All figures are
reproduced unmodified from the repository; the three marked \emph{animation
frame} are single frames pulled from animated GIFs, which a printed page
cannot show moving.

\subsubsection{\ce{H2}: hydrogen}
\label{sec:h2}


% GENERATED by gen_tables.py -- do not edit by hand (tab:data:h2:bond)
\begin{table*}[!t]
\caption{\ce{H2} bond length scan: Hartree--Fock, VQE and exact energies (Ha) at every geometry; $\delta = (E_{\mathrm{VQE}}-E_{\mathrm{exact}})\times 10^3$ (mHa).}
\label{tab:data:h2:bond}
\centering\small\setlength{\tabcolsep}{3.5pt}
\adjustbox{max width=\textwidth}{%
\begin{tabular}{rrrrr@{\hskip 1.2em}rrrrr}
\toprule
$r$ (\AA) & $E_{\mathrm{HF}}$ & $E_{\mathrm{VQE}}$ & $E_{\mathrm{exact}}$ & $\delta$ & $r$ (\AA) & $E_{\mathrm{HF}}$ & $E_{\mathrm{VQE}}$ & $E_{\mathrm{exact}}$ & $\delta$ \\
\midrule
0.3 & -0.593828 & -0.601804 & -0.601804 & 0.000 & 1.7 & -0.854338 & -0.971427 & -0.971427 & 0.000 \\
0.35 & -0.780455 & -0.789269 & -0.789269 & 0.000 & 1.75 & -0.841349 & -0.966335 & -0.966335 & 0.000 \\
0.4 & -0.904361 & -0.914150 & -0.914150 & 0.000 & 1.8 & -0.828848 & -0.961817 & -0.961817 & 0.000 \\
0.45 & -0.987513 & -0.998416 & -0.998416 & 0.000 & 1.85 & -0.816842 & -0.957833 & -0.957833 & 0.000 \\
0.5 & -1.042996 & -1.055160 & -1.055160 & 0.000 & 1.9 & -0.805333 & -0.954339 & -0.954339 & 0.000 \\
0.55 & -1.079051 & -1.092630 & -1.092630 & 0.000 & 1.95 & -0.794318 & -0.951290 & -0.951290 & 0.000 \\
0.6 & -1.101128 & -1.116286 & -1.116286 & 0.000 & 2 & -0.783793 & -0.948641 & -0.948641 & 0.000 \\
0.65 & -1.112997 & -1.129905 & -1.129905 & 0.000 & 2.05 & -0.773749 & -0.946350 & -0.946350 & 0.000 \\
0.7 & -1.117349 & -1.136189 & -1.136189 & 0.000 & 2.1 & -0.764178 & -0.944375 & -0.944375 & 0.000 \\
0.75 & -1.116151 & -1.137117 & -1.137117 & 0.000 & 2.15 & -0.755066 & -0.942678 & -0.942678 & 0.000 \\
0.8 & -1.110850 & -1.134148 & -1.134148 & 0.000 & 2.2 & -0.746401 & -0.941224 & -0.941224 & 0.000 \\
0.85 & -1.102511 & -1.128362 & -1.128362 & 0.000 & 2.25 & -0.738169 & -0.939982 & -0.939982 & 0.000 \\
0.9 & -1.091914 & -1.120560 & -1.120560 & 0.000 & 2.3 & -0.730353 & -0.938922 & -0.938922 & 0.000 \\
0.95 & -1.079637 & -1.111339 & -1.111339 & 0.000 & 2.35 & -0.722939 & -0.938021 & -0.938021 & 0.000 \\
1 & -1.066109 & -1.101150 & -1.101150 & 0.000 & 2.4 & -0.715910 & -0.937255 & -0.937255 & 0.000 \\
1.05 & -1.051657 & -1.090342 & -1.090342 & 0.000 & 2.45 & -0.709250 & -0.936605 & -0.936605 & 0.000 \\
1.1 & -1.036539 & -1.079193 & -1.079193 & 0.000 & 2.5 & -0.702944 & -0.936055 & -0.936055 & 0.000 \\
1.15 & -1.020964 & -1.067930 & -1.067930 & 0.000 & 2.55 & -0.696975 & -0.935589 & -0.935589 & 0.000 \\
1.2 & -1.005107 & -1.056741 & -1.056741 & 0.000 & 2.6 & -0.691328 & -0.935196 & -0.935196 & 0.000 \\
1.25 & -0.989114 & -1.045783 & -1.045783 & 0.000 & 2.65 & -0.685988 & -0.934864 & -0.934864 & 0.000 \\
1.3 & -0.973111 & -1.035186 & -1.035186 & 0.000 & 2.7 & -0.680941 & -0.934584 & -0.934584 & 0.000 \\
1.35 & -0.957203 & -1.025054 & -1.025054 & 0.000 & 2.75 & -0.676172 & -0.934349 & -0.934349 & 0.000 \\
1.4 & -0.941481 & -1.015468 & -1.015468 & 0.000 & 2.8 & -0.671669 & -0.934151 & -0.934151 & 0.000 \\
1.45 & -0.926017 & -1.006487 & -1.006487 & 0.000 & 2.85 & -0.667417 & -0.933985 & -0.933985 & 0.000 \\
1.5 & -0.910874 & -0.998149 & -0.998149 & 0.000 & 2.9 & -0.663405 & -0.933846 & -0.933846 & 0.000 \\
1.55 & -0.896099 & -0.990476 & -0.990476 & 0.000 & 2.95 & -0.659619 & -0.933729 & -0.933729 & 0.000 \\
1.6 & -0.881732 & -0.983473 & -0.983473 & 0.000 & 3 & -0.656048 & -0.933632 & -0.933632 & 0.000 \\
1.65 & -0.867804 & -0.977130 & -0.977130 & 0.000 &  & & & &  \\
\bottomrule
\end{tabular}}
\end{table*}

\paragraph{Setup.}
\ce{H2} has two electrons in two spin-orbitals, so it needs no active-space
truncation at all: the STO-3G calculation \emph{is} the full configuration
interaction. The bond was scanned from $0.30$ to $3.00\Ang$ in $0.05\Ang$
steps, $55$ geometries, by a wide margin the densest sampling in the
campaign, affordable only because each point costs a fraction of a second.

\paragraph{Variational results.}
Table~\ref{tab:data:h2:bond} is the complete scan. The variational solution
is essentially exact: the largest deviation from exact diagonalisation
anywhere on the curve is $3\times10^{-4}\mHa$, roughly $5000$ times below the
chemical-accuracy threshold, and the mean rounds to zero at four decimal
places. The energy minimum falls at $r=0.75\Ang$ with $E=-1.137117\Ha$,
matching the accepted STO-3G result.

The correlation column is the interesting one pedagogically. Near
equilibrium, Hartree--Fock misses only a few tens of milli-Hartree, but by
$r=3.0\Ang$ the deficit has grown to $277.6\mHa$. This is the textbook
failure of restricted Hartree--Fock at dissociation: the single-determinant
reference forces both electrons to share a delocalised molecular orbital even
as the atoms separate, which describes two isolated hydrogen atoms poorly.
The variational treatment tracks the exact curve throughout, exactly the
regime where correlated methods earn their cost.


\paragraph{Reconstruction.}
None was done. The repository has no \texttt{h2\_polynomial/} or
\texttt{h2\_spline/} directory, so \ce{H2} contributes no rows to
Table~\ref{tab:interp} and no figures to this report. The gap is defensible
on its own terms, with $55$ nodes at $0.05\Ang$ spacing there is not much for
an interpolant to recover, and the piecewise-linear reconstruction is already
adequate, but it does mean the densest scan in the campaign is the one scan
with no interpolation evidence behind it. We flag it as a gap rather than
paper over it.

% ---------------------------------------------------------------------
\MoleculeBarrier
\subsubsection{\ce{LiH}: lithium hydride, and a negative result}
\label{sec:lih}

\paragraph{Setup.}
\ce{LiH} has four electrons, two of which are correlated in a
$(2\text{e},5\text{o})$ active space over eight qubits with a $24$-parameter
ansatz. The bond was scanned from $0.9$ to $3.9\Ang$ in $0.25\Ang$ steps,
$13$ nodes, a deliberately coarse grid, and optimised with SLSQP, which
converged in an average of $167$ energy evaluations.

\paragraph{Variational results.}
Table~\ref{tab:data:lih:bond} gives the scan. Accuracy is excellent: mean
variational error $0.0001\mHa$, maximum $0.0006\mHa$, against a correlation
energy that reaches $153.2\mHa$ at the stretched end. The lowest computed
node is $-7.880947\Ha$ at $r=1.65\Ang$. The efficiency is notable too:
$167$ evaluations on a $24$-parameter landscape converges comfortably, and it
is the reason this system is accurate where the deeper-circuit molecules of
Sections~\ref{sec:n2} and \ref{sec:o2} are not.

\figone{lih_polynomial_curve.png}{\ce{LiH} dissociation curve reconstructed by
a single global polynomial through all $13$ nodes (degree $12$). The fit
passes through every computed point by construction, with an in-sample
residual of $1.4\times10^{-11}\mHa$, machine precision. That residual is
\emph{not} evidence of accuracy, which is the whole point of the
out-of-sample analysis in Fig.~\ref{fig:lih_poly_err}.}{lih_poly_curve}

\figone{lih_polynomial_degree_sweep.png}{\ce{LiH} polynomial degree sweep:
leave-one-out error against fitted degree, tabulated in
Table~\ref{tab:deg:lih:bond}. The Runge signature we expected, a minimum
followed by a rise, does not appear. Error falls essentially monotonically to
the highest testable degree.}{lih_poly_sweep}
% GENERATED by gen_tables.py -- do not edit by hand (tab:deg:lih:bond)
\begin{table}[!htbp]
\caption{\ce{LiH} bond length polynomial degree sweep: leave-one-out error (mHa) against fitted degree $d$; the best degree is in bold.}
\label{tab:deg:lih:bond}
\centering\footnotesize\setlength{\tabcolsep}{3pt}
\adjustbox{max width=\columnwidth}{%
\begin{tabular}{rrrr}
\toprule
Degree $d$ & max & RMS & MAE \\
\midrule
1 & 58.6150 & 30.7998 & 24.0411 \\
2 & 58.4160 & 31.6399 & 27.4945 \\
3 & 33.8929 & 18.7903 & 16.6363 \\
4 & 16.5561 & 7.3950 & 6.3092 \\
5 & 7.4935 & 3.1427 & 2.4570 \\
6 & 3.9421 & 1.7138 & 1.3773 \\
7 & 1.6812 & 0.7436 & 0.5334 \\
8 & 0.5390 & 0.2104 & 0.1416 \\
9 & 0.5286 & 0.2088 & 0.1290 \\
10 & 0.5323 & 0.2307 & 0.1321 \\
\textbf{11} & \textbf{0.1091} & \textbf{0.0474} & \textbf{0.0254} \\
\bottomrule
\end{tabular}}
\end{table}



\figone{lih_polynomial_error.png}{\ce{LiH} leave-one-out error at each
interior node for the polynomial and for the piecewise-linear baseline. The
polynomial sits below the baseline everywhere, by two to three orders of
magnitude.}{lih_poly_err}


% GENERATED by gen_tables.py -- do not edit by hand (tab:data:lih:bond)
\begin{table}[!htbp]
\caption{\ce{LiH} bond length scan: Hartree--Fock, VQE and exact energies (Ha) at every geometry; $\delta = (E_{\mathrm{VQE}}-E_{\mathrm{exact}})\times 10^3$ (mHa).}
\label{tab:data:lih:bond}
\centering\footnotesize\setlength{\tabcolsep}{3pt}
\adjustbox{max width=\columnwidth}{%
\begin{tabular}{rrrrr}
\toprule
$r$ (\AA) & $E_{\mathrm{HF}}$ & $E_{\mathrm{VQE}}$ & $E_{\mathrm{exact}}$ & $\delta$ \\
\midrule
0.9 & -7.705753 & -7.722835 & -7.722835 & 0.000 \\
1.15 & -7.823711 & -7.840175 & -7.840175 & 0.000 \\
1.4 & -7.860539 & -7.878230 & -7.878231 & 0.001 \\
1.65 & -7.859852 & -7.880947 & -7.880947 & 0.000 \\
1.9 & -7.841112 & -7.867989 & -7.867989 & 0.000 \\
2.15 & -7.813947 & -7.849305 & -7.849305 & 0.000 \\
2.4 & -7.783382 & -7.830343 & -7.830343 & 0.000 \\
2.65 & -7.752218 & -7.814172 & -7.814172 & 0.000 \\
2.9 & -7.722219 & -7.802147 & -7.802147 & 0.000 \\
3.15 & -7.694586 & -7.794129 & -7.794129 & 0.000 \\
3.4 & -7.670060 & -7.789145 & -7.789145 & 0.000 \\
3.65 & -7.648946 & -7.786160 & -7.786160 & 0.000 \\
3.9 & -7.631193 & -7.784406 & -7.784406 & 0.000 \\
\bottomrule
\end{tabular}}
\end{table}


\paragraph{The result the experiment did not expect.}
The \ce{LiH} polynomial experiment was written to reproduce Runge's
phenomenon. It did not, and saying so plainly is more useful than shipping
the story we expected to tell.

Table~\ref{tab:deg:lih:bond} and Fig.~\ref{fig:lih_poly_sweep} show
leave-one-out RMS falling from $31.64\mHa$ at degree $2$ to $0.047\mHa$ at
degree $11$, the \emph{highest degree the $13$-node scan can test}, with no
rise anywhere. The full degree-$12$ interpolant has exactly one interior
minimum, so it invents no spurious features at all. The headline comparison
comes out backwards from what the textbook predicts:
\begin{center}
\begin{tabular}{lr}
\toprule
Method & LOO RMS \\
\midrule
Piecewise linear      & $13.373\mHa$ \\
Cubic spline          & $1.815\mHa$ \\
Global polynomial ($d=11$) & $\mathbf{0.047\mHa}$ \\
\bottomrule
\end{tabular}
\end{center}
The global polynomial beats the cubic spline by a factor of $38$, and the
linear baseline by a factor of $282$.

\paragraph{Why.}
Runge's phenomenon is a property of the \emph{function being interpolated},
not an automatic consequence of high degree. It bites when the function has a
singularity in the complex plane close to the real interval, Runge's own
$1/(1+25x^2)$ has poles at $\pm i/5$, right next to $[-1,1]$. The \ce{LiH}
CAS$(2,5)$ curve over $[0.9, 3.9]\Ang$ is smooth and gently varying with no
nearby singularity, so a high-degree polynomial approximates it very well.
Equation~\eqref{eq:polyerror} agrees: when $f^{(n)}$ stays bounded, the
$1/n!$ factor wins.

\paragraph{Where the trap does start.}
The trap is real, it just begins at a higher node density. Re-running the
same interval at $0.125\Ang$ spacing, $25$ nodes instead of $13$, gives best
degree $13$ at RMS $0.0004\mHa$, degree $18$ at $0.0005\mHa$, and degree
$23$ at $0.0065\mHa$: error bottoms out and then climbs by a factor of $16$.
That is the Runge signature. The scale is worth noting, though: the whole
effect plays out below one hundredth of a milli-Hartree, a clean textbook
trend rather than the catastrophe the phenomenon is usually illustrated
with. Contrast this with the \ce{N2O} surface (Section~\ref{sec:n2o}), where
the polynomial failure is genuinely catastrophic.

\figpair{lih_spline_curve.png}{Cubic spline reconstruction of the \ce{LiH}
curve.}{lih_spline_curve}{lih_spline_error.png}{Spline and linear
leave-one-out error. The spline improves on the linear baseline by
$7.4\times$ in RMS, a solid result the polynomial nonetheless
beats.}{lih_spline_err}

\paragraph{Geometry.}
Both reconstructions locate the minimum in essentially the same place:
$1.5474\Ang$ (polynomial) and $1.5496\Ang$ (spline), against a best computed
node of $1.65\Ang$. The polynomial value reproduces the model's own optimised
bond length of $1.5474\Ang$ to four decimal places. From a $0.25\Ang$ grid,
the interpolation recovers a geometry accurate to the model's own reference,
roughly two orders of magnitude of resolution gained for zero extra quantum
cost. Both sit $0.045$--$0.048\Ang$ below the experimental $1.5949\Ang$, and
that gap belongs to the STO-3G/CAS$(2,5)$ model, not to the interpolation.

% ---------------------------------------------------------------------
\MoleculeBarrier
\subsubsection{\ce{BeH2}: beryllium hydride}
\label{sec:beh2}

\paragraph{Setup.}
\ce{BeH2} has six electrons in a $(4\text{e},4\text{o})$ active space, six
qubits, $26$ parameters. It is the most thoroughly explored system in the
campaign: a symmetric stretch ($16$ nodes, $0.9$--$2.4\Ang$), a bend ($19$
nodes, $90$--$180^\circ$) and a full two-dimensional surface on an
$11\times10$ grid, $110$ variational solves, the joint-largest surface here.

\figwide{beh2_surface_anim_still.png}{\ce{BeH2} ground-state energy surface
(animation frame). $110$ VQE energies rise out of a bond-length $\times$
bond-angle grid in scan order while the molecular rendering redraws to match.
The trough runs along the angle axis: bending is soft, stretching is stiff.
The reconstruction analysis below turns this into a number.}{beh2_anim}

\figwide{beh2_surface_scan_3d.png}{\ce{BeH2} completed energy surface, static
render. The minimum sits at the linear geometry $\theta = 180^\circ$, as
expected for a $\mathrm{D}_{\infty h}$ molecule.}{beh2_surface3d}

\paragraph{Variational results.}
The bend scan is excellent: mean error $0.0044\mHa$, maximum $0.0062\mHa$,
four orders of magnitude inside chemical accuracy
(Table~\ref{tab:data:beh2:bend}). The stretch scan is noticeably worse, mean
$0.5595\mHa$ and maximum $3.4778\mHa$ (Table~\ref{tab:data:beh2:stretch}),
and the surface inherits this at mean $0.4541\mHa$, maximum $4.6537\mHa$.
The pattern makes physical sense: bending at a fixed bond length barely
touches the electronic structure, while stretching towards dissociation
drives the system multireference, where a $26$-parameter ansatz started at
Hartree--Fock struggles most. The correlation energy tells the same story,
rising from $5.9\mHa$ across the bend to $68.1\mHa$ across the stretch.


% GENERATED by gen_tables.py -- do not edit by hand (tab:data:beh2:bend)
\begin{table}[!htbp]
\caption{\ce{BeH2} bend angle scan: Hartree--Fock, VQE and exact energies (Ha) at every geometry; $\delta = (E_{\mathrm{VQE}}-E_{\mathrm{exact}})\times 10^3$ (mHa).}
\label{tab:data:beh2:bend}
\centering\footnotesize\setlength{\tabcolsep}{3pt}
\adjustbox{max width=\columnwidth}{%
\begin{tabular}{rrrrr}
\toprule
$\theta$ (deg) & $E_{\mathrm{HF}}$ & $E_{\mathrm{VQE}}$ & $E_{\mathrm{exact}}$ & $\delta$ \\
\midrule
90 & -15.471340 & -15.475542 & -15.475547 & 0.004 \\
95 & -15.482637 & -15.487013 & -15.487018 & 0.005 \\
100 & -15.492809 & -15.497363 & -15.497368 & 0.005 \\
105 & -15.501972 & -15.506701 & -15.506707 & 0.006 \\
110 & -15.510230 & -15.515127 & -15.515133 & 0.006 \\
115 & -15.517671 & -15.522726 & -15.522732 & 0.006 \\
120 & -15.524368 & -15.529570 & -15.529576 & 0.006 \\
125 & -15.530383 & -15.535720 & -15.535726 & 0.006 \\
130 & -15.535765 & -15.541223 & -15.541229 & 0.006 \\
135 & -15.540554 & -15.546118 & -15.546124 & 0.006 \\
140 & -15.544781 & -15.550437 & -15.550442 & 0.005 \\
145 & -15.548470 & -15.554201 & -15.554207 & 0.005 \\
150 & -15.551640 & -15.557431 & -15.557436 & 0.004 \\
155 & -15.554305 & -15.560140 & -15.560144 & 0.004 \\
160 & -15.556475 & -15.562341 & -15.562343 & 0.003 \\
165 & -15.558157 & -15.564041 & -15.564042 & 0.002 \\
170 & -15.559355 & -15.565249 & -15.565250 & 0.001 \\
175 & -15.560073 & -15.565971 & -15.565972 & 0.001 \\
180 & -15.560312 & -15.566211 & -15.566212 & 0.001 \\
\bottomrule
\end{tabular}}
\end{table}


% GENERATED by gen_tables.py -- do not edit by hand (tab:data:beh2:stretch)
\begin{table}[!htbp]
\caption{\ce{BeH2} symmetric stretch scan: Hartree--Fock, VQE and exact energies (Ha) at every geometry; $\delta = (E_{\mathrm{VQE}}-E_{\mathrm{exact}})\times 10^3$ (mHa).}
\label{tab:data:beh2:stretch}
\centering\footnotesize\setlength{\tabcolsep}{3pt}
\adjustbox{max width=\columnwidth}{%
\begin{tabular}{rrrrr}
\toprule
$r$ (\AA) & $E_{\mathrm{HF}}$ & $E_{\mathrm{VQE}}$ & $E_{\mathrm{exact}}$ & $\delta$ \\
\midrule
0.9 & -15.338455 & -15.346488 & -15.346490 & 0.001 \\
1 & -15.455668 & -15.462870 & -15.462871 & 0.001 \\
1.1 & -15.521999 & -15.528632 & -15.528633 & 0.001 \\
1.2 & -15.553586 & -15.559818 & -15.559819 & 0.001 \\
1.3 & -15.561278 & -15.567233 & -15.567233 & 0.001 \\
1.4 & -15.552460 & -15.558235 & -15.558237 & 0.002 \\
1.5 & -15.532213 & -15.537897 & -15.537897 & 0.001 \\
1.6 & -15.504086 & -15.509753 & -15.509754 & 0.001 \\
1.7 & -15.470606 & -15.476333 & -15.476334 & 0.001 \\
1.8 & -15.433626 & -15.439488 & -15.439488 & 0.001 \\
1.9 & -15.394538 & -15.400613 & -15.400614 & 0.001 \\
2 & -15.354417 & -15.384666 & -15.385345 & 0.678 \\
2.1 & -15.314113 & -15.350025 & -15.351020 & 0.996 \\
2.2 & -15.274312 & -15.317412 & -15.318911 & 1.499 \\
2.3 & -15.235575 & -15.287969 & -15.290260 & 2.291 \\
2.4 & -15.198367 & -15.262967 & -15.266445 & 3.478 \\
\bottomrule
\end{tabular}}
\end{table}


\paragraph{Bend reconstruction.}
The bend is the easiest reconstruction problem in the campaign. The curve is
smooth and nearly harmonic, and every method does well
(Table~\ref{tab:interp}): linear $0.344\mHa$ RMS, the polynomial at best
degree $9$ reaching $0.0001\mHa$, the spline $0.0008\mHa$. The degree sweep
(Fig.~\ref{fig:beh2_bend_sweep}, Table~\ref{tab:deg:beh2:bend}) shows the
classic bowl shape, with a worst-case RMS of $7.59\mHa$ at low degree,
underfitting rather than oscillation.

\figpair{beh2_bend_polynomial_curve.png}{\ce{BeH2} bend, polynomial
reconstruction.}{beh2_bend_poly_curve}{beh2_bend_polynomial_degree_sweep.png}%
{\ce{BeH2} bend degree sweep. Error falls to $0.0001\mHa$ at degree $9$.}%
{beh2_bend_sweep}
% GENERATED by gen_tables.py -- do not edit by hand (tab:deg:beh2:bend)
\begin{table}[!htbp]
\caption{\ce{BeH2} bend angle polynomial degree sweep: leave-one-out error (mHa) against fitted degree $d$; the best degree is in bold.}
\label{tab:deg:beh2:bend}
\centering\footnotesize\setlength{\tabcolsep}{3pt}
\adjustbox{max width=\columnwidth}{%
\begin{tabular}{rrrr}
\toprule
Degree $d$ & max & RMS & MAE \\
\midrule
1 & 13.1025 & 7.5881 & 6.7476 \\
2 & 1.0822 & 0.6700 & 0.6044 \\
3 & 0.2973 & 0.1796 & 0.1579 \\
4 & 0.0540 & 0.0334 & 0.0296 \\
5 & 0.0043 & 0.0018 & 0.0015 \\
6 & 0.0040 & 0.0016 & 0.0013 \\
7 & 0.0011 & 0.0004 & 0.0003 \\
8 & 0.0007 & 0.0003 & 0.0002 \\
\textbf{9} & \textbf{0.0002} & \textbf{0.0001} & \textbf{0.0001} \\
10 & 0.0004 & 0.0001 & 0.0001 \\
11 & 0.0005 & 0.0002 & 0.0001 \\
12 & 0.0005 & 0.0002 & 0.0001 \\
13 & 0.0005 & 0.0002 & 0.0001 \\
14 & 0.0038 & 0.0011 & 0.0005 \\
15 & 0.0351 & 0.0093 & 0.0036 \\
16 & 0.1607 & 0.0493 & 0.0196 \\
17 & 0.5580 & 0.1928 & 0.0755 \\
\bottomrule
\end{tabular}}
\end{table}



\figpair{beh2_bend_polynomial_error.png}{\ce{BeH2} bend, polynomial versus
linear leave-one-out error.}{beh2_bend_poly_err}{beh2_bend_spline_curve.png}%
{\ce{BeH2} bend, cubic spline reconstruction.}{beh2_bend_spline_curve}

\figone{beh2_bend_spline_error.png}{\ce{BeH2} bend, spline leave-one-out
error. RMS $0.0008\mHa$ against a linear baseline of $0.344\mHa$, a
$430\times$ improvement on an already easy problem.}{beh2_bend_spline_err}

\paragraph{Stretch reconstruction, and a spurious minimum.}
The stretch is harder, and it gives the campaign's clearest example of why a
good fit statistic does not mean a usable curve. Polynomial and spline are
almost tied statistically, $4.496\mHa$ versus $4.647\mHa$ RMS, both roughly
$2.2\times$ better than the $10.12\mHa$ linear baseline. But the recovered
geometries disagree completely: the spline places the minimum at
$1.289\Ang$, close to the reference $1.3264\Ang$, while the full-degree
polynomial places it at $0.9305\Ang$, at the extreme left edge of the
scanned interval, $0.40\Ang$ off, and physically meaningless.

This is Runge's phenomenon caught in the act. The degree-$15$ interpolant
oscillates near the interval boundary, and the oscillation dips below the
true curve at the edge, creating a spurious global minimum. The degree sweep
(Table~\ref{tab:deg:beh2:stretch}) confirms the mechanism: worst-case RMS
reaches $572.5\mHa$, four orders of magnitude above the best degree. The
lesson generalises to every polynomial row in Table~\ref{tab:geom}: an
aggregate error metric can look fine while the interpolant is qualitatively
wrong exactly where someone would read it.

\figpair{beh2_stretch_polynomial_curve.png}{\ce{BeH2} symmetric stretch,
polynomial reconstruction. Note the behaviour at the left edge, which
produces the spurious minimum discussed in the
text.}{beh2_str_poly_curve}{beh2_stretch_polynomial_degree_sweep.png}%
{\ce{BeH2} stretch degree sweep. Worst-case RMS reaches $572.5\mHa$.}%
{beh2_str_sweep}
% GENERATED by gen_tables.py -- do not edit by hand (tab:deg:beh2:stretch)
\begin{table}[!htbp]
\caption{\ce{BeH2} symmetric stretch polynomial degree sweep: leave-one-out error (mHa) against fitted degree $d$; the best degree is in bold.}
\label{tab:deg:beh2:stretch}
\centering\footnotesize\setlength{\tabcolsep}{3pt}
\adjustbox{max width=\columnwidth}{%
\begin{tabular}{rrrr}
\toprule
Degree $d$ & max & RMS & MAE \\
\midrule
1 & 91.4818 & 59.4448 & 50.2601 \\
2 & 59.9675 & 36.3232 & 31.4312 \\
3 & 21.7420 & 14.0259 & 12.8937 \\
4 & 10.4896 & 4.6810 & 3.7688 \\
5 & 11.6111 & 4.6444 & 3.4286 \\
6 & 12.0656 & 5.1597 & 3.8667 \\
\textbf{7} & \textbf{10.6220} & \textbf{4.4964} & \textbf{3.2178} \\
8 & 15.2741 & 6.1943 & 3.8921 \\
9 & 18.9213 & 8.2993 & 5.8991 \\
10 & 19.1465 & 7.5641 & 5.7692 \\
11 & 109.0775 & 31.7178 & 15.6726 \\
12 & 334.6386 & 107.3525 & 51.8664 \\
13 & 763.7332 & 268.9519 & 124.8257 \\
14 & 1498.5750 & 572.5081 & 256.2855 \\
\bottomrule
\end{tabular}}
\end{table}



\figpair{beh2_stretch_polynomial_error.png}{\ce{BeH2} stretch, polynomial
versus linear leave-one-out
error.}{beh2_str_poly_err}{beh2_stretch_spline_curve.png}{\ce{BeH2} stretch,
cubic spline reconstruction. The spline stays monotone at the edges and
places the minimum at $1.289\Ang$.}{beh2_str_spline_curve}

\figone{beh2_stretch_spline_error.png}{\ce{BeH2} stretch, spline
leave-one-out error, RMS $4.647\mHa$.}{beh2_str_spline_err}

\paragraph{Surface reconstruction.}
The half-grid holdout earns its keep here. Fitting on a $6\times5$ subgrid,
one quarter of the $110$ computed points, and scoring on the remaining three
quarters gives bicubic spline RMS $3.098\mHa$ against a bilinear baseline of
$17.97\mHa$, a $5.8\times$ improvement. The total-degree polynomial reaches
$3.966\mHa$ but, unlike the 1-D case, carries a nonzero in-sample residual of
$8.58\mHa$: it cannot represent the surface exactly even on the points it was
fitted to.

The practical reading is direct: \emph{we could have run a quarter of this
surface and reconstructed the rest to within about $3\mHa$ RMS}. At two to
four minutes per point, that saves several hours, and it is the result that
justifies coarse grids on more expensive molecules.

\figwide{beh2_surface_polynomial_surface.png}{\ce{BeH2} surface, total-degree
polynomial reconstruction.}{beh2_surf_poly}
\figwide{beh2_surface_polynomial_error.png}{\ce{BeH2} surface, polynomial holdout error distribution.}{beh2_surf_poly_err}

\figwide{beh2_surface_spline_surface.png}{\ce{BeH2} surface, bicubic spline
reconstruction from the half grid.}{beh2_surf_spline}
\figwide{beh2_surface_spline_error.png}{\ce{BeH2} surface, spline holdout error.
RMS $3.098\mHa$ against $17.97\mHa$ bilinear.}{beh2_surf_spline_err}

\paragraph{Geometry.}
The spline surface locates the minimum at $(1.292\Ang, 180.0^\circ)$ against
the reference $(1.3264\Ang, 180.0^\circ)$: the angle exactly right, the bond
$0.034\Ang$ short, recovered from a grid whose bond spacing is $0.15\Ang$.
The polynomial surface gives $(1.2995\Ang,\allowbreak 176.38^\circ)$, missing the linear
geometry by nearly four degrees, a qualitative error for a molecule whose
linearity is a symmetry requirement.

% ---------------------------------------------------------------------
\MoleculeBarrier
\subsubsection{\ce{H2O}: water}
\label{sec:h2o}

\paragraph{Setup.}
\ce{H2O} has ten electrons in a $(4\text{e},4\text{o})$ active space, six
qubits, $26$ parameters. The repository holds one scan: the H--O--H bend from
$70$ to $160^\circ$ in $5^\circ$ steps at a fixed bond length, $19$
geometries. A symmetric-stretch coordinate is implemented in the driver and
documented in the project README, but no saved stretch scan is present, so we
cannot report on it here.

\paragraph{Variational results.}
Table~\ref{tab:data:h2o:bend} gives the complete scan. Accuracy is among the
best in the campaign: mean error $0.0007\mHa$, maximum $0.0015\mHa$, three
orders of magnitude inside chemical accuracy. The minimum falls at $\theta =
100^\circ$ with $E = -74.971353\Ha$. The experimental H--O--H angle is
$104.5^\circ$; the $4.5^\circ$ gap belongs to the minimal basis and small
active space, and no improvement to the solver would remove it. Correlation
energy across the scan reaches $11.1\mHa$, small, because bending at a fixed
bond length is a gentle perturbation of the electronic structure, exactly as
we saw for \ce{BeH2}.


% GENERATED by gen_tables.py -- do not edit by hand (tab:data:h2o:bend)
\begin{table}[!htbp]
\caption{\ce{H2O} bend angle scan: Hartree--Fock, VQE and exact energies (Ha) at every geometry; $\delta = (E_{\mathrm{VQE}}-E_{\mathrm{exact}})\times 10^3$ (mHa).}
\label{tab:data:h2o:bend}
\centering\footnotesize\setlength{\tabcolsep}{3pt}
\adjustbox{max width=\columnwidth}{%
\begin{tabular}{rrrrr}
\toprule
$\theta$ (deg) & $E_{\mathrm{HF}}$ & $E_{\mathrm{VQE}}$ & $E_{\mathrm{exact}}$ & $\delta$ \\
\midrule
70 & -74.912380 & -74.923506 & -74.923507 & 0.000 \\
75 & -74.928523 & -74.938913 & -74.938913 & 0.000 \\
80 & -74.941332 & -74.951073 & -74.951073 & 0.000 \\
85 & -74.951034 & -74.960204 & -74.960204 & 0.000 \\
90 & -74.957841 & -74.966501 & -74.966502 & 0.000 \\
95 & -74.961953 & -74.970155 & -74.970156 & 0.000 \\
100 & -74.963568 & -74.971353 & -74.971353 & 0.000 \\
105 & -74.962882 & -74.970280 & -74.970280 & 0.000 \\
110 & -74.960090 & -74.967125 & -74.967126 & 0.001 \\
115 & -74.955392 & -74.962083 & -74.962084 & 0.001 \\
120 & -74.948995 & -74.955358 & -74.955359 & 0.001 \\
125 & -74.941120 & -74.947170 & -74.947170 & 0.001 \\
130 & -74.932011 & -74.937762 & -74.937763 & 0.001 \\
135 & -74.921941 & -74.927413 & -74.927413 & 0.001 \\
140 & -74.911219 & -74.916437 & -74.916438 & 0.001 \\
145 & -74.900201 & -74.905201 & -74.905202 & 0.001 \\
150 & -74.889289 & -74.894117 & -74.894119 & 0.001 \\
155 & -74.878927 & -74.883640 & -74.883641 & 0.001 \\
160 & -74.869584 & -74.874243 & -74.874244 & 0.001 \\
\bottomrule
\end{tabular}}
\end{table}


\paragraph{Reconstruction.}
As with \ce{H2}, none was done: the repository has no
\texttt{h2o\_polynomial/} or \texttt{h2o\_spline/} directory. \ce{H2O}
therefore contributes no rows to Table~\ref{tab:interp} and no figures. Since
\ce{H2O} is the molecule named in the project README as the worked example,
this is the most conspicuous gap in the campaign, and we flag it rather than
let it pass quietly.

% ---------------------------------------------------------------------
\MoleculeBarrier
\subsubsection{\ce{NH3}: ammonia}
\label{sec:nh3}

\paragraph{Setup.}
\ce{NH3} has ten electrons in a $(6\text{e},5\text{o})$ active space, eight
qubits, $54$ parameters, $\mathrm{C}_{3v}$ symmetry with two free
coordinates. SLSQP converged in $255$--$348$ evaluations depending on the
scan. Three scans were run: an umbrella bend ($9$ nodes,
$96$--$120^\circ$), a symmetric N--H stretch ($11$ nodes, $0.85$--$1.6\Ang$)
and a $7\times7$ surface ($49$ points).

\ce{NH3} is the one molecule whose raw scan plots survive in the repository
alongside the interpolation figures, so Figs.~\ref{fig:nh3_bend_raw} to
\ref{fig:nh3_surf_raw} show the campaign's output in its original form.

\figpair{nh3_bend_curve.png}{\ce{NH3} umbrella bend, raw VQE scan.}%
{nh3_bend_raw}{nh3_stretch_curve.png}{\ce{NH3} symmetric N--H stretch, raw
VQE scan.}{nh3_str_raw}

\figwide{nh3_surface.png}{\ce{NH3} raw energy surface over the
bond-length $\times$ umbrella-angle grid, $49$ VQE
solves.}{nh3_surf_raw}

\paragraph{Variational results.}
All three scans reach chemical accuracy. The bend attains mean $0.0158\mHa$
(maximum $0.0494$), the stretch mean $0.1937\mHa$ (maximum $0.7605$), and the
surface mean $0.3343\mHa$, although the surface maximum, $4.5362\mHa$,
exceeds the threshold at its worst point, presumably the most stretched
corner of the grid. The stretch-versus-bend contrast mirrors \ce{BeH2}
exactly.


% GENERATED by gen_tables.py -- do not edit by hand (tab:data:nh3:bend)
\begin{table}[!htbp]
\caption{\ce{NH3} bend angle scan: Hartree--Fock, VQE and exact energies (Ha) at every geometry; $\delta = (E_{\mathrm{VQE}}-E_{\mathrm{exact}})\times 10^3$ (mHa).}
\label{tab:data:nh3:bend}
\centering\footnotesize\setlength{\tabcolsep}{3pt}
\adjustbox{max width=\columnwidth}{%
\begin{tabular}{rrrrr}
\toprule
$\theta$ (deg) & $E_{\mathrm{HF}}$ & $E_{\mathrm{VQE}}$ & $E_{\mathrm{exact}}$ & $\delta$ \\
\midrule
96 & -55.446703 & -55.467822 & -55.467829 & 0.007 \\
99 & -55.450917 & -55.471373 & -55.471381 & 0.008 \\
102 & -55.453414 & -55.473341 & -55.473349 & 0.008 \\
105 & -55.454280 & -55.473791 & -55.473799 & 0.008 \\
108 & -55.453602 & -55.472728 & -55.472736 & 0.008 \\
111 & -55.451472 & -55.470284 & -55.470295 & 0.011 \\
114 & -55.447996 & -55.466601 & -55.466619 & 0.018 \\
117 & -55.443295 & -55.462119 & -55.462143 & 0.024 \\
120 & -55.437509 & -55.456768 & -55.456818 & 0.049 \\
\bottomrule
\end{tabular}}
\end{table}


% GENERATED by gen_tables.py -- do not edit by hand (tab:data:nh3:stretch)
\begin{table}[!htbp]
\caption{\ce{NH3} symmetric stretch scan: Hartree--Fock, VQE and exact energies (Ha) at every geometry; $\delta = (E_{\mathrm{VQE}}-E_{\mathrm{exact}})\times 10^3$ (mHa).}
\label{tab:data:nh3:stretch}
\centering\footnotesize\setlength{\tabcolsep}{3pt}
\adjustbox{max width=\columnwidth}{%
\begin{tabular}{rrrrr}
\toprule
$r$ (\AA) & $E_{\mathrm{HF}}$ & $E_{\mathrm{VQE}}$ & $E_{\mathrm{exact}}$ & $\delta$ \\
\midrule
0.85 & -55.320316 & -55.332929 & -55.332933 & 0.004 \\
0.925 & -55.416129 & -55.431383 & -55.431387 & 0.004 \\
1 & -55.452322 & -55.470863 & -55.470873 & 0.009 \\
1.075 & -55.449123 & -55.472133 & -55.472148 & 0.015 \\
1.15 & -55.419851 & -55.448498 & -55.448528 & 0.030 \\
1.225 & -55.373570 & -55.409923 & -55.409975 & 0.052 \\
1.3 & -55.316579 & -55.362264 & -55.362374 & 0.110 \\
1.375 & -55.253251 & -55.310751 & -55.310961 & 0.210 \\
1.45 & -55.186578 & -55.257680 & -55.258050 & 0.370 \\
1.525 & -55.118575 & -55.206832 & -55.207397 & 0.565 \\
1.6 & -55.050603 & -55.155814 & -55.156574 & 0.760 \\
\bottomrule
\end{tabular}}
\end{table}


\paragraph{The inversion barrier.}
\ce{NH3} carries a real physical observable that only a $C^2$ reconstruction
can deliver. The molecule tunnels between two pyramidal configurations
through a planar transition state at $120^\circ$, and the height of that
barrier is a classic spectroscopic quantity. From the spline reconstruction
of the bend curve, the barrier is
\begin{equation}
\Delta E_{\text{inv}} = 17.06\mHa \approx 10.7\,\text{kcal}\,\text{mol}^{-1},
\end{equation}
measured from the interpolated minimum at $104.38^\circ$ to the planar
geometry. The experimental value is roughly
$5.8\,\text{kcal}\,\text{mol}^{-1}$; the overestimate belongs to the minimal
basis and restricted active space rather than to the reconstruction. The
methodological point is what matters here: the barrier is read off the
\emph{interpolated} curve, because the planar geometry $120^\circ$ is an
endpoint of a $3^\circ$ grid and the minimum at $104.38^\circ$ falls between
nodes. A piecewise-linear reconstruction would have snapped the minimum to
$105^\circ$ and distorted the barrier along with it.

\figpair{nh3_bend_polynomial_curve.png}{\ce{NH3} bend, polynomial
reconstruction.}{nh3_bend_poly}{nh3_bend_polynomial_degree_sweep.png}%
{\ce{NH3} bend degree sweep; best degree $4$ at RMS $0.061\mHa$.}%
{nh3_bend_sweep}
% GENERATED by gen_tables.py -- do not edit by hand (tab:deg:nh3:bend)
\begin{table}[!htbp]
\caption{\ce{NH3} bend angle polynomial degree sweep: leave-one-out error (mHa) against fitted degree $d$; the best degree is in bold.}
\label{tab:deg:nh3:bend}
\centering\footnotesize\setlength{\tabcolsep}{3pt}
\adjustbox{max width=\columnwidth}{%
\begin{tabular}{rrrr}
\toprule
Degree $d$ & max & RMS & MAE \\
\midrule
1 & 4.9654 & 3.3600 & 3.1111 \\
2 & 0.6015 & 0.3645 & 0.3319 \\
3 & 0.1545 & 0.0973 & 0.0895 \\
\textbf{4} & \textbf{0.1172} & \textbf{0.0605} & \textbf{0.0431} \\
5 & 0.1596 & 0.0836 & 0.0682 \\
6 & 0.1489 & 0.0738 & 0.0593 \\
7 & 0.2826 & 0.1590 & 0.1199 \\
\bottomrule
\end{tabular}}
\end{table}



\figpair{nh3_bend_polynomial_error.png}{\ce{NH3} bend, polynomial versus
linear leave-one-out error.}{nh3_bend_poly_err}{nh3_bend_spline_curve.png}%
{\ce{NH3} bend, cubic spline reconstruction, the curve the inversion barrier
is measured from.}{nh3_bend_spline}

\figone{nh3_bend_spline_error.png}{\ce{NH3} bend, spline leave-one-out error.
RMS $0.067\mHa$ against a linear baseline of $0.653\mHa$.}{nh3_bend_spline_err}

\paragraph{Reconstruction.}
With only $9$ and $11$ nodes respectively, the \ce{NH3} curves are the
sparsest 1-D data in the campaign, and the results are correspondingly
instructive. On the bend, polynomial ($0.061\mHa$ at degree $4$) and spline
($0.067\mHa$) are statistically indistinguishable, and both beat linear
($0.653\mHa$) by an order of magnitude. On the stretch the polynomial wins
more clearly, $0.466\mHa$ at degree $6$ against $0.920\mHa$ for the spline
and $12.79\mHa$ for linear, a $27\times$ improvement over the baseline.

The surface is the hardest case. With a $7\times7$ grid, the half-grid
holdout leaves a $4\times4$ subgrid, the bare minimum that supports a bicubic
spline on both axes. Spline RMS is $8.28\mHa$ and total-degree polynomial
$8.49\mHa$, against a bilinear baseline of $30.11\mHa$: a $3.6\times$
improvement, but well above chemical accuracy in absolute terms. This is the
clearest evidence in the campaign that the half-grid experiment has a floor:
below roughly five points per axis there is not enough information left to
reconstruct a surface accurately, whichever method is used. We also did not
generate a half-grid error map for the \ce{NH3} surface the way we did for
\ce{BeH2}, \ce{NO2} and \ce{N2O}, so only the reconstructed surfaces
themselves are shown below, not an error plot.

\figpair{nh3_stretch_polynomial_curve.png}{\ce{NH3} stretch, polynomial
reconstruction.}{nh3_str_poly}{nh3_stretch_polynomial_degree_sweep.png}%
{\ce{NH3} stretch degree sweep; best degree $6$.}{nh3_str_sweep}
% GENERATED by gen_tables.py -- do not edit by hand (tab:deg:nh3:stretch)
\begin{table}[!htbp]
\caption{\ce{NH3} symmetric stretch polynomial degree sweep: leave-one-out error (mHa) against fitted degree $d$; the best degree is in bold.}
\label{tab:deg:nh3:stretch}
\centering\footnotesize\setlength{\tabcolsep}{3pt}
\adjustbox{max width=\columnwidth}{%
\begin{tabular}{rrrr}
\toprule
Degree $d$ & max & RMS & MAE \\
\midrule
1 & 79.9923 & 53.4177 & 48.1491 \\
2 & 43.8140 & 27.4471 & 24.3178 \\
3 & 12.9864 & 8.1997 & 7.2974 \\
4 & 3.3268 & 1.6579 & 1.4349 \\
5 & 1.8593 & 0.8965 & 0.6895 \\
\textbf{6} & \textbf{1.0562} & \textbf{0.4660} & \textbf{0.3318} \\
7 & 1.6842 & 0.7732 & 0.5741 \\
8 & 2.1305 & 0.8105 & 0.5427 \\
9 & 9.5652 & 4.6410 & 2.9185 \\
\bottomrule
\end{tabular}}
\end{table}



\figpair{nh3_stretch_polynomial_error.png}{\ce{NH3} stretch, polynomial
versus linear leave-one-out
error.}{nh3_str_poly_err}{nh3_stretch_spline_curve.png}{\ce{NH3} stretch,
cubic spline reconstruction.}{nh3_str_spline}

\figone{nh3_stretch_spline_error.png}{\ce{NH3} stretch, spline leave-one-out
error, RMS $0.920\mHa$.}{nh3_str_spline_err}

\figwide{nh3_surface_polynomial.png}{\ce{NH3} surface, total-degree
polynomial reconstruction (degree $6$).}{nh3_surf_poly}
\figwide{nh3_surface_spline.png}{\ce{NH3} surface, bicubic spline reconstruction
from the $4\times4$ half grid.}{nh3_surf_spline}
% GENERATED by gen_tables.py -- do not edit by hand (tab:deg:nh3:surface)
\begin{table}[!htbp]
\caption{\ce{NH3} 2-D surface polynomial degree sweep: leave-one-out error (mHa) against fitted degree $d$; the best degree is in bold.}
\label{tab:deg:nh3:surface}
\centering\footnotesize\setlength{\tabcolsep}{3pt}
\adjustbox{max width=\columnwidth}{%
\begin{tabular}{rrrr}
\toprule
Degree $d$ & max & RMS & MAE \\
\midrule
1 & 119.2770 & 74.4766 & 65.5689 \\
2 & 60.6062 & 34.3305 & 29.5216 \\
\textbf{3} & \textbf{16.0243} & \textbf{8.4926} & \textbf{7.0075} \\
\bottomrule
\end{tabular}}
\end{table}



\paragraph{Geometry.}
The reconstructions recover the model's own optimised geometry almost
exactly. The bend minimum sits at $104.38^\circ$ (both methods) against a
model reference of $104.27^\circ$, an error of $0.11^\circ$ from a grid
spaced at $3^\circ$. The stretch minimum sits at $1.0391$ (polynomial) and
$1.0394\Ang$ (spline) against a model reference of $1.0401\Ang$: errors of
$0.0010$ and $0.0007\Ang$ from a grid spaced at $0.075\Ang$, a hundredfold
gain in effective resolution and the single most convincing demonstration in
the campaign of what interpolation buys. The surface agrees, placing the
minimum at $(1.0454\Ang, 102.49^\circ)$ against the model surface reference
$(1.05\Ang, 102.05^\circ)$.

Against \emph{experiment}, $1.0124\Ang$ and $106.67^\circ$, the gaps are
larger, $0.027\Ang$ and $2.3^\circ$, and again these belong to the model, not
to anything interpolation could fix.

% ---------------------------------------------------------------------
\MoleculeBarrier
\subsubsection{\ce{N2}: nitrogen}
\label{sec:n2}

\paragraph{Setup.}
\ce{N2} has fourteen electrons in a $(6\text{e},6\text{o})$ active space, ten
qubits, and the largest ansatz in the campaign at $117$ parameters. Nine
nodes from $0.8$ to $2.4\Ang$ in $0.2\Ang$ steps, optimised with COBYLA under
a $1000$-evaluation ceiling.

\paragraph{Variational results.}
This is the campaign's hardest optimisation and its least accurate result:
mean error $10.48\mHa$, maximum $36.35\mHa$, an order of magnitude outside
chemical accuracy (Table~\ref{tab:data:n2:bond}).

The diagnosis is clear. COBYLA spends its full $1000$-evaluation budget at
every geometry, which on a $117$-dimensional landscape is fewer than nine
evaluations per parameter, not enough for a derivative-free method to map the
surface, let alone converge on it. The variational bound is never violated,
so the algorithm behaves correctly; it is simply stopped early. Triple-bond
breaking is also the most strongly multireference process in the campaign,
so the landscape being searched is the roughest one here.

Context matters, though: Hartree--Fock misses $777.0\mHa$ of correlation
energy on this curve, the largest deficit anywhere in the campaign. Missing
$10.5\mHa$ of that still means recovering $98.6\%$ of the correlation the
mean-field reference leaves out. The calculation is far from useless; it is
simply not converged to spectroscopic standards.


% GENERATED by gen_tables.py -- do not edit by hand (tab:data:n2:bond)
\begin{table}[!htbp]
\caption{\ce{N2} bond length scan: Hartree--Fock, VQE and exact energies (Ha) at every geometry; $\delta = (E_{\mathrm{VQE}}-E_{\mathrm{exact}})\times 10^3$ (mHa).}
\label{tab:data:n2:bond}
\centering\footnotesize\setlength{\tabcolsep}{3pt}
\adjustbox{max width=\columnwidth}{%
\begin{tabular}{rrrrr}
\toprule
$r$ (\AA) & $E_{\mathrm{HF}}$ & $E_{\mathrm{VQE}}$ & $E_{\mathrm{exact}}$ & $\delta$ \\
\midrule
0.8 & -106.680802 & -106.713261 & -106.714551 & 1.289 \\
1 & -107.419532 & -107.519020 & -107.520461 & 1.441 \\
1.2 & -107.487784 & -107.642667 & -107.645024 & 2.357 \\
1.4 & -107.357815 & -107.585331 & -107.591182 & 5.851 \\
1.6 & -107.184846 & -107.510796 & -107.513395 & 2.599 \\
1.8 & -107.017327 & -107.456615 & -107.459506 & 2.891 \\
2 & -106.871504 & -107.407238 & -107.437024 & 29.786 \\
2.2 & -106.751831 & -107.420818 & -107.432611 & 11.793 \\
2.4 & -106.656608 & -107.397268 & -107.433615 & 36.347 \\
\bottomrule
\end{tabular}}
\end{table}


\paragraph{Reconstruction.}
With only nine nodes at $0.2\Ang$ spacing across a bond-breaking curve, this
is the sparsest and most sharply curved 1-D reconstruction problem in the
campaign, and all the errors are correspondingly large. Linear interpolation
gives $134.18\mHa$ RMS, a value too large for any chemical purpose. The
polynomial at its best degree of $6$ reaches $20.34\mHa$ and the spline
$37.59\mHa$; the polynomial wins by a factor of $1.8$ and improves on the
baseline by $6.6\times$.

None of these reconstructions reaches chemical accuracy. Nine nodes across a
triple-bond dissociation is too coarse a grid for any interpolant to rescue,
regardless of which one comes out ahead in the comparison. The degree sweep (Table~\ref{tab:deg:n2:bond})
supports this: even the worst degree reaches only $207.5\mHa$, because with
seven testable degrees there is no room for the dramatic Runge failures seen
elsewhere.

\figpair{n2_polynomial_curve.png}{\ce{N2} dissociation curve, polynomial
reconstruction.}{n2_poly_curve}{n2_polynomial_degree_sweep.png}{\ce{N2}
degree sweep; best degree $6$ at $20.34\mHa$
RMS.}{n2_sweep}
% GENERATED by gen_tables.py -- do not edit by hand (tab:deg:n2:bond)
\begin{table}[!htbp]
\caption{\ce{N2} bond length polynomial degree sweep: leave-one-out error (mHa) against fitted degree $d$; the best degree is in bold.}
\label{tab:deg:n2:bond}
\centering\footnotesize\setlength{\tabcolsep}{3pt}
\adjustbox{max width=\columnwidth}{%
\begin{tabular}{rrrr}
\toprule
Degree $d$ & max & RMS & MAE \\
\midrule
1 & 362.6409 & 206.1169 & 171.3931 \\
2 & 377.8793 & 207.4911 & 172.1065 \\
3 & 253.2753 & 132.1034 & 111.9389 \\
4 & 117.4039 & 61.8030 & 53.5987 \\
5 & 81.6669 & 45.2077 & 36.6940 \\
\textbf{6} & \textbf{40.8968} & \textbf{20.3413} & \textbf{16.3684} \\
7 & 75.5146 & 42.4990 & 32.0552 \\
\bottomrule
\end{tabular}}
\end{table}



\figpair{n2_polynomial_error.png}{\ce{N2} polynomial versus linear
leave-one-out error.}{n2_poly_err}{n2_spline_curve.png}{\ce{N2} cubic spline
reconstruction.}{n2_spline_curve}

\figone{n2_spline_error.png}{\ce{N2} spline leave-one-out error, RMS
$37.59\mHa$. Compare the vertical scale with
Fig.~\ref{fig:beh2_bend_spline_err}: this is a two-orders-of-magnitude harder
problem.}{n2_spline_err}

\paragraph{Geometry.}
The polynomial places the minimum at $1.191\Ang$ and the spline at
$1.1685\Ang$, against an experimental bond length of $1.098\Ang$. Both
overestimate by $0.07$--$0.09\Ang$. Most of that error belongs to the minimal
basis and restricted active space, but with a $10.5\mHa$ mean variational
error sitting on top of a curve whose well depth is being read to
sub-milli-Hartree precision, some of it belongs to the incomplete
optimisation as well. This is the one molecule in the campaign where solver
error and model error cannot be cleanly separated, and we do not claim to
separate them.

% ---------------------------------------------------------------------
\MoleculeBarrier
\subsubsection{\ce{O2}: oxygen}
\label{sec:o2}

\paragraph{Setup.}
\ce{O2} has sixteen electrons in an $(8\text{e},6\text{o})$ active space with
five $\alpha$ and three $\beta$ electrons, ten qubits, $68$ parameters.
\ce{O2} is a ground-state \emph{triplet} ($2S = 2$), which needed a
restricted open-shell Hartree--Fock reference rather than the closed-shell
RHF used for most of the campaign. Nine nodes from $0.9$ to $2.5\Ang$ in
$0.2\Ang$ steps, COBYLA.

\paragraph{Variational results.}
Accuracy is poor for the same reason as \ce{N2}: mean $10.55\mHa$, maximum
$28.55\mHa$, with COBYLA saturating its $1000$-evaluation budget at every
geometry (Table~\ref{tab:data:o2:bond}). The open-shell reference adds a
further difficulty: the ansatz has to represent a state of definite spin
multiplicity, and nothing in the hardware-efficient construction enforces
that, so part of the optimiser's work goes into maintaining a symmetry that a
spin-adapted ansatz would get for free. Hartree--Fock misses $576.5\mHa$
here, so VQE again recovers roughly $98\%$ of the correlation energy despite
missing chemical accuracy.


% GENERATED by gen_tables.py -- do not edit by hand (tab:data:o2:bond)
\begin{table}[!htbp]
\caption{\ce{O2} bond length scan: Hartree--Fock, VQE and exact energies (Ha) at every geometry; $\delta = (E_{\mathrm{VQE}}-E_{\mathrm{exact}})\times 10^3$ (mHa).}
\label{tab:data:o2:bond}
\centering\footnotesize\setlength{\tabcolsep}{3pt}
\adjustbox{max width=\columnwidth}{%
\begin{tabular}{rrrrr}
\toprule
$r$ (\AA) & $E_{\mathrm{HF}}$ & $E_{\mathrm{VQE}}$ & $E_{\mathrm{exact}}$ & $\delta$ \\
\midrule
0.9 & -147.139439 & -147.171132 & -147.171162 & 0.030 \\
1.1 & -147.590894 & -147.656338 & -147.656442 & 0.103 \\
1.3 & -147.618812 & -147.735735 & -147.736896 & 1.161 \\
1.5 & -147.518566 & -147.697642 & -147.702491 & 4.849 \\
1.7 & -147.388737 & -147.637722 & -147.649223 & 11.501 \\
1.9 & -147.268103 & -147.601774 & -147.616729 & 14.955 \\
2.1 & -147.168643 & -147.579407 & -147.607961 & 28.555 \\
2.3 & -147.090786 & -147.594686 & -147.607685 & 12.999 \\
2.5 & -147.032001 & -147.587711 & -147.608502 & 20.790 \\
\bottomrule
\end{tabular}}
\end{table}


\paragraph{Reconstruction.}
The polynomial wins clearly: $7.99\mHa$ RMS at degree $6$, against
$22.00\mHa$ for the spline and $80.54\mHa$ for linear, improvements of
$10.1\times$ and $3.7\times$ on the baseline. As with \ce{N2}, nine nodes is
too sparse for any method to reach chemical accuracy, so read this comparison
as a ranking, not an endorsement.

The pattern across \ce{N2}, \ce{O2} and \ce{F2} is itself the finding: on
sparse, smooth, single-coordinate dissociation curves the global polynomial
outperforms the spline, and more decisively as the node count rises. With
nine nodes the margin is $2.8\times$ (\ce{O2}); with seventeen it becomes
$109\times$ (\ce{F2}, next section).

\figpair{o2_polynomial_curve.png}{\ce{O2} dissociation curve, polynomial
reconstruction.}{o2_poly_curve}{o2_polynomial_degree_sweep.png}{\ce{O2}
degree sweep; best degree $6$ at $7.99\mHa$ RMS, worst $130.17\mHa$.}%
{o2_sweep}
% GENERATED by gen_tables.py -- do not edit by hand (tab:deg:o2:bond)
\begin{table}[!htbp]
\caption{\ce{O2} bond length polynomial degree sweep: leave-one-out error (mHa) against fitted degree $d$; the best degree is in bold.}
\label{tab:deg:o2:bond}
\centering\footnotesize\setlength{\tabcolsep}{3pt}
\adjustbox{max width=\columnwidth}{%
\begin{tabular}{rrrr}
\toprule
Degree $d$ & max & RMS & MAE \\
\midrule
1 & 227.2053 & 126.7282 & 102.3330 \\
2 & 227.5155 & 130.1700 & 109.7675 \\
3 & 147.8719 & 80.2569 & 70.0870 \\
4 & 58.7578 & 30.7725 & 26.1430 \\
5 & 37.5126 & 21.5989 & 17.5765 \\
\textbf{6} & \textbf{16.2480} & \textbf{7.9871} & \textbf{6.3197} \\
7 & 69.0158 & 38.8415 & 29.2965 \\
\bottomrule
\end{tabular}}
\end{table}



\figpair{o2_polynomial_error.png}{\ce{O2} polynomial versus linear
leave-one-out error.}{o2_poly_err}{o2_spline_curve.png}{\ce{O2} cubic spline
reconstruction.}{o2_spline_curve}

\figone{o2_spline_error.png}{\ce{O2} spline leave-one-out error, RMS
$22.00\mHa$.}{o2_spline_err}

\paragraph{Geometry.}
Polynomial $1.3028\Ang$, spline $1.2754\Ang$, experiment $1.208\Ang$. The
spline is the closer of the two, $0.067\Ang$ above experiment. Both sit above
the true value, consistent with the general tendency of minimal-basis CAS
calculations to over-lengthen bonds.

% ---------------------------------------------------------------------
\MoleculeBarrier
\subsubsection{\ce{F2}: fluorine}
\label{sec:f2}

\paragraph{Setup.}
\ce{F2} has eighteen electrons in a $(10\text{e},6\text{o})$ active space,
the largest electron count in any active space here, ten qubits, but only
$35$ ansatz parameters. Seventeen nodes from $0.9$ to $2.5\Ang$ in $0.1\Ang$
steps.

\paragraph{Variational results.}
\ce{F2} is the campaign's clearest demonstration that circuit \emph{depth},
not molecular size, governs VQE accuracy. It has more active electrons than
\ce{N2} or \ce{O2} and the same qubit count, yet its ansatz carries $35$
parameters against their $117$ and $68$, and it reaches mean error
$0.0589\mHa$ and maximum $0.6568\mHa$, comfortably inside chemical accuracy
(Table~\ref{tab:data:f2:bond}). Single-bond dissociation is also far less
multireference than breaking a triple bond. Hartree--Fock misses $359.7\mHa$
on this curve; VQE recovers $99.98\%$ of it.


% GENERATED by gen_tables.py -- do not edit by hand (tab:data:f2:bond)
\begin{table}[!htbp]
\caption{\ce{F2} bond length scan: Hartree--Fock, VQE and exact energies (Ha) at every geometry; $\delta = (E_{\mathrm{VQE}}-E_{\mathrm{exact}})\times 10^3$ (mHa).}
\label{tab:data:f2:bond}
\centering\footnotesize\setlength{\tabcolsep}{3pt}
\adjustbox{max width=\columnwidth}{%
\begin{tabular}{rrrrr}
\toprule
$r$ (\AA) & $E_{\mathrm{HF}}$ & $E_{\mathrm{VQE}}$ & $E_{\mathrm{exact}}$ & $\delta$ \\
\midrule
0.9 & -195.219409 & -195.228895 & -195.228895 & 0.000 \\
1 & -195.638041 & -195.652560 & -195.652560 & 0.000 \\
1.1 & -195.853469 & -195.876050 & -195.876050 & 0.000 \\
1.2 & -195.951873 & -195.986310 & -195.986310 & 0.000 \\
1.3 & -195.981228 & -196.031793 & -196.031793 & 0.000 \\
1.4 & -195.970461 & -196.041633 & -196.041633 & 0.000 \\
1.5 & -195.938144 & -196.034258 & -196.034258 & 0.000 \\
1.6 & -195.896120 & -196.020869 & -196.020870 & 0.000 \\
1.7 & -195.851442 & -196.007365 & -196.007366 & 0.000 \\
1.8 & -195.807933 & -195.996136 & -195.996137 & 0.001 \\
1.9 & -195.767510 & -195.987687 & -195.987690 & 0.003 \\
2 & -195.731098 & -195.981740 & -195.981746 & 0.007 \\
2.1 & -195.699107 & -195.977800 & -195.977801 & 0.001 \\
2.2 & -195.671626 & -195.975327 & -195.975335 & 0.007 \\
2.3 & -195.648491 & -195.973870 & -195.973900 & 0.030 \\
2.4 & -195.629339 & -195.973075 & -195.973370 & 0.295 \\
2.5 & -195.613682 & -195.972682 & -195.973339 & 0.657 \\
\bottomrule
\end{tabular}}
\end{table}


\paragraph{Reconstruction: the polynomial's best case.}
This is where global polynomial interpolation is vindicated most clearly. At
degree $13$ the leave-one-out RMS is $0.0221\mHa$, against $2.42\mHa$ for the
cubic spline and $31.28\mHa$ for the linear baseline. The polynomial is
$109\times$ better than the spline and $1415\times$ better than the
baseline, and it is the only reconstruction in the campaign to land more
than an order of magnitude \emph{below} the variational error of the data it
is fitting. At that point the interpolation has stopped being the limiting
source of error at all.

The conditions that make this work are the ones identified in
Section~\ref{sec:lih}: a smooth, gently varying, single-minimum curve with no
nearby complex singularity, sampled densely enough ($17$ nodes) to support a
high degree. The degree sweep (Table~\ref{tab:deg:f2:bond}) still shows a
worst-case RMS of $122.14\mHa$, confirming that the good behaviour comes from
choosing the right degree, not from high degree in general.

\figpair{f2_polynomial_curve.png}{\ce{F2} dissociation curve, polynomial
reconstruction at degree $13$, the most accurate reconstruction in the
campaign.}{f2_poly_curve}{f2_polynomial_degree_sweep.png}{\ce{F2} degree
sweep. Best degree $13$ at $0.022\mHa$ RMS.}{f2_sweep}
% GENERATED by gen_tables.py -- do not edit by hand (tab:deg:f2:bond)
\begin{table}[!htbp]
\caption{\ce{F2} bond length polynomial degree sweep: leave-one-out error (mHa) against fitted degree $d$; the best degree is in bold.}
\label{tab:deg:f2:bond}
\centering\footnotesize\setlength{\tabcolsep}{3pt}
\adjustbox{max width=\columnwidth}{%
\begin{tabular}{rrrr}
\toprule
Degree $d$ & max & RMS & MAE \\
\midrule
1 & 208.1189 & 122.1387 & 104.6113 \\
2 & 191.0836 & 100.2701 & 84.1729 \\
3 & 116.2164 & 60.0132 & 52.5021 \\
4 & 50.5991 & 25.5932 & 22.4370 \\
5 & 20.1042 & 8.1653 & 6.8857 \\
6 & 8.0573 & 2.8694 & 2.1300 \\
7 & 4.6188 & 1.6211 & 1.1736 \\
8 & 2.7988 & 1.0129 & 0.7095 \\
9 & 1.0813 & 0.4218 & 0.2804 \\
10 & 0.1979 & 0.0568 & 0.0317 \\
11 & 0.4207 & 0.1551 & 0.0822 \\
12 & 0.0668 & 0.0238 & 0.0119 \\
\textbf{13} & \textbf{0.0777} & \textbf{0.0221} & \textbf{0.0105} \\
14 & 0.6768 & 0.2349 & 0.1033 \\
15 & 1.0644 & 0.3923 & 0.1672 \\
\bottomrule
\end{tabular}}
\end{table}



\figpair{f2_polynomial_error.png}{\ce{F2} polynomial versus linear
leave-one-out error. Note the logarithmic separation between the two
methods.}{f2_poly_err}{f2_spline_curve.png}{\ce{F2} cubic spline
reconstruction, RMS $2.42\mHa$.}{f2_spline_curve}

\figone{f2_spline_error.png}{\ce{F2} spline leave-one-out error. A
$12.9\times$ improvement on the linear baseline, respectable on its own, but
two orders of magnitude behind the polynomial on this
curve.}{f2_spline_err}

\paragraph{Geometry.}
Polynomial $1.3932\Ang$ and spline $1.3931\Ang$ agree to four decimal places,
against an experimental $1.412\Ang$. The $0.019\Ang$ gap is the smallest
geometry error of any diatomic in the campaign and belongs entirely to the
model. That two structurally very different interpolants agree this closely
is itself evidence that the reconstruction is well determined here.

% ---------------------------------------------------------------------
\MoleculeBarrier
\subsubsection{\ce{H2O2}: hydrogen peroxide}
\label{sec:h2o2}

\paragraph{Setup.}
\ce{H2O2} has eighteen electrons in an $(8\text{e},6\text{o})$ active space,
ten qubits, $92$ parameters, COBYLA at $1000$ evaluations. It is the
geometrically richest system in the campaign, the only molecule with a
genuine \emph{torsional} degree of freedom. Four scans were run: an O--O
stretch ($21$ nodes, $1.0$--$3.0\Ang$), an O--O--H bend ($13$ nodes,
$80$--$140^\circ$), a dihedral twist ($19$ nodes, $0$--$180^\circ$) and a
$10\times5$ dihedral $\times$ angle surface ($50$ points).

\figwide{h2o2_surface_anim_still.png}{\ce{H2O2} energy surface over the
dihedral $\times$ bend grid (animation frame). The torsional coordinate is
the soft direction; the bend is stiffer.}{h2o2_anim}

\figwide{h2o2_surface_scan_3d.png}{\ce{H2O2} completed energy surface, static
render of the $10\times5$ grid.}{h2o2_surf3d}

\figwide{h2o2_twist_anim_still.png}{\ce{H2O2} dihedral twist (animation
frame). The molecule rotates about the O--O axis through the full
$0$--$180^\circ$ range while the energy curve is traced
alongside.}{h2o2_twist_anim}

\paragraph{Variational results.}
Accuracy is moderate and depends on the scan: twist mean $0.2545\mHa$,
surface $0.2830\mHa$, stretch $0.6311\mHa$, bend $0.8882\mHa$, all four
inside chemical accuracy on the mean, with maxima of $0.45$, $0.92$, $2.10$
and $2.97\mHa$ respectively. The $92$-parameter ansatz with a
$1000$-evaluation budget is close to the practical limit of what COBYLA can
handle, and this system sits right at the boundary between the
well-converged and stagnating groups from Section~\ref{sec:protocol}.

The torsional potential is physically the most interesting curve in the
campaign. The computed minimum falls at a dihedral of $90^\circ$
(Table~\ref{tab:data:h2o2:twist}), reflecting the competition between
lone-pair repulsion, which favours a perpendicular arrangement, and
hyperconjugation, which favours planarity. The experimental gas-phase
dihedral is $111.5^\circ$; the model's preference for exactly $90^\circ$
reflects the minimal basis, and the curve is flat enough near there that the
$0.45\mHa$ variational noise is a non-trivial fraction of the feature being
located.


% GENERATED by gen_tables.py -- do not edit by hand (tab:data:h2o2:bend)
\begin{table}[!htbp]
\caption{\ce{H2O2} bend angle scan: Hartree--Fock, VQE and exact energies (Ha) at every geometry; $\delta = (E_{\mathrm{VQE}}-E_{\mathrm{exact}})\times 10^3$ (mHa).}
\label{tab:data:h2o2:bend}
\centering\footnotesize\setlength{\tabcolsep}{3pt}
\adjustbox{max width=\columnwidth}{%
\begin{tabular}{rrrrr}
\toprule
$\theta$ (deg) & $E_{\mathrm{HF}}$ & $E_{\mathrm{VQE}}$ & $E_{\mathrm{exact}}$ & $\delta$ \\
\midrule
80 & -148.714504 & -148.767895 & -148.769078 & 1.183 \\
85 & -148.734474 & -148.786281 & -148.786483 & 0.202 \\
90 & -148.748650 & -148.797650 & -148.798071 & 0.420 \\
95 & -148.757081 & -148.800381 & -148.803354 & 2.973 \\
100 & -148.759960 & -148.794902 & -148.795601 & 0.699 \\
105 & -148.757557 & -148.787872 & -148.789150 & 1.278 \\
110 & -148.750190 & -148.778444 & -148.778812 & 0.368 \\
115 & -148.738201 & -148.764620 & -148.764906 & 0.286 \\
120 & -148.721952 & -148.747332 & -148.747617 & 0.285 \\
125 & -148.701825 & -148.725833 & -148.727161 & 1.328 \\
130 & -148.678225 & -148.703047 & -148.703855 & 0.808 \\
135 & -148.651589 & -148.676904 & -148.678126 & 1.223 \\
140 & -148.622398 & -148.650035 & -148.650529 & 0.494 \\
\bottomrule
\end{tabular}}
\end{table}


% GENERATED by gen_tables.py -- do not edit by hand (tab:data:h2o2:stretch)
\begin{table}[!htbp]
\caption{\ce{H2O2} symmetric stretch scan: Hartree--Fock, VQE and exact energies (Ha) at every geometry; $\delta = (E_{\mathrm{VQE}}-E_{\mathrm{exact}})\times 10^3$ (mHa).}
\label{tab:data:h2o2:stretch}
\centering\footnotesize\setlength{\tabcolsep}{3pt}
\adjustbox{max width=\columnwidth}{%
\begin{tabular}{rrrrr}
\toprule
$r_{\mathrm{OO}}$ (\AA) & $E_{\mathrm{HF}}$ & $E_{\mathrm{VQE}}$ & $E_{\mathrm{exact}}$ & $\delta$ \\
\midrule
1 & -148.356555 & -148.375308 & -148.375804 & 0.496 \\
1.1 & -148.581109 & -148.595731 & -148.597497 & 1.766 \\
1.2 & -148.698717 & -148.712703 & -148.713400 & 0.697 \\
1.3 & -148.751629 & -148.765353 & -148.765879 & 0.526 \\
1.4 & -148.764399 & -148.789524 & -148.790117 & 0.594 \\
1.5 & -148.752258 & -148.792732 & -148.793111 & 0.379 \\
1.6 & -148.725418 & -148.785811 & -148.786349 & 0.538 \\
1.7 & -148.690988 & -148.774269 & -148.776365 & 2.097 \\
1.8 & -148.653786 & -148.764213 & -148.766074 & 1.860 \\
1.9 & -148.616828 & -148.756419 & -148.756861 & 0.442 \\
2 & -148.581809 & -148.748621 & -148.749290 & 0.669 \\
2.1 & -148.549568 & -148.742545 & -148.743419 & 0.874 \\
2.2 & -148.520461 & -148.738556 & -148.739059 & 0.502 \\
2.3 & -148.494595 & -148.735829 & -148.735937 & 0.108 \\
2.4 & -148.471936 & -148.732292 & -148.732553 & 0.261 \\
2.5 & -148.452357 & -148.731186 & -148.731434 & 0.248 \\
2.6 & -148.435646 & -148.730688 & -148.730768 & 0.081 \\
2.7 & -148.421533 & -148.730267 & -148.730403 & 0.136 \\
2.8 & -148.409706 & -148.730031 & -148.730224 & 0.193 \\
2.9 & -148.399842 & -148.729802 & -148.730184 & 0.382 \\
3 & -148.391626 & -148.729797 & -148.730203 & 0.406 \\
\bottomrule
\end{tabular}}
\end{table}


% GENERATED by gen_tables.py -- do not edit by hand (tab:data:h2o2:twist)
\begin{table}[!htbp]
\caption{\ce{H2O2} dihedral twist scan: Hartree--Fock, VQE and exact energies (Ha) at every geometry; $\delta = (E_{\mathrm{VQE}}-E_{\mathrm{exact}})\times 10^3$ (mHa).}
\label{tab:data:h2o2:twist}
\centering\footnotesize\setlength{\tabcolsep}{3pt}
\adjustbox{max width=\columnwidth}{%
\begin{tabular}{rrrrr}
\toprule
$\phi$ (deg) & $E_{\mathrm{HF}}$ & $E_{\mathrm{VQE}}$ & $E_{\mathrm{exact}}$ & $\delta$ \\
\midrule
0 & -148.746723 & -148.784878 & -148.785332 & 0.454 \\
10 & -148.747089 & -148.785451 & -148.785657 & 0.206 \\
20 & -148.748122 & -148.786349 & -148.786570 & 0.221 \\
30 & -148.749656 & -148.787716 & -148.787906 & 0.191 \\
40 & -148.751465 & -148.789252 & -148.789446 & 0.194 \\
50 & -148.753325 & -148.790769 & -148.790971 & 0.202 \\
60 & -148.755056 & -148.792095 & -148.792307 & 0.211 \\
70 & -148.756540 & -148.793116 & -148.793337 & 0.221 \\
80 & -148.757725 & -148.793759 & -148.794002 & 0.244 \\
90 & -148.758608 & -148.794034 & -148.794287 & 0.253 \\
100 & -148.759219 & -148.793946 & -148.794196 & 0.251 \\
110 & -148.759610 & -148.793467 & -148.793739 & 0.272 \\
120 & -148.759833 & -148.792648 & -148.792919 & 0.271 \\
130 & -148.759942 & -148.791457 & -148.791723 & 0.267 \\
140 & -148.759979 & -148.789859 & -148.790137 & 0.278 \\
150 & -148.759977 & -148.787891 & -148.788181 & 0.290 \\
160 & -148.759961 & -148.785767 & -148.786023 & 0.257 \\
170 & -148.759946 & -148.783879 & -148.784158 & 0.279 \\
180 & -148.759940 & -148.783110 & -148.783385 & 0.275 \\
\bottomrule
\end{tabular}}
\end{table}


\paragraph{Reconstruction.}
\ce{H2O2} is the cleanest spline win in the campaign: the spline takes all
three 1-D scans.

On the \emph{bend}, spline RMS is $0.776\mHa$ against $1.252\mHa$ for the
polynomial and $2.440\mHa$ for linear, a $3.1\times$ improvement on the
baseline. On the \emph{stretch}, the margin widens: $1.030\mHa$ spline,
$1.335\mHa$ polynomial, $14.617\mHa$ linear, a $14.2\times$ improvement. The
stretch degree sweep (Table~\ref{tab:deg:h2o2:stretch}) reaches a worst-case
RMS of $1446.8\mHa$, the second-largest polynomial failure in the campaign,
and the reason is visible in Table~\ref{tab:geom}: the full-degree polynomial
places the stretch minimum at $2.975\Ang$, more than an \AA ngstr\"om past
the true value of $1.458\Ang$ and near the far edge of the scanned interval.
That is the \ce{BeH2} edge-oscillation failure again, in its most extreme
1-D form.

The \emph{twist} is the one case in the campaign where the global polynomial
loses outright to the piecewise-linear baseline, $2.044\mHa$ against
$1.421\mHa$, with the spline narrowly ahead of both at $1.365\mHa$. A
torsional potential is periodic with a broad flat minimum, exactly the shape
a global polynomial handles worst; the degree sweep peaks at $858.5\mHa$.
The polynomial's interpolated minimum, $167.1^\circ$, sits nearly $80^\circ$
from the computed minimum, and even the spline's $134.3^\circ$ is off by
$44^\circ$ because the well is so shallow. Where the function is flat, tiny
energy errors turn into large positional errors, a statement about
conditioning, not about the interpolation method being wrong.

\paragraph{No surface reconstruction.}
Although the $50$-point surface was computed, the repository has no surface
entry in either the \texttt{h2o2\_polynomial/} or \texttt{h2o2\_spline/}
output; only the bend, stretch and twist were post-processed. \ce{H2O2}
therefore contributes no surface row to Table~\ref{tab:interp}. The
$10\times5$ grid would have left only $5\times3$ points after halving, below
the bicubic threshold on one axis, so the omission is defensible, but it is
still an omission and we record it as one.

\figpair{h2o2_bend_polynomial_curve.png}{\ce{H2O2} O--O--H bend, polynomial
reconstruction.}{h2o2_bend_poly}{h2o2_bend_polynomial_degree_sweep.png}%
{\ce{H2O2} bend degree sweep; best degree $5$.}{h2o2_bend_sweep}
% GENERATED by gen_tables.py -- do not edit by hand (tab:deg:h2o2:bend)
\begin{table}[!htbp]
\caption{\ce{H2O2} bend angle polynomial degree sweep: leave-one-out error (mHa) against fitted degree $d$; the best degree is in bold.}
\label{tab:deg:h2o2:bend}
\centering\footnotesize\setlength{\tabcolsep}{3pt}
\adjustbox{max width=\columnwidth}{%
\begin{tabular}{rrrr}
\toprule
Degree $d$ & max & RMS & MAE \\
\midrule
1 & 28.2146 & 20.9762 & 18.9605 \\
2 & 6.8346 & 3.5269 & 2.7413 \\
3 & 3.3955 & 1.8661 & 1.5398 \\
4 & 2.3948 & 1.4589 & 1.2840 \\
\textbf{5} & \textbf{2.5453} & \textbf{1.2523} & \textbf{1.0290} \\
6 & 3.2040 & 1.5291 & 1.3109 \\
7 & 3.5541 & 1.3976 & 1.0236 \\
8 & 4.9817 & 1.9929 & 1.4113 \\
9 & 7.0599 & 2.3425 & 1.4044 \\
10 & 23.2549 & 9.6742 & 5.5487 \\
11 & 21.0222 & 9.1279 & 4.9010 \\
\bottomrule
\end{tabular}}
\end{table}



\figpair{h2o2_bend_polynomial_error.png}{\ce{H2O2} bend, polynomial versus
linear leave-one-out
error.}{h2o2_bend_poly_err}{h2o2_bend_spline_curve.png}{\ce{H2O2} bend,
cubic spline reconstruction, RMS $0.776\mHa$.}{h2o2_bend_spline}

\figone{h2o2_bend_spline_error.png}{\ce{H2O2} bend, spline leave-one-out
error.}{h2o2_bend_spline_err}

\figpair{h2o2_stretch_polynomial_curve.png}{\ce{H2O2} O--O stretch,
polynomial reconstruction. The behaviour near the right edge produces the
spurious $2.975\Ang$ minimum discussed in the
text.}{h2o2_str_poly}{h2o2_stretch_polynomial_degree_sweep.png}{\ce{H2O2}
stretch degree sweep; worst-case RMS $1446.8\mHa$.}{h2o2_str_sweep}
% GENERATED by gen_tables.py -- do not edit by hand (tab:deg:h2o2:stretch)
\begin{table}[!htbp]
\caption{\ce{H2O2} symmetric stretch polynomial degree sweep: leave-one-out error (mHa) against fitted degree $d$; the best degree is in bold.}
\label{tab:deg:h2o2:stretch}
\centering\footnotesize\setlength{\tabcolsep}{3pt}
\adjustbox{max width=\columnwidth}{%
\begin{tabular}{rrrr}
\toprule
Degree $d$ & max & RMS & MAE \\
\midrule
1 & 103.8823 & 56.9259 & 46.0353 \\
2 & 106.5620 & 53.6256 & 44.8182 \\
3 & 78.1055 & 37.0579 & 32.1022 \\
4 & 40.4276 & 19.2324 & 16.4437 \\
5 & 17.7383 & 7.7269 & 6.6567 \\
6 & 6.9635 & 2.7242 & 2.1064 \\
7 & 4.4838 & 1.5437 & 1.0730 \\
8 & 3.9464 & 1.4061 & 1.1239 \\
\textbf{9} & \textbf{3.3708} & \textbf{1.3354} & \textbf{1.0558} \\
10 & 7.5316 & 2.1751 & 1.2736 \\
11 & 13.7322 & 3.5903 & 1.8743 \\
12 & 20.8655 & 5.6848 & 2.8409 \\
13 & 19.9068 & 6.1381 & 2.9960 \\
14 & 14.7073 & 3.4917 & 1.4574 \\
15 & 69.8646 & 16.4383 & 5.4986 \\
16 & 225.2926 & 73.7564 & 28.8422 \\
17 & 56.5335 & 16.5151 & 6.4874 \\
18 & 919.5614 & 259.8630 & 94.8950 \\
19 & 4434.1966 & 1446.8255 & 527.1237 \\
\bottomrule
\end{tabular}}
\end{table}



\figpair{h2o2_stretch_polynomial_error.png}{\ce{H2O2} stretch, polynomial
versus linear leave-one-out
error.}{h2o2_str_poly_err}{h2o2_stretch_spline_curve.png}{\ce{H2O2} stretch,
cubic spline reconstruction, RMS $1.030\mHa$.}{h2o2_str_spline}

\figone{h2o2_stretch_spline_error.png}{\ce{H2O2} stretch, spline
leave-one-out error.}{h2o2_str_spline_err}

\figpair{h2o2_twist_polynomial_curve.png}{\ce{H2O2} dihedral twist,
polynomial reconstruction. This is the one scan in the campaign where the
polynomial is beaten by piecewise-linear
interpolation.}{h2o2_twist_poly}{h2o2_twist_polynomial_degree_sweep.png}%
{\ce{H2O2} twist degree sweep; worst-case RMS $858.5\mHa$.}{h2o2_twist_sweep}
% GENERATED by gen_tables.py -- do not edit by hand (tab:deg:h2o2:twist)
\begin{table}[!htbp]
\caption{\ce{H2O2} dihedral twist polynomial degree sweep: leave-one-out error (mHa) against fitted degree $d$; the best degree is in bold.}
\label{tab:deg:h2o2:twist}
\centering\footnotesize\setlength{\tabcolsep}{3pt}
\adjustbox{max width=\columnwidth}{%
\begin{tabular}{rrrr}
\toprule
Degree $d$ & max & RMS & MAE \\
\midrule
1 & 5.4405 & 2.2236 & 1.6504 \\
2 & 5.9582 & 2.3140 & 1.6655 \\
3 & 6.0792 & 2.4288 & 1.7216 \\
4 & 5.7467 & 2.3789 & 1.8953 \\
5 & 4.8826 & 2.2276 & 1.5793 \\
6 & 4.7984 & 2.2888 & 1.7817 \\
7 & 4.3455 & 2.1211 & 1.7158 \\
8 & 7.3751 & 2.7163 & 1.9878 \\
9 & 5.6338 & 2.5426 & 1.9562 \\
\textbf{10} & \textbf{4.7536} & \textbf{2.0442} & \textbf{1.6239} \\
11 & 15.1352 & 4.6323 & 2.8286 \\
12 & 20.7084 & 5.4532 & 2.8669 \\
13 & 66.1262 & 16.9995 & 7.0632 \\
14 & 198.9544 & 57.1625 & 24.0925 \\
15 & 357.5401 & 122.0698 & 50.4368 \\
16 & 168.6051 & 41.3574 & 12.6249 \\
17 & 2485.2755 & 858.5445 & 336.1907 \\
\bottomrule
\end{tabular}}
\end{table}



\figpair{h2o2_twist_polynomial_error.png}{\ce{H2O2} twist, polynomial versus
linear leave-one-out error.}{h2o2_twist_poly_err}%
{h2o2_twist_spline_curve.png}{\ce{H2O2} twist, cubic spline reconstruction of
the torsional potential, RMS $1.365\mHa$.}{h2o2_twist_spline}

\figone{h2o2_twist_spline_error.png}{\ce{H2O2} twist, spline leave-one-out
error.}{h2o2_twist_spline_err}

% ---------------------------------------------------------------------
\MoleculeBarrier
\subsubsection{\ce{NO2}: nitrogen dioxide}
\label{sec:no2}

\paragraph{Setup.}
\ce{NO2} has twenty-three electrons, an \emph{odd} number, which makes it an
open-shell doublet radical ($2S=1$) needing an ROHF reference. Two active
spaces were used: $(5\text{e},4\text{o})$ on six qubits with $20$ parameters
for the bend and the surface, and $(7\text{e},5\text{o})$ on eight qubits
with $37$ parameters for the stretch. Three scans: bend ($19$ nodes,
$90$--$180^\circ$), symmetric stretch ($16$ nodes, $0.9$--$2.4\Ang$) and an
$11\times10$ surface ($110$ points, joint-largest in the campaign).

\figwide{no2_surface_scan_3d.png}{\ce{NO2} energy surface over the bond-length
$\times$ bond-angle grid, $110$ VQE solves.}{no2_surf3d}

\paragraph{Variational results.}
\ce{NO2} is the campaign's least accurate system after \ce{N2} and \ce{O2}:
bend mean $18.93\mHa$ (maximum $38.58$), stretch mean $6.34\mHa$ (maximum
$30.32$), surface mean $9.61\mHa$ (maximum $42.10$). All three miss chemical
accuracy by a wide margin.

The cause here is different from \ce{N2}. The ansatz is small, $20$
parameters for the bend, and COBYLA still uses $831$ evaluations, plenty for
that dimensionality. The difficulty is the open-shell doublet: a
$20$-parameter hardware-efficient circuit started from an ROHF reference has
limited ability to represent the correct spin state, so the variational
minimum it can reach genuinely sits above the exact ground state. This is
\emph{ansatz-expressivity} error rather than optimiser stagnation, and it is
the one place in the campaign where the two causes separate cleanly:
\ce{N2} had a large ansatz and too few evaluations; \ce{NO2} has enough
evaluations and too small an ansatz.


% GENERATED by gen_tables.py -- do not edit by hand (tab:data:no2:bend)
\begin{table}[!htbp]
\caption{\ce{NO2} bend angle scan: Hartree--Fock, VQE and exact energies (Ha) at every geometry; $\delta = (E_{\mathrm{VQE}}-E_{\mathrm{exact}})\times 10^3$ (mHa).}
\label{tab:data:no2:bend}
\centering\footnotesize\setlength{\tabcolsep}{3pt}
\adjustbox{max width=\columnwidth}{%
\begin{tabular}{rrrrr}
\toprule
$\theta$ (deg) & $E_{\mathrm{HF}}$ & $E_{\mathrm{VQE}}$ & $E_{\mathrm{exact}}$ & $\delta$ \\
\midrule
90 & -201.192217 & -201.249413 & -201.249413 & 0.000 \\
95 & -201.202353 & -201.263014 & -201.263014 & 0.000 \\
100 & -201.162788 & -201.248025 & -201.249429 & 1.403 \\
105 & -201.197736 & -201.253590 & -201.264120 & 10.530 \\
110 & -201.223419 & -201.252311 & -201.288507 & 36.196 \\
115 & -201.241848 & -201.268871 & -201.305991 & 37.120 \\
120 & -201.254512 & -201.280054 & -201.317954 & 37.900 \\
125 & -201.262530 & -201.286870 & -201.325445 & 38.575 \\
130 & -201.266755 & -201.291332 & -201.329274 & 37.942 \\
135 & -201.267863 & -201.294375 & -201.330082 & 35.707 \\
140 & -201.266406 & -201.301472 & -201.328399 & 26.927 \\
145 & -201.262865 & -201.301727 & -201.324693 & 22.966 \\
150 & -201.257690 & -201.299797 & -201.319415 & 19.618 \\
155 & -201.251353 & -201.296612 & -201.313040 & 16.429 \\
160 & -201.244391 & -201.292950 & -201.306124 & 13.174 \\
165 & -201.237464 & -201.289407 & -201.299342 & 9.935 \\
170 & -201.231400 & -201.286562 & -201.293512 & 6.950 \\
175 & -201.227163 & -201.284887 & -201.289520 & 4.633 \\
180 & -201.225626 & -201.284426 & -201.288094 & 3.669 \\
\bottomrule
\end{tabular}}
\end{table}


% GENERATED by gen_tables.py -- do not edit by hand (tab:data:no2:stretch)
\begin{table}[!htbp]
\caption{\ce{NO2} symmetric stretch scan: Hartree--Fock, VQE and exact energies (Ha) at every geometry; $\delta = (E_{\mathrm{VQE}}-E_{\mathrm{exact}})\times 10^3$ (mHa).}
\label{tab:data:no2:stretch}
\centering\footnotesize\setlength{\tabcolsep}{3pt}
\adjustbox{max width=\columnwidth}{%
\begin{tabular}{rrrrr}
\toprule
$r$ (\AA) & $E_{\mathrm{HF}}$ & $E_{\mathrm{VQE}}$ & $E_{\mathrm{exact}}$ & $\delta$ \\
\midrule
0.9 & -200.277971 & -200.322049 & -200.322049 & 0.000 \\
1 & -200.893684 & -200.947393 & -200.947393 & 0.000 \\
1.1 & -201.174147 & -201.236666 & -201.236666 & 0.000 \\
1.2 & -201.270208 & -201.339202 & -201.339202 & 0.000 \\
1.3 & -201.266809 & -201.350007 & -201.369011 & 19.004 \\
1.4 & -201.212115 & -201.305845 & -201.324876 & 19.031 \\
1.5 & -201.135184 & -201.226384 & -201.242355 & 15.971 \\
1.6 & -201.060449 & -201.119861 & -201.130914 & 11.053 \\
1.7 & -201.003196 & -201.034530 & -201.037817 & 3.287 \\
1.8 & -200.964938 & -200.974190 & -200.974660 & 0.470 \\
1.9 & -200.940909 & -200.942936 & -200.942974 & 0.039 \\
2 & -200.926344 & -200.927545 & -200.927553 & 0.009 \\
2.1 & -200.914352 & -200.914754 & -200.914755 & 0.001 \\
2.2 & -200.882229 & -200.887391 & -200.917712 & 30.320 \\
2.3 & -200.888450 & -200.888998 & -200.890390 & 1.392 \\
2.4 & -200.886437 & -200.887404 & -200.888322 & 0.918 \\
\bottomrule
\end{tabular}}
\end{table}


\paragraph{Reconstruction.}
The \ce{NO2} bend is the campaign's only case where the cubic spline is
beaten by the \emph{linear} baseline: spline $7.258\mHa$ against linear
$5.049\mHa$, with the polynomial best at $3.924\mHa$. The explanation is the
$18.93\mHa$ mean variational error. The spline's advantage rests on the data
being a smooth function sampled exactly; when each node carries an
independent error comparable to the curvature being fitted, the $C^2$
constraint spreads that noise into neighbouring intervals, while
piecewise-linear interpolation keeps each error inside its own two
subintervals. The general caveat that follows is important: spline
superiority assumes accurate data, and that assumption can be tested.

On the stretch, where the variational error is a third as large, normal
service resumes: polynomial $5.233\mHa$, spline $8.469\mHa$, linear
$54.194\mHa$, the polynomial improving on the baseline by $10.4\times$.

On the surface the spline wins clearly, $38.28\mHa$ against $46.23\mHa$ for
the total-degree polynomial and $112.96\mHa$ for bilinear, and the
polynomial's in-sample residual of $122.5\mHa$ shows it cannot even represent
the fitted points, let alone the held-out ones.

\paragraph{Geometry.}
The bend reconstructions disagree sharply. The spline places the minimum at
$142.27^\circ$ against a reference of $134.3^\circ$, an error of $8.0^\circ$.
The full-degree polynomial places it at $173.43^\circ$, an error of
$39.1^\circ$, close enough to the linear geometry to be qualitatively wrong
about the molecule's shape. The stretch shows the same pattern in
exaggerated form: spline $1.265\Ang$ versus reference $1.194\Ang$ (error
$0.071\Ang$), polynomial $2.373\Ang$ (error $1.179\Ang$, at the far edge of
the interval). The degree sweeps reach worst-case RMS values of $2170.2\mHa$
and $1071.8\mHa$ respectively, the largest 1-D polynomial failures in the
campaign.

\figpair{no2_bend_polynomial_curve.png}{\ce{NO2} bend, polynomial
reconstruction.}{no2_bend_poly}{no2_bend_polynomial_degree_sweep.png}%
{\ce{NO2} bend degree sweep; worst-case RMS $2170.2\mHa$.}{no2_bend_sweep}
% GENERATED by gen_tables.py -- do not edit by hand (tab:deg:no2:bend)
\begin{table}[!htbp]
\caption{\ce{NO2} bend angle polynomial degree sweep: leave-one-out error (mHa) against fitted degree $d$; the best degree is in bold.}
\label{tab:deg:no2:bend}
\centering\footnotesize\setlength{\tabcolsep}{3pt}
\adjustbox{max width=\columnwidth}{%
\begin{tabular}{rrrr}
\toprule
Degree $d$ & max & RMS & MAE \\
\midrule
1 & 19.8146 & 13.0253 & 11.5843 \\
2 & 19.5311 & 8.7503 & 6.7960 \\
3 & 16.3382 & 7.6784 & 6.3879 \\
4 & 17.6620 & 5.2658 & 2.8707 \\
5 & 17.9401 & 5.5656 & 3.4448 \\
6 & 16.7375 & 5.7283 & 3.8371 \\
7 & 15.0567 & 5.5441 & 4.2174 \\
8 & 10.1705 & 4.5566 & 3.6836 \\
\textbf{9} & \textbf{9.8915} & \textbf{3.9242} & \textbf{2.6250} \\
10 & 20.1489 & 6.5159 & 4.0064 \\
11 & 57.8952 & 15.7515 & 7.8352 \\
12 & 100.9952 & 29.3955 & 14.3663 \\
13 & 130.0082 & 37.4490 & 17.1481 \\
14 & 237.0671 & 64.1154 & 26.3076 \\
15 & 639.7099 & 201.0126 & 82.6372 \\
16 & 1101.9309 & 314.9069 & 122.6804 \\
17 & 6282.2910 & 2170.2329 & 849.8245 \\
\bottomrule
\end{tabular}}
\end{table}



\figpair{no2_bend_polynomial_error.png}{\ce{NO2} bend, polynomial versus
linear leave-one-out error.}{no2_bend_poly_err}{no2_bend_spline_curve.png}%
{\ce{NO2} bend, cubic spline reconstruction.}{no2_bend_spline}

\figone{no2_bend_spline_error.png}{\ce{NO2} bend, spline leave-one-out error.
At $7.26\mHa$ RMS this is \emph{worse} than the $5.05\mHa$ linear baseline,
the campaign's one spline defeat, caused by the $18.9\mHa$ variational noise
in the underlying data.}{no2_bend_spline_err}

\figpair{no2_stretch_polynomial_curve.png}{\ce{NO2} symmetric stretch,
polynomial reconstruction.}{no2_str_poly}%
{no2_stretch_polynomial_degree_sweep.png}{\ce{NO2} stretch degree sweep.}%
{no2_str_sweep}
% GENERATED by gen_tables.py -- do not edit by hand (tab:deg:no2:stretch)
\begin{table}[!htbp]
\caption{\ce{NO2} symmetric stretch polynomial degree sweep: leave-one-out error (mHa) against fitted degree $d$; the best degree is in bold.}
\label{tab:deg:no2:stretch}
\centering\footnotesize\setlength{\tabcolsep}{3pt}
\adjustbox{max width=\columnwidth}{%
\begin{tabular}{rrrr}
\toprule
Degree $d$ & max & RMS & MAE \\
\midrule
1 & 335.3914 & 180.4174 & 148.4173 \\
2 & 314.1513 & 177.4883 & 153.4797 \\
3 & 176.2432 & 100.7122 & 89.5247 \\
4 & 70.0349 & 31.8396 & 25.7123 \\
5 & 34.8118 & 15.9166 & 12.7452 \\
6 & 37.3378 & 17.3319 & 14.6221 \\
7 & 42.2889 & 14.6234 & 10.3301 \\
8 & 36.6522 & 15.5889 & 10.4453 \\
\textbf{9} & \textbf{10.3533} & \textbf{5.2332} & \textbf{4.2782} \\
10 & 50.3125 & 15.2589 & 8.7097 \\
11 & 83.0536 & 31.5000 & 17.0648 \\
12 & 43.3194 & 12.5409 & 6.7376 \\
13 & 407.8067 & 119.1430 & 51.8706 \\
14 & 2805.3755 & 1071.7516 & 479.7738 \\
\bottomrule
\end{tabular}}
\end{table}



\figpair{no2_stretch_polynomial_error.png}{\ce{NO2} stretch, polynomial
versus linear leave-one-out error.}{no2_str_poly_err}%
{no2_stretch_spline_curve.png}{\ce{NO2} stretch, cubic spline
reconstruction.}{no2_str_spline}

\figone{no2_stretch_spline_error.png}{\ce{NO2} stretch, spline leave-one-out
error, RMS $8.47\mHa$.}{no2_str_spline_err}

\figwide{no2_surface_polynomial_surface.png}{\ce{NO2} surface, total-degree
polynomial reconstruction. The in-sample residual alone is $122.5\mHa$.}{no2_surf_poly}
\figwide{no2_surface_polynomial_error.png}{\ce{NO2} surface,
polynomial holdout error.}{no2_surf_poly_err}

\figwide{no2_surface_spline_surface.png}{\ce{NO2} surface, bicubic spline
reconstruction from the half grid.}{no2_surf_spline}
\figwide{no2_surface_spline_error.png}{\ce{NO2} surface, spline holdout error, RMS
$38.28\mHa$ against $112.96\mHa$ bilinear.}{no2_surf_spline_err}

% ---------------------------------------------------------------------
\MoleculeBarrier
\subsubsection{\ce{N2O}: nitrous oxide}
\label{sec:n2o}

\paragraph{Setup.}
\ce{N2O} has twenty-two electrons in a linear triatomic with two inequivalent
bonds, which makes it the most thoroughly decomposed system in the campaign:
the N--N and N--O stretches were scanned \emph{separately}, plus a bend, plus
a two-dimensional $r_{\text{NN}} \times r_{\text{NO}}$ surface on a
$10\times6$ grid at fixed linear geometry. The bend and surface used
$(4\text{e},4\text{o})$ on six qubits ($26$ parameters); the two stretches
used $(6\text{e},5\text{o})$ on eight qubits ($54$ parameters).

\figwide{n2o_surface_scan_3d.png}{\ce{N2O} energy surface over the two
independent bond lengths at fixed linear geometry, $60$ VQE
solves.}{n2o_surf3d}

\paragraph{Variational results.}
Accuracy is poor throughout: bend mean $9.07\mHa$, N--N stretch $8.69\mHa$,
N--O stretch $2.72\mHa$, surface $6.98\mHa$. None reaches chemical accuracy.
The $(4\text{e},4\text{o})$ active space used for the bend and surface is
plainly too small for a $22$-electron molecule with two multiple bonds; it
correlates only four electrons out of twenty-two. The N--O stretch, which
used the larger $(6\text{e},5\text{o})$ space, is correspondingly the most
accurate of the four, and that correlation between active-space size and
accuracy is the clearest single-molecule evidence in the campaign that the
active space, not the optimiser, is the binding constraint here.


% GENERATED by gen_tables.py -- do not edit by hand (tab:data:n2o:bend)
\begin{table}[!htbp]
\caption{\ce{N2O} bend angle scan: Hartree--Fock, VQE and exact energies (Ha) at every geometry; $\delta = (E_{\mathrm{VQE}}-E_{\mathrm{exact}})\times 10^3$ (mHa).}
\label{tab:data:n2o:bend}
\centering\footnotesize\setlength{\tabcolsep}{3pt}
\adjustbox{max width=\columnwidth}{%
\begin{tabular}{rrrrr}
\toprule
$\theta$ (deg) & $E_{\mathrm{HF}}$ & $E_{\mathrm{VQE}}$ & $E_{\mathrm{exact}}$ & $\delta$ \\
\midrule
90 & -181.020241 & -181.088133 & -181.095279 & 7.145 \\
95 & -181.010381 & -181.075638 & -181.087151 & 11.513 \\
100 & -181.006934 & -181.065925 & -181.074079 & 8.154 \\
105 & -181.015267 & -181.064291 & -181.082783 & 18.491 \\
110 & -181.032663 & -181.081917 & -181.101767 & 19.850 \\
115 & -181.054324 & -181.112432 & -181.115480 & 3.048 \\
120 & -181.076910 & -181.143284 & -181.149129 & 5.844 \\
125 & -181.098411 & -181.164259 & -181.176590 & 12.331 \\
130 & -181.117703 & -181.189586 & -181.199002 & 9.417 \\
135 & -181.134324 & -181.209673 & -181.217386 & 7.713 \\
140 & -181.148277 & -181.225636 & -181.232551 & 6.915 \\
145 & -181.159787 & -181.238507 & -181.245117 & 6.610 \\
150 & -181.169136 & -181.249004 & -181.255556 & 6.551 \\
155 & -181.176593 & -181.257649 & -181.264224 & 6.574 \\
160 & -181.182389 & -181.264675 & -181.271365 & 6.689 \\
165 & -181.186709 & -181.270273 & -181.277101 & 6.828 \\
170 & -181.189696 & -181.270574 & -181.281411 & 10.837 \\
175 & -181.191449 & -181.273158 & -181.284138 & 10.980 \\
180 & -181.192027 & -181.278247 & -181.285079 & 6.832 \\
\bottomrule
\end{tabular}}
\end{table}


% GENERATED by gen_tables.py -- do not edit by hand (tab:data:n2o:nn_stretch)
\begin{table}[!htbp]
\caption{\ce{N2O} N--N stretch scan: Hartree--Fock, VQE and exact energies (Ha) at every geometry; $\delta = (E_{\mathrm{VQE}}-E_{\mathrm{exact}})\times 10^3$ (mHa).}
\label{tab:data:n2o:nn_stretch}
\centering\footnotesize\setlength{\tabcolsep}{3pt}
\adjustbox{max width=\columnwidth}{%
\begin{tabular}{rrrrr}
\toprule
$r_{\mathrm{NN}}$ (\AA) & $E_{\mathrm{HF}}$ & $E_{\mathrm{VQE}}$ & $E_{\mathrm{exact}}$ & $\delta$ \\
\midrule
0.9 & -180.847700 & -180.894792 & -180.903393 & 8.601 \\
1 & -181.093002 & -181.155884 & -181.165067 & 9.183 \\
1.1 & -181.183836 & -181.266783 & -181.273831 & 7.048 \\
1.2 & -181.190761 & -181.297568 & -181.300893 & 3.325 \\
1.3 & -181.155524 & -181.287329 & -181.288350 & 1.022 \\
1.4 & -181.102893 & -181.255671 & -181.257457 & 1.786 \\
1.5 & -181.047528 & -181.210977 & -181.216738 & 5.761 \\
1.6 & -180.997330 & -181.159528 & -181.174080 & 14.552 \\
1.7 & -180.955065 & -181.116350 & -181.139861 & 23.511 \\
1.8 & -180.920363 & -181.098984 & -181.117843 & 18.859 \\
1.9 & -180.891819 & -181.092540 & -181.104964 & 12.424 \\
2 & -180.868085 & -181.089781 & -181.097549 & 7.767 \\
2.1 & -180.848130 & -181.088461 & -181.093260 & 4.799 \\
2.2 & -180.831204 & -181.087759 & -181.090751 & 2.992 \\
\bottomrule
\end{tabular}}
\end{table}


% GENERATED by gen_tables.py -- do not edit by hand (tab:data:n2o:no_stretch)
\begin{table}[!htbp]
\caption{\ce{N2O} N--O stretch scan: Hartree--Fock, VQE and exact energies (Ha) at every geometry; $\delta = (E_{\mathrm{VQE}}-E_{\mathrm{exact}})\times 10^3$ (mHa).}
\label{tab:data:n2o:no_stretch}
\centering\footnotesize\setlength{\tabcolsep}{3pt}
\adjustbox{max width=\columnwidth}{%
\begin{tabular}{rrrrr}
\toprule
$r_{\mathrm{NO}}$ (\AA) & $E_{\mathrm{HF}}$ & $E_{\mathrm{VQE}}$ & $E_{\mathrm{exact}}$ & $\delta$ \\
\midrule
0.9 & -180.749006 & -180.859944 & -180.861662 & 1.718 \\
1 & -181.026828 & -181.139886 & -181.140782 & 0.896 \\
1.1 & -181.151195 & -181.256193 & -181.258391 & 2.198 \\
1.2 & -181.195496 & -181.281341 & -181.288287 & 6.947 \\
1.3 & -181.202329 & -181.262330 & -181.274458 & 12.128 \\
1.4 & -181.195642 & -181.230850 & -181.241697 & 10.847 \\
1.5 & -181.187179 & -181.204728 & -181.209500 & 4.773 \\
1.6 & -181.180884 & -181.189138 & -181.190220 & 1.082 \\
1.7 & -181.176903 & -181.181594 & -181.181757 & 0.163 \\
1.8 & -181.174373 & -181.178067 & -181.178089 & 0.023 \\
1.9 & -181.172581 & -181.176094 & -181.176097 & 0.004 \\
2 & -181.171143 & -181.174678 & -181.174679 & 0.001 \\
2.1 & -181.169889 & -181.173480 & -181.173481 & 0.001 \\
2.2 & -181.168758 & -181.172401 & -181.172401 & 0.000 \\
2.3 & -181.167736 & -181.171419 & -181.171419 & 0.000 \\
\bottomrule
\end{tabular}}
\end{table}


\paragraph{Reconstruction: the campaign's worst failure.}
The \ce{N2O} surface produces the most dramatic failure in this report. The
total-degree polynomial has an in-sample residual of $30.7\mHa$, already
poor, and a \emph{half-grid holdout RMS of $5186.4\mHa$}, with a maximum
error of $10362.8\mHa$. That is an error of more than ten Hartree on a
surface whose total energy variation is a fraction of one Hartree. The
reconstruction is off by four orders of magnitude, and $69\times$
\emph{worse} than the trivial bilinear baseline of $74.7\mHa$.

The mechanism is the two-dimensional version of edge oscillation. A
total-degree fit on a $5\times3$ subgrid, what remains of a $10\times6$ grid
after halving, has more coefficients than the data can pin down in the short
direction, so the fit is effectively unconstrained along one axis and drifts
away from the fitted points. The bicubic spline on the same subgrid achieves
$43.2\mHa$, comfortably beating bilinear's $74.7\mHa$. This is the strongest
argument in the report for using splines by default on surfaces: the spline
degrades gracefully where the polynomial fails catastrophically, and on an
expensive surface there is no way to find out which regime you are in after
the fact.

\paragraph{1-D reconstructions.}
The one-dimensional scans behave far better, and, unusually, split their
verdicts three ways. On the \emph{bend}, the spline wins at $1.895\mHa$
against $2.627\mHa$ polynomial and $3.502\mHa$ linear. On the \emph{N--N
stretch}, the spline wins again at $2.526\mHa$ against $3.812\mHa$ and
$25.902\mHa$, a $10.3\times$ improvement on the baseline. But on the
\emph{N--O stretch} the polynomial wins decisively: $0.119\mHa$ at degree
$12$ against $1.741\mHa$ spline and $26.824\mHa$ linear, a factor of $225$
over the baseline.

The same molecule, scanned along two chemically similar coordinates with the
same active space and node spacing, gives opposite verdicts, and that is the
sharpest demonstration in the report of its central methodological claim:
\emph{which interpolant wins is an empirical property of the particular
curve, and has to be measured rather than assumed}. The N--O curve happens to
be smooth enough over its interval to support a degree-$12$ fit; the N--N
curve, spanning a stronger bond over a wider effective range, is not.

\paragraph{Geometry.}
The N--O stretch reconstructions are excellent, placing the minimum at
$1.1924$ (polynomial) and $1.1927\Ang$ (spline) against a reference of
$1.185\Ang$, errors below $0.008\Ang$ from a $0.15\Ang$ grid. The N--N
stretch is worse, $1.213$ and $1.216\Ang$ against $1.128\Ang$, errors of
roughly $0.086\Ang$ that reflect the small active space rather than the
interpolation. On the bend, the spline correctly returns the linear geometry
$180.0^\circ$ exactly, while the polynomial reports $173.48^\circ$, once
again inventing a spurious interior minimum where the true minimum sits at
an endpoint.

\figpair{n2o_bend_polynomial_curve.png}{\ce{N2O} bend, polynomial
reconstruction.}{n2o_bend_poly}{n2o_bend_polynomial_degree_sweep.png}%
{\ce{N2O} bend degree sweep; worst-case RMS $1959.1\mHa$.}{n2o_bend_sweep}
% GENERATED by gen_tables.py -- do not edit by hand (tab:deg:n2o:bend)
\begin{table}[!htbp]
\caption{\ce{N2O} bend angle polynomial degree sweep: leave-one-out error (mHa) against fitted degree $d$; the best degree is in bold.}
\label{tab:deg:n2o:bend}
\centering\footnotesize\setlength{\tabcolsep}{3pt}
\adjustbox{max width=\columnwidth}{%
\begin{tabular}{rrrr}
\toprule
Degree $d$ & max & RMS & MAE \\
\midrule
1 & 42.2695 & 22.6748 & 19.3340 \\
2 & 39.8830 & 19.0849 & 16.0842 \\
3 & 25.9818 & 12.5738 & 10.9506 \\
4 & 10.5461 & 5.2566 & 4.3113 \\
5 & 12.0648 & 5.0184 & 3.5877 \\
6 & 10.5839 & 5.0119 & 4.0296 \\
7 & 6.0250 & 2.9613 & 2.4321 \\
\textbf{8} & \textbf{4.7589} & \textbf{2.6268} & \textbf{2.1837} \\
9 & 12.8454 & 4.0973 & 2.7974 \\
10 & 21.1114 & 7.1502 & 4.5029 \\
11 & 7.8085 & 3.0411 & 2.1155 \\
12 & 16.0187 & 5.4564 & 3.4003 \\
13 & 74.4241 & 24.3174 & 11.3541 \\
14 & 64.6002 & 16.0240 & 5.9935 \\
15 & 396.7793 & 103.9210 & 39.5718 \\
16 & 1921.8225 & 600.7182 & 239.2475 \\
17 & 5671.2321 & 1959.1411 & 767.1648 \\
\bottomrule
\end{tabular}}
\end{table}



\figpair{n2o_bend_polynomial_error.png}{\ce{N2O} bend, polynomial versus
linear leave-one-out error.}{n2o_bend_poly_err}{n2o_bend_spline_curve.png}%
{\ce{N2O} bend, cubic spline reconstruction.}{n2o_bend_spline}

\figone{n2o_bend_spline_error.png}{\ce{N2O} bend, spline leave-one-out error,
RMS $1.895\mHa$.}{n2o_bend_spline_err}

\figpair{n2o_nn_stretch_polynomial_curve.png}{\ce{N2O} N--N stretch,
polynomial reconstruction.}{n2o_nn_poly}%
{n2o_nn_stretch_polynomial_degree_sweep.png}{\ce{N2O} N--N stretch degree
sweep; best degree $7$.}{n2o_nn_sweep}
% GENERATED by gen_tables.py -- do not edit by hand (tab:deg:n2o:nn_stretch)
\begin{table}[!htbp]
\caption{\ce{N2O} N--N stretch polynomial degree sweep: leave-one-out error (mHa) against fitted degree $d$; the best degree is in bold.}
\label{tab:deg:n2o:nn_stretch}
\centering\footnotesize\setlength{\tabcolsep}{3pt}
\adjustbox{max width=\columnwidth}{%
\begin{tabular}{rrrr}
\toprule
Degree $d$ & max & RMS & MAE \\
\midrule
1 & 145.0041 & 78.7828 & 65.8487 \\
2 & 128.1584 & 78.0175 & 68.5992 \\
3 & 60.5396 & 38.5869 & 34.7940 \\
4 & 23.1560 & 9.4348 & 7.3997 \\
5 & 17.1673 & 7.9891 & 6.3352 \\
6 & 15.1735 & 8.0507 & 6.6646 \\
\textbf{7} & \textbf{8.0099} & \textbf{3.8120} & \textbf{3.1057} \\
8 & 19.3879 & 7.2945 & 5.0476 \\
9 & 19.9984 & 7.1784 & 4.7093 \\
10 & 34.2069 & 10.9983 & 6.0242 \\
11 & 106.0466 & 40.6286 & 21.5088 \\
12 & 173.3713 & 71.8440 & 36.0606 \\
\bottomrule
\end{tabular}}
\end{table}



\figpair{n2o_nn_stretch_polynomial_error.png}{\ce{N2O} N--N stretch,
polynomial versus linear leave-one-out error.}{n2o_nn_poly_err}%
{n2o_nn_stretch_spline_curve.png}{\ce{N2O} N--N stretch, cubic spline
reconstruction, RMS $2.526\mHa$.}{n2o_nn_spline}

\figone{n2o_nn_stretch_spline_error.png}{\ce{N2O} N--N stretch, spline
leave-one-out error.}{n2o_nn_spline_err}

\figpair{n2o_no_stretch_polynomial_curve.png}{\ce{N2O} N--O stretch,
polynomial reconstruction at degree $12$, RMS $0.119\mHa$, the second-best
reconstruction in the campaign.}{n2o_no_poly}%
{n2o_no_stretch_polynomial_degree_sweep.png}{\ce{N2O} N--O stretch degree
sweep. Contrast the shape with the N--N sweep
(Fig.~\ref{fig:n2o_nn_sweep}), the same molecule, opposite
verdicts.}{n2o_no_sweep}
% GENERATED by gen_tables.py -- do not edit by hand (tab:deg:n2o:no_stretch)
\begin{table}[!htbp]
\caption{\ce{N2O} N--O stretch polynomial degree sweep: leave-one-out error (mHa) against fitted degree $d$; the best degree is in bold.}
\label{tab:deg:n2o:no_stretch}
\centering\footnotesize\setlength{\tabcolsep}{3pt}
\adjustbox{max width=\columnwidth}{%
\begin{tabular}{rrrr}
\toprule
Degree $d$ & max & RMS & MAE \\
\midrule
1 & 138.0157 & 64.3447 & 46.4402 \\
2 & 143.1219 & 69.1727 & 57.3741 \\
3 & 97.8883 & 52.0650 & 44.7846 \\
4 & 40.5530 & 24.1991 & 21.4597 \\
5 & 10.3508 & 4.8643 & 4.0952 \\
6 & 3.2071 & 1.6978 & 1.4323 \\
7 & 2.1444 & 0.8976 & 0.7764 \\
8 & 4.2406 & 1.5355 & 1.0663 \\
9 & 2.9648 & 1.1689 & 0.7307 \\
10 & 1.1288 & 0.3554 & 0.1969 \\
11 & 2.1007 & 0.8345 & 0.4423 \\
\textbf{12} & \textbf{0.3536} & \textbf{0.1194} & \textbf{0.0584} \\
13 & 1.2711 & 0.5049 & 0.2386 \\
\bottomrule
\end{tabular}}
\end{table}



\figpair{n2o_no_stretch_polynomial_error.png}{\ce{N2O} N--O stretch,
polynomial versus linear leave-one-out error.}{n2o_no_poly_err}%
{n2o_no_stretch_spline_curve.png}{\ce{N2O} N--O stretch, cubic spline
reconstruction, RMS $1.741\mHa$.}{n2o_no_spline}

\figone{n2o_no_stretch_spline_error.png}{\ce{N2O} N--O stretch, spline
leave-one-out error.}{n2o_no_spline_err}



\figwide{n2o_surface_polynomial_surface.png}{\ce{N2O} surface, total-degree
polynomial reconstruction. Holdout RMS $5186.4\mHa$, the campaign's worst
failure, $69\times$ worse than bilinear
interpolation.}{n2o_surf_poly}
\figwide{n2o_surface_polynomial_error.png}{\ce{N2O}
surface, polynomial holdout error. Note the scale.}{n2o_surf_poly_err}

\figwide{n2o_surface_spline_surface.png}{\ce{N2O} surface, bicubic spline
reconstruction from the $5\times3$ half grid, the same data on which the
polynomial fails.}{n2o_surf_spline}
\figwide{n2o_surface_spline_error.png}{\ce{N2O}
surface, spline holdout error, RMS $43.2\mHa$ against $74.7\mHa$
bilinear.}{n2o_surf_spline_err}


\MoleculeBarrier
\subsection{Campaign-wide accuracy}
\label{sec:accuracy}
Table~\ref{tab:vqeacc} gives the variational error against exact
diagonalisation for every scan. Three things stand out.

\subsubsection{The variational bound holds everywhere}
Across all $692$ optimisations, $E_{\text{VQE}} \geq E_{\text{exact}}$ with
no exception. This is the correctness check Section~\ref{sec:theory}
promised: an energy below the exact reference would point to a Hamiltonian
construction error or a mismatched active space, and we found none.

\subsubsection{Accuracy tracks ansatz depth, not molecule size}
The scans split cleanly into two groups. Systems with modest parameter
counts, \ce{H2} ($\overline{\Delta}<0.001\mHa$), \ce{H2O}
($0.0007\mHa$), \ce{LiH} ($0.0001\mHa$), the \ce{BeH2} bend ($0.0044\mHa$),
the \ce{NH3} bend ($0.0158\mHa$), reach chemical accuracy by three or four
orders of magnitude. Systems with the deepest circuits, \ce{N2}
($p=117$, $\overline{\Delta}=10.48\mHa$), \ce{O2} ($p=68$, $10.55\mHa$),
the \ce{NO2} bend ($18.93\mHa$), fall short of it.

This is not a failure of the variational principle. It is optimiser
stagnation: COBYLA gets $1000$ evaluations, and on a $117$-parameter
landscape that budget runs out before convergence. The evidence is direct:
the scans that miss chemical accuracy are exactly the ones whose evaluation
counts saturate at the $1000$-evaluation ceiling, while SLSQP on \ce{LiH}
converges in $167$ evaluations and on \ce{NH3} in $255$--$348$, comfortably
inside budget, and both reach chemical accuracy. The fix is more evaluations
or a better-conditioned ansatz, not a different algorithm.

\subsubsection{Even the worst variational errors are small next to the
correlation energy}
The $\Delta^{\text{HF}}_{\max}$ column puts the errors in context. For
\ce{N2}, Hartree--Fock misses $777\mHa$ of correlation energy, while VQE
misses $10.5\mHa$ on average, so the variational calculation still recovers
about $98.6\%$ of what mean-field theory left out. For \ce{F2} the numbers
are $359.7\mHa$ missed by Hartree--Fock against $0.059\mHa$ by VQE, a
recovery of $99.98\%$. The method is doing real work even where it misses the
chemical-accuracy threshold.

% GENERATED by gen_tables.py -- do not edit by hand (tab:vqeacc)
\begin{table*}[!t]
\caption{Accuracy of the variational solution against exact diagonalisation of the same active-space Hamiltonian, in milli-Hartree. $\overline{\Delta}$ and $\Delta_{\max}$ are the mean and maximum of $|E_{\mathrm{VQE}}-E_{\mathrm{exact}}|$ over the scan; $\Delta^{\mathrm{HF}}_{\max}$ is the largest correlation energy the Hartree--Fock reference misses on the same points, i.e.\ the error the variational calculation is trying to recover. The chemical-accuracy threshold is $1.6$~mHa.}
\label{tab:vqeacc}
\centering\small\setlength{\tabcolsep}{3.5pt}

\adjustbox{max width=\textwidth}{%
\begin{tabular}{llrrrrrl}
\toprule
Molecule & Coordinate & $\overline{\Delta}$ & $\Delta_{\max}$ & $\Delta^{\mathrm{HF}}_{\max}$ & $E_{\min}$ (Ha) & evals & Chem.\ acc.? \\
\midrule
\ce{H2} & bond length & 0.0000 & 0.0003 & 277.58 & -1.137117 & -- & \textbf{yes} \\
\addlinespace[1pt]
\ce{LiH} & bond length & 0.0001 & 0.0006 & 153.21 & -7.880947 & 167 & \textbf{yes} \\
\addlinespace[1pt]
\ce{BeH2} & bend angle & 0.0044 & 0.0062 & 5.90 & -15.566211 & -- & \textbf{yes} \\
\ce{BeH2} & symmetric stretch & 0.5595 & 3.4778 & 68.08 & -15.567233 & -- & \textbf{yes} \\
\ce{BeH2} & 2-D surface & 0.4541 & 4.6537 & 68.08 & -15.564426 & -- & \textbf{yes} \\
\addlinespace[1pt]
\ce{H2O} & bend angle & 0.0007 & 0.0015 & 11.13 & -74.971353 & -- & \textbf{yes} \\
\addlinespace[1pt]
\ce{NH3} & bend angle & 0.0158 & 0.0494 & 21.13 & -55.473791 & 255 & \textbf{yes} \\
\ce{NH3} & symmetric stretch & 0.1937 & 0.7605 & 105.97 & -55.472133 & 348 & \textbf{yes} \\
\ce{NH3} & 2-D surface & 0.3343 & 4.5362 & 108.27 & -55.470597 & 347 & \textbf{yes} \\
\addlinespace[1pt]
\ce{N2} & bond length & 10.4839 & 36.3471 & 777.01 & -107.642667 & 1000 & no \\
\addlinespace[1pt]
\ce{O2} & bond length & 10.5491 & 28.5546 & 576.50 & -147.735735 & 1000 & no \\
\addlinespace[1pt]
\ce{F2} & bond length & 0.0589 & 0.6568 & 359.66 & -196.041633 & -- & \textbf{yes} \\
\addlinespace[1pt]
\ce{H2O2} & bend angle & 0.8882 & 2.9732 & 54.57 & -148.800381 & 1000 & \textbf{yes} \\
\ce{H2O2} & symmetric stretch & 0.6311 & 2.0967 & 338.58 & -148.792732 & 1000 & \textbf{yes} \\
\ce{H2O2} & 2-D surface & 0.2830 & 0.9229 & 62.30 & -148.811427 & 1000 & \textbf{yes} \\
\ce{H2O2} & dihedral twist & 0.2545 & 0.4539 & 38.61 & -148.794034 & 1000 & \textbf{yes} \\
\addlinespace[1pt]
\ce{NO2} & bend angle & 18.9302 & 38.5753 & 86.64 & -201.301727 & 831 & no \\
\ce{NO2} & symmetric stretch & 6.3435 & 30.3202 & 112.76 & -201.350007 & 955 & no \\
\ce{NO2} & 2-D surface & 9.6097 & 42.0953 & 188.72 & -201.304513 & 777 & no \\
\addlinespace[1pt]
\ce{N2O} & bend angle & 9.0698 & 19.8504 & 93.05 & -181.278247 & -- & no \\
\ce{N2O} & N--N stretch & 8.6878 & 23.5106 & 259.55 & -181.297568 & -- & no \\
\ce{N2O} & N--O stretch & 2.7185 & 12.1282 & 113.95 & -181.281341 & -- & no \\
\ce{N2O} & 2-D surface & 6.9779 & 27.0626 & 265.91 & -181.287840 & -- & no \\
\addlinespace[1pt]
\bottomrule
\end{tabular}}
\end{table*}


\subsection{Campaign-wide reconstruction}
\label{sec:recon}
Table~\ref{tab:interp} collects the out-of-sample reconstruction errors for
every post-processed scan, and Table~\ref{tab:geom} the recovered
equilibrium geometries; Section~\ref{sec:results} already interpreted these
numbers molecule by molecule, and Section~\ref{sec:discussion} draws the
conclusions that only make sense across all of them.

Two things worth restating here. First, \ce{H2} and \ce{H2O} were scanned but
never post-processed: the repository has no \texttt{h2\_polynomial/},
\texttt{h2\_spline/}, \texttt{h2o\_polynomial/} or \texttt{h2o\_spline/}
directory, and correspondingly no interpolation figures for either. Their
sections above present scan data only, and we said so there rather than
quietly skipping the molecules. Second, the \emph{Winner} column of
Table~\ref{tab:interp} is decided on out-of-sample RMS alone; where the
margin is narrow, the per-molecule discussion already said so.

% GENERATED by gen_tables.py -- do not edit by hand (tab:interp)
\begin{table*}[!t]
\caption{Out-of-sample reconstruction error for every scan that was post-processed, in milli-Hartree. One-dimensional scans are scored by leave-one-out at the interior nodes; surfaces by the half-grid holdout of Section~\ref{sec:holdout}. \emph{Linear} is the piecewise-linear baseline implied by every raw scan plot. A dash means the method was not applied to that scan; \emph{n/a} means the full interpolating degree left too few points to score.}
\label{tab:interp}
\centering\small\setlength{\tabcolsep}{3.5pt}

\adjustbox{max width=\textwidth}{%
\begin{tabular}{llrrrrrrrr}
\toprule
 & & \multicolumn{2}{c}{Linear baseline} & \multicolumn{3}{c}{Global polynomial} & \multicolumn{2}{c}{Cubic spline} & \\
\cmidrule(lr){3-4}\cmidrule(lr){5-7}\cmidrule(lr){8-9}
Molecule & Coordinate & RMS & max & best $d$ & RMS & worst RMS & RMS & max & Winner \\
\midrule
\ce{LiH} & bond length & 13.373 & 39.643 & 11 & 0.047 & 31.64 & 1.815 & 5.147 & polynomial \\
\addlinespace[1pt]
\ce{BeH2} & symmetric stretch & 10.119 & 25.310 & 7 & 4.496 & 572.51 & 4.647 & 10.984 & polynomial \\
\ce{BeH2} & bend angle & 0.344 & 0.561 & 9 & 0.000 & 7.59 & 0.001 & 0.002 & polynomial \\
\ce{BeH2} & 2-D surface & 17.969 & 48.857 & -- & 3.966 & -- & 3.097 & 10.052 & spline \\
\addlinespace[1pt]
\ce{NH3} & symmetric stretch & 12.787 & 29.486 & 6 & 0.466 & 53.42 & 0.920 & 1.608 & polynomial \\
\ce{NH3} & bend angle & 0.653 & 0.792 & 4 & 0.060 & 3.36 & 0.067 & 0.140 & polynomial \\
\ce{NH3} & 2-D surface & 30.108 & 65.067 & 6 & 8.493 & 74.48 & 8.284 & 14.359 & spline \\
\addlinespace[1pt]
\ce{N2} & bond length & 134.178 & 341.056 & 6 & 20.341 & 207.49 & 37.592 & 72.518 & polynomial \\
\addlinespace[1pt]
\ce{O2} & bond length & 80.541 & 202.904 & 6 & 7.987 & 130.17 & 22.002 & 43.203 & polynomial \\
\addlinespace[1pt]
\ce{F2} & bond length & 31.283 & 100.088 & 13 & 0.022 & 122.14 & 2.420 & 7.959 & polynomial \\
\addlinespace[1pt]
\ce{H2O2} & symmetric stretch & 14.617 & 51.725 & 9 & 1.335 & 1446.83 & 1.030 & 3.538 & spline \\
\ce{H2O2} & bend angle & 2.440 & 4.319 & 5 & 1.252 & 20.98 & 0.776 & 1.487 & spline \\
\ce{H2O2} & dihedral twist & 1.421 & 3.720 & 10 & 2.044 & 858.54 & 1.365 & 3.044 & spline \\
\addlinespace[1pt]
\ce{NO2} & symmetric stretch & 54.194 & 168.036 & 9 & 5.233 & 1071.75 & 8.469 & 21.643 & polynomial \\
\ce{NO2} & bend angle & 5.049 & 14.295 & 9 & 3.924 & 2170.23 & 7.258 & 22.116 & polynomial \\
\ce{NO2} & 2-D surface & 112.955 & 346.984 & -- & 46.230 & -- & 38.278 & 108.663 & spline \\
\addlinespace[1pt]
\ce{N2O} & N--N stretch & 25.902 & 75.096 & 7 & 3.812 & 78.78 & 2.525 & 6.210 & spline \\
\ce{N2O} & N--O stretch & 26.824 & 81.818 & 12 & 0.119 & 69.17 & 1.741 & 5.193 & polynomial \\
\ce{N2O} & bend angle & 3.502 & 9.629 & 8 & 2.627 & 1959.14 & 1.895 & 3.988 & spline \\
\ce{N2O} & 2-D surface & 74.718 & 231.025 & -- & 5186.448 & -- & 43.208 & 96.800 & spline \\
\addlinespace[1pt]
\bottomrule
\end{tabular}}
\end{table*}

% GENERATED by gen_tables.py -- do not edit by hand (tab:geom)
\begin{table*}[!t]
\caption{Equilibrium geometries recovered from the reconstructed curves and surfaces, compared with the best node actually computed and with the reference value (experimental where quoted in the scan metadata, otherwise the model's own optimised geometry). Lengths in \AA, angles in degrees. The gap between the \emph{best node} and the \emph{interpolated} column is the resolution the interpolant buys over the raw grid.}
\label{tab:geom}
\centering\small\setlength{\tabcolsep}{3.5pt}

\adjustbox{max width=\textwidth}{%
\begin{tabular}{lllrrrr}
\toprule
Molecule & Coordinate & Method & Best node & Interpolated & Reference & $|$Interp.$-$Ref.$|$ \\
\midrule
\ce{LiH} & bond length & polynomial & 1.65 & 1.5474 & 1.5949 & 0.0475 \\
\ce{LiH} & bond length & spline & 1.65 & 1.5496 & 1.5949 & 0.0453 \\
\ce{BeH2} & bend angle & polynomial & 180 & 180 & 180 & 0 \\
\ce{BeH2} & bend angle & spline & 180 & 180 & 180 & 0 \\
\ce{BeH2} & symmetric stretch & polynomial & 1.3 & 0.9305 & 1.3264 & 0.3959 \\
\ce{BeH2} & symmetric stretch & spline & 1.3 & 1.289 & 1.3264 & 0.0374 \\
\ce{NH3} & bend angle & polynomial & 105 & 104.3755 & 104.272 & 0.1035 \\
\ce{NH3} & bend angle & spline & 105 & 104.3845 & 104.272 & 0.1125 \\
\ce{NH3} & symmetric stretch & polynomial & 1.075 & 1.0391 & 1.0401 & 0.001 \\
\ce{NH3} & symmetric stretch & spline & 1.075 & 1.0394 & 1.0401 & 0.0007 \\
\ce{N2} & bond length & polynomial & 1.2 & 1.191 & 1.098 & 0.093 \\
\ce{N2} & bond length & spline & 1.2 & 1.1685 & 1.098 & 0.0705 \\
\ce{O2} & bond length & polynomial & 1.3 & 1.3028 & 1.208 & 0.0948 \\
\ce{O2} & bond length & spline & 1.3 & 1.2754 & 1.208 & 0.0674 \\
\ce{F2} & bond length & polynomial & 1.4 & 1.3932 & 1.412 & 0.0188 \\
\ce{F2} & bond length & spline & 1.4 & 1.3931 & 1.412 & 0.0189 \\
\bottomrule
\end{tabular}\hspace{1.5em}\begin{tabular}{lllrrrr}
\toprule
Molecule & Coordinate & Method & Best node & Interpolated & Reference & $|$Interp.$-$Ref.$|$ \\
\midrule
\ce{H2O2} & bend angle & polynomial & 95 & 93.6465 & 101.9 & 8.2535 \\
\ce{H2O2} & bend angle & spline & 95 & 93.938 & 101.9 & 7.962 \\
\ce{H2O2} & symmetric stretch & polynomial & 1.5 & 2.9748 & 1.458 & 1.5168 \\
\ce{H2O2} & symmetric stretch & spline & 1.5 & 1.4671 & 1.458 & 0.0091 \\
\ce{H2O2} & dihedral twist & polynomial & 150 & 167.0633 & 111.5 & 55.5633 \\
\ce{H2O2} & dihedral twist & spline & 150 & 134.278 & 111.5 & 22.778 \\
\ce{NO2} & bend angle & polynomial & 145 & 173.4323 & 134.3 & 39.1323 \\
\ce{NO2} & bend angle & spline & 145 & 142.269 & 134.3 & 7.969 \\
\ce{NO2} & symmetric stretch & polynomial & 1.3 & 2.3728 & 1.194 & 1.1788 \\
\ce{NO2} & symmetric stretch & spline & 1.3 & 1.2649 & 1.194 & 0.0709 \\
\ce{N2O} & bend angle & polynomial & 180 & 173.4797 & 180 & 6.5203 \\
\ce{N2O} & bend angle & spline & 180 & 180 & 180 & 0 \\
\ce{N2O} & N--N stretch & polynomial & 1.2 & 1.2132 & 1.128 & 0.0852 \\
\ce{N2O} & N--N stretch & spline & 1.2 & 1.2157 & 1.128 & 0.0877 \\
\ce{N2O} & N--O stretch & polynomial & 1.2 & 1.1924 & 1.185 & 0.0074 \\
\ce{N2O} & N--O stretch & spline & 1.2 & 1.1927 & 1.185 & 0.0077 \\
\bottomrule
\end{tabular}}
\end{table*}



\subsection{Discussion}
\label{sec:discussion}
\subsubsection{When does the variational solver succeed?}

Across $692$ optimisations, accuracy tracks ansatz conditioning rather than
molecular size. Table~\ref{tab:vqeacc} splits cleanly into three groups.

\textbf{Well-converged (14 of 23 scans).} Small to moderate parameter counts
with an optimiser that converges inside budget. \ce{H2}, \ce{H2O}, \ce{LiH},
\ce{F2}, the \ce{BeH2} bend and all three \ce{NH3} scans reach mean errors
between $10^{-4}$ and $0.2\mHa$, three to four orders of magnitude inside
chemical accuracy. \ce{F2} makes the point most sharply: ten qubits and ten
active electrons, yet only $35$ parameters, and a mean error of $0.059\mHa$.

\textbf{Optimiser-limited.} \ce{N2} ($p=117$) and \ce{O2} ($p=68$) both
saturate COBYLA's $1000$-evaluation ceiling at every geometry and land at
$10.5\mHa$ mean error. Fewer than nine evaluations per parameter is not
enough for a derivative-free method. The fix is more budget, a
gradient-based optimiser (the parameter-shift rule of
Section~\ref{sec:theory} exists for exactly this), or, as the
\ce{O2} anecdote in Section~\ref{sec:theory} shows, simply warm-starting from
a neighbouring geometry. It is not a different algorithm.

\textbf{Ansatz-limited.} \ce{NO2} has only $20$ parameters and uses $831$
evaluations, so it is not starved of budget; its $18.9\mHa$ bend error is a
genuine limit of what that circuit can represent for an open-shell doublet.
\ce{N2O} is limited the same way by a $(4\text{e},4\text{o})$ active space
over $22$ electrons, and the fact that its one scan with a larger active
space, the N--O stretch, $(6\text{e},5\text{o})$, is also its most accurate,
makes the diagnosis concrete.

The distinction matters because the two failure modes have opposite
remedies. Adding evaluations helps an optimiser-limited run and does nothing
for an ansatz-limited one; enlarging the active space helps the latter and
makes the former worse.

\subsubsection{When does each interpolant win?}

Counting the winners in Table~\ref{tab:interp}: the global polynomial takes
$9$ of $18$ scored scans, the spline $8$, and piecewise-linear interpolation
$1$ (the \ce{NO2} bend). No method dominates, and the exceptions are
informative.

\textbf{The polynomial wins on smooth, densely sampled 1-D curves.} \ce{F2}
($0.022\mHa$ at $17$ nodes), \ce{LiH} ($0.047\mHa$ at $13$ nodes) and the
\ce{N2O} N--O stretch ($0.119\mHa$ at $15$ nodes) are the three best
reconstructions in the campaign, and all three are global polynomials. The
common features are a single smooth minimum, no nearby complex singularity,
and enough nodes to support a high degree. Where these hold, the polynomial
wins by one to two orders of magnitude.

\textbf{The spline wins where the curve is shaped awkwardly.} All three
\ce{H2O2} scans, the \ce{N2O} N--N stretch, the \ce{N2O} bend, and every
surface but one go to the spline. Torsional potentials, flat wells and
functions with an endpoint minimum all defeat a global polynomial, and the
spline's locality handles them without incident.

\textbf{The spline wins decisively on surfaces.} Four of five surfaces go to
the spline, and the fifth, \ce{NH3}, is a statistical tie. The \ce{N2O}
surface is the crucial data point: the total-degree polynomial reaches a
holdout RMS of $5186\mHa$, $69\times$ worse than doing nothing clever at all,
while the bicubic spline on the same subgrid achieves $43.2\mHa$. Downside
risk of that size settles the default choice on its own.

\textbf{Linear interpolation wins when the data are noisy.} The \ce{NO2}
bend, carrying $18.9\mHa$ of variational error per node, is the one scan
where the spline loses to the baseline. Spline superiority assumes the
samples are accurate; when node error approaches the curvature being fitted,
the $C^2$ constraint spreads noise rather than suppressing it. That this
happened exactly once, in exactly the scan with the largest variational
error, is a reassuring consistency check.

\subsubsection{What interpolation buys, quantitatively}

Two concrete returns justify the post-processing layer.

\paragraph{Sub-grid geometry resolution.}
Table~\ref{tab:geom} shows equilibrium geometries recovered to
$0.0007$--$0.09\Ang$ from grids spaced at $0.075$--$0.25\Ang$. The best cases
stand out: \ce{NH3}'s stretch minimum lands within $0.0007\Ang$ of the model
reference from a $0.075\Ang$ grid, a hundredfold gain; \ce{LiH} reproduces
its model bond length to four decimals from a $0.25\Ang$ grid. Since the
piecewise-linear minimum is always a grid node by construction, none of these
gains would exist without a smooth reconstruction.

\paragraph{Avoiding quadratic cost on surfaces.}
The half-grid holdout answers a question with a direct budget attached. On
\ce{BeH2}, fitting a bicubic spline to one quarter of the $110$-point surface
reconstructs the remaining three quarters at $3.10\mHa$ RMS. On \ce{NO2},
$38.3\mHa$; on \ce{N2O}, $43.2\mHa$; on \ce{NH3}, $8.28\mHa$. At two to four
minutes per variational solve, running quarter grids would have saved
several hours per surface.

The caveat is the floor. On \ce{NH3}, halving a $7\times7$ grid leaves
$4\times4$, the bare minimum for a bicubic spline, and the holdout error of
$8.28\mHa$ sits well outside chemical accuracy. Below roughly five points per
axis there is not enough information left, whatever method is used. The
practical rule that follows: sample at least ten points per axis if the
surface is going to be halved, and use a spline.

\subsubsection{On reporting the \ce{LiH} result}

Section~\ref{sec:lih} documents an experiment written to demonstrate Runge's
phenomenon that failed to demonstrate it. We report this as the headline
finding for that molecule rather than bury it, for three reasons.

First, it is true, and the measurement is unambiguous: leave-one-out RMS
falls monotonically to the highest testable degree, and the full interpolant
has exactly one interior minimum.

Second, the reason is more instructive than the expected result would have
been. Runge's phenomenon is often taught as a consequence of high degree on
equally spaced nodes, which oversimplifies it: it is a consequence of the
interpolated function having a singularity near the real interval. A
numerical-analysis course is better served by a worked example that exposes
that distinction than by one more copy of the standard picture.

Third, the campaign has plenty of genuine Runge failures to balance it: the
\ce{BeH2} stretch, where the degree-$15$ interpolant invents a minimum at the
interval edge and misplaces the bond length by $0.40\Ang$; the \ce{H2O2}
stretch, misplacing it by $1.52\Ang$; the \ce{NO2} stretch, by $1.18\Ang$;
and the \ce{N2O} surface, failing by four orders of magnitude. The
phenomenon is real. It just is not automatic, and only measuring out of
sample tells the two cases apart.

\subsubsection{Supplementary: execution on IBM Quantum hardware}
\label{sec:ibm}

A separate Qiskit implementation queues the same \ce{LiH} calculation on real
IBM Quantum hardware, using \texttt{PySCFDriver}, a UCCSD ansatz mapped with
a \texttt{ParityMapper}, and SLSQP. The optimisation loop itself always runs
against a free, exact statevector estimator; only the already-optimised
circuit for each bond length is sent to real hardware for one final energy
evaluation. Reproducing this part of the project needs an IBM Quantum
account, whose credentials (a Cloud Resource Name and an API key) go in a
\texttt{.env} file copied from \texttt{.env.example}.

Running real hardware costs real quota, and the script is written around
that constraint: it picks whichever backend \texttt{service.least\_busy}
reports as least busy, submits every bond length as a single batched job
rather than one job per point, asks for explicit confirmation before
submitting, and fixes the shot count at $2000$ rather than leaving it at the
service default. The shot-count choice has a documented reason: an earlier
run left the shot count unset and burned about four minutes of a ten-minute
account quota on a single seven-point job before it had to be cancelled.
Error mitigation is switched off (\texttt{resilience\_level = 0}) to keep
runtime and cost down. That last choice matters for reading
Table~\ref{tab:ibm}: the huge hardware error below is device noise measured
with mitigation deliberately turned off to fit inside a small quota, not a
measurement of the hardware at its best.

\begin{table}[!htbp]
\caption{\ce{LiH} on an exact statevector simulator versus IBM Quantum
hardware, from \texttt{VQE demo on IBM Quantum Computer/}. Errors are
$|E_{\text{VQE}} - E_{\text{exact}}|$ in mHa. The simulator is accurate to
sub-micro-Hartree; the hardware, run with error mitigation off to control
cost, is off by more than two Hartree.}
\label{tab:ibm}
\centering\footnotesize
\adjustbox{max width=\columnwidth}{%
\begin{tabular}{rrrrr}
\toprule
$r$ (\AA) & $E_{\text{exact}}$ & $E_{\text{sim}}$ & $\delta_{\text{sim}}$
& $\delta_{\text{hw}}$ \\
\midrule
0.50 & $-7.047910$ & $-7.047910$ & $0.000287$ & $2143.85$ \\
0.75 & $-7.574579$ & $-7.574579$ & $0.000110$ & $2127.61$ \\
1.00 & $-7.784021$ & $-7.784021$ & $0.000028$ & --- \\
1.25 & $-7.861612$ & $-7.861612$ & $0.000043$ & --- \\
1.50 & $-7.882140$ & $-7.882140$ & $0.000105$ & --- \\
1.75 & $-7.876947$ & $-7.876947$ & $0.000075$ & --- \\
2.00 & $-7.860828$ & $-7.860828$ & $0.000049$ & --- \\
\bottomrule
\end{tabular}}
\end{table}

Only two of the seven bond lengths finished on hardware; the rest were left
out to stay inside the account's quota. The simulator reproduces exact
diagonalisation to better than $3\times10^{-4}\mHa$ at every point, four
orders of magnitude inside chemical accuracy. The two completed hardware runs
are off by about $2.14\Ha$, roughly $1340\times$ the chemical-accuracy
threshold, and larger than the entire correlation energy being sought. Gate
errors, readout errors, decoherence, and the decision to skip error
mitigation to save quota, together overwhelm the quantity of interest.

This is not a criticism of the hardware, it is a statement of where the
technology stands today with mitigation off and a tight quota, and it
justifies the architectural choice made throughout the rest of this project:
the physics and the numerical analysis are both better served by exact
classical simulation of the quantum circuits. Every result in
Sections~\ref{sec:protocol} and \ref{sec:results} is a statement about the
VQE \emph{algorithm}, not about a particular noisy device.

\subsubsection{Comparison with the base paper}
\label{sec:comparison}

The base paper~\cite{peruzzo2014} established that a variational loop of
short quantum circuits and a classical optimiser can recover a molecular
dissociation curve, demonstrated on one two-qubit molecule. Our campaign
confirms its two load-bearing claims at a much larger scale. The variational
bound it relies on held in all $692$ optimisations, and the algorithm reaches
chemical accuracy on $14$ of $23$ scans with up to ten qubits and $117$
parameters. Where it falls short, the cause is the classical half of the loop
the base paper introduced, the optimiser and the ansatz, not the variational
principle itself (Section~\ref{sec:discussion}). Our supplementary hardware
run (Section~\ref{sec:ibm}) also echoes the base paper's motivation from the
other side: on present-day hardware, without error mitigation, device noise is
larger than the entire correlation energy being sought. Finally, where the
base paper reports energies only at the sampled bond lengths, our
reconstruction results show that the curve between those samples can be
recovered to sub-grid accuracy, provided the reconstruction method is chosen
by out-of-sample measurement rather than assumed.

% =====================================================================
\FloatBarrier
\FloatBarrier
\section{Conclusion and Future Work}
\label{sec:conclusion}

\subsection{Key takeaways}
This project treated building molecular potential energy surfaces as what it
is from a numerical-analysis standpoint: reconstructing a smooth function
from a small number of very expensive samples. We implemented the
Variational Quantum Eigensolver over CUDA-Q and PySCF, ran $23$ scans across
$11$ molecules totalling $692$ variational optimisations, checked every
energy against exact diagonalisation of the identical active-space
Hamiltonian, and reconstructed each scan by three interpolation methods
scored strictly out of sample.

The main findings are these.

\begin{enumerate}[leftmargin=1.5em]
\item \textbf{VQE accuracy tracks ansatz conditioning, not molecular size.}
Fourteen of $23$ scans reach chemical accuracy, several by three to four
orders of magnitude. \ce{F2}, with ten active electrons and only $35$
parameters, outperforms \ce{N2}, with six active electrons and $117$. The
failures split cleanly into optimiser-limited and ansatz-limited cases with
opposite remedies.

\item \textbf{The variational bound held in all $692$ optimisations}, the
strongest evidence available that the Hamiltonian construction and
active-space bookkeeping are correct throughout.

\item \textbf{Splines are the safer default; polynomials are the better
gamble when the curve allows it.} The polynomial takes $9$ of $18$ scored
scans and produces the three best reconstructions in the campaign ($0.022$,
$0.047$ and $0.119\mHa$). But it also produces every catastrophic failure,
including a surface reconstruction $69\times$ worse than doing nothing. The
spline never fails badly. On an expensive surface, that asymmetry is worth
more than the polynomial's occasional win.

\item \textbf{Interpolation buys roughly an order of magnitude of effective
grid resolution for no extra quantum cost.} Equilibrium geometries were
located to $0.0007$--$0.09\Ang$ from grids spaced at $0.075$--$0.25\Ang$, and
the half-grid holdout shows that quarter-cost surfaces reconstruct to
$3$--$43\mHa$ RMS, provided at least about ten points per axis are sampled.

\item \textbf{Runge's phenomenon is a property of the function, not of the
degree.} On \ce{LiH} the global polynomial outperformed the cubic spline by
$38\times$ with no oscillation at any testable degree, and we found the node
density at which the expected failure does appear. The campaign still
contains four clear Runge failures, so the phenomenon is real, it is simply
not automatic, and only out-of-sample measurement tells the two cases apart.
\end{enumerate}

The idea that ties these together, and the one we would most want carried
forward, is that in-sample residuals carry no information for an
interpolant. Our polynomial fits reach residuals as low as
$1.4\times10^{-11}\mHa$ while being wrong by four orders of magnitude between
nodes. Only leave-one-out cross-validation and the half-grid holdout tell a
reconstruction that works apart from one that merely fits.

\subsection{Limitations of the current approach}
\label{sec:limitations}
We state the limitations of this work plainly.

\textbf{Minimal basis throughout.} Every calculation uses STO-3G. Absolute
energies are far from the basis-set limit, and equilibrium geometries carry
model errors of up to $0.09\Ang$ and several degrees (Table~\ref{tab:geom}).
No conclusion about \emph{chemistry} should be drawn from these numbers; the
conclusions here are about \emph{numerical methods}, for which internal
consistency against exact diagonalisation is what matters.

\textbf{Small active spaces.} Active spaces range from $(2\text{e},5\text{o})$
to $(10\text{e},6\text{o})$. For \ce{N2O}, $(4\text{e},4\text{o})$ correlates
four electrons out of twenty-two, inadequate by any chemical standard, and
that shows up in the results.

\textbf{Optimiser budgets not converged for the largest ans\"atze.}
\ce{N2}, \ce{O2}, \ce{NO2} and \ce{N2O} did not reach chemical accuracy. For
the first two the cause is a $1000$-evaluation ceiling, and those scans
should be rerun with a larger budget, or with warm-starting, before their
numbers are quoted as converged results.

\textbf{Coverage gaps.} \ce{H2} and \ce{H2O} were scanned but never
post-processed, so the campaign's densest scan and its canonical example
molecule contribute no interpolation evidence. The \ce{H2O2} surface was
computed but not reconstructed. The \ce{H2O} symmetric stretch is implemented
in the driver and documented in the project README, but no saved scan exists
in the repository. These are gaps in the campaign, not choices we can justify
after the fact.

\textbf{Half-grid holdout is a proxy.} The 2-D validation withholds three
quarters of the points from a rectangular subgrid. That is the right question
for cost planning, but it is not leave-one-out, and its error estimates run
pessimistic relative to what a well-chosen adaptive sampling scheme would
achieve.

\textbf{No treatment of noise or error mitigation in the main campaign.} The
main campaign simulates circuits exactly. Shot noise, gate error, and the
error-mitigation techniques that would matter on hardware sit outside its
scope, as Section~\ref{sec:ibm} makes clear, and the one hardware run we do
report was taken with error mitigation switched off to control cost.

\textbf{Reference values are heterogeneous.} Table~\ref{tab:geom} compares
some molecules against experimental values and others against the model's
own optimised geometry, depending on what the scan metadata recorded. The
column is labelled accordingly, but cross-molecule comparison of the final
column should keep that in mind.

\subsection{Future work}
The clearest next steps: rerun \ce{N2}, \ce{O2} and \ce{NO2} with
gradient-based optimisation and the parameter-shift rule, to separate
optimiser error from ansatz error definitively; close the coverage gaps by
post-processing \ce{H2}, \ce{H2O} and the \ce{H2O2} surface; replace the
uniform grids with Chebyshev-distributed nodes, which would remove the
edge-oscillation failures behind every bad polynomial geometry in
Table~\ref{tab:geom}; and extend the half-grid study into adaptive sampling,
placing new expensive points where the current reconstruction is least
certain instead of on a fixed lattice.

\begin{acks}
The authors thank their supervisor, Ishrat Jahan, for guidance throughout
CSE 402. This work builds on NVIDIA CUDA-Q Solvers, PySCF, Qiskit, NumPy,
SciPy and Matplotlib.
\end{acks}

% =====================================================================
\FloatBarrier
\section{Code and Installation}
\label{sec:code}

\ProjectBanner
\medskip

The complete implementation is in a public GitHub repository~\cite{projectrepo}.
It contains the VQE core, every per-molecule experiment driver, the numerical
post-processing scripts, the saved scan files from which every table and
figure in this report was generated, the supplementary IBM Quantum script and
the LaTeX source of this report. The code is Python and builds on CUDA-Q
Solvers~\cite{cudaq}, PySCF~\cite{pyscf}, Qiskit~\cite{qiskit},
NumPy~\cite{numpy}, SciPy~\cite{scipy} and Matplotlib~\cite{matplotlib}.

\paragraph{Installation.}
Clone the repository, create a virtual environment (Python 3.10 or newer) and
install the pinned requirements, as in Listing~\ref{lst:install}.

\begin{lstlisting}[language=bash,caption={Installation.},label={lst:install}]
git clone https://github.com/QuitttCat/CSE-402-Numerical-Analysis-Simulation-and-Modeling-Project.git
cd CSE-402-Numerical-Analysis-Simulation-and-Modeling-Project
python -m venv .venv
source .venv/bin/activate      # Windows: .venv\Scripts\activate
pip install -r requirements.txt
\end{lstlisting}

The per-molecule experiment drivers need only Qiskit, Qiskit Nature and PySCF.
\texttt{cudaq-solvers}, used by the \texttt{src/} core and the demo scripts, is
declared in \texttt{requirements.txt} as a Linux/WSL-only dependency
(\texttt{sys\_platform != "win32"}), so those parts need WSL or a Linux
machine, while the post-processing scripts run anywhere Python and SciPy do.
The supplementary IBM Quantum run additionally needs an IBM Quantum account,
whose credentials go in a \texttt{.env} file copied from
\texttt{.env.example} (Section~\ref{sec:ibm}).

\subsection{Reproduction}
\label{sec:repro}
Every figure and table in this report can be regenerated without re-running
any quantum simulation, because every scan is saved with its metadata. Note
that \texttt{cudaq-solvers} is a Linux/WSL-only dependency, so the VQE scans
themselves need a Linux environment even though the post-processing scripts
run anywhere Python and SciPy do.

\begin{lstlisting}[language=bash,caption={Reproducing the campaign.}]
# environment
python -m venv .venv && source .venv/bin/activate
pip install -r requirements.txt

# --- re-run a scan from scratch (expensive) ---
python experiment/h2/h2_ground_state_estimation.py \
       --start 0.3 --stop 2.5 --step 0.1
python experiment/h2o/h2o_ground_state_estimation.py \
       --coordinate bend --start 70 --stop 160 --step 5

# --- reopen a saved scan interactively (seconds) ---
python experiment/beh2/beh2_ground_state_estimation.py \
       --load beh2_surface_scan.json
python experiment/h2o/h2o_ground_state_estimation.py \
       --load h2o_bend_scan.json

# --- numerical post-processing (runs no chemistry) ---
python experiment/beh2_spline/beh2_spline_interpolation.py
python experiment/lih_polynomial/lih_polynomial_interpolation.py

# --- reproduce the LiH Runge finding ---
python experiment/lih/lih_ground_state_estimation.py \
       --step 0.125 --save lih_fine
python experiment/lih_polynomial/lih_polynomial_interpolation.py \
       --load lih_fine.json

# --- quick tour of every ansatz, optimiser and visualisation
#     backend, without running a full molecular scan ---
python examples/molecule_demo.py
python examples/vqe_demo.py --quantum

# --- regenerate this report's tables from the saved scans ---
python CSE402_report_acm/gen_tables.py
\end{lstlisting}

\subsection{Repository structure}
\label{sec:structure}

\begin{lstlisting}[language=bash,caption={Layout of the project repository.}]
src/
  molecule.py     # CUDA-Q / PySCF wrapper + 3D visualisation (1887 lines)
  vqe.py          # ansatze, optimisers, parameter-shift rule (740 lines)
examples/
  molecule_demo.py, vqe_demo.py   # runnable tour of src/
out/
  *.html          # visualisation backends exercised by molecule_demo.py
experiment/
  scan_io.py      # shared save/load contract, format_version guard
  <mol>/          # ground-state estimation driver + saved scans
  <mol>_polynomial/   # global polynomial reconstruction + LOO
  <mol>_spline/       # cubic/bicubic spline reconstruction + LOO
NEW_MOLECULE_PROMPT.md
  # the project's own notes for adding a new molecule: DOF-counting,
  # reusable code patterns, and a list of traps found the hard way
VQE demo on IBM Quantum Computer/
  qiskit_vqe.py, lih_energies.csv, lih_energies_simulator.csv
Presentation slides/
  Molecular Simulation - Proposal Presentation.pptx, Project Plan.pdf
CSE402_report_acm/
  report.tex, refs.bib, gen_tables.py, figures/, tables/   # this report
\end{lstlisting}

Molecules with an \texttt{experiment/} directory: \ce{H2}, \ce{H2O},
\ce{LiH}, \ce{BeH2}, \ce{N2}, \ce{O2}, \ce{F2}, \ce{NH3}, \ce{H2O2},
\ce{NO2}, \ce{N2O}. Of these, all but \ce{H2} and \ce{H2O} additionally have
\texttt{\_polynomial} and \texttt{\_spline} post-processing directories.

% =====================================================================
\bibliographystyle{ACM-Reference-Format}
\bibliography{refs}

\end{document}